\input{preamble}

% Bon, commençons ici.
%
\begin{document}

\title{Lissage des morphismes d'anneaux}


\maketitle

\phantomsection
\label{section-phantom}

\tableofcontents

\section{Introduction}
\label{section-introduction}

\noindent
Le résultat principal de ce chapitre est le suivant :
$$
\fbox{Un morphisme d'anneaux régulier entre anneaux noethériens est une
colimite filtrante de morphismes lisses.}
$$
Ce théorème est dû à Popescu, voir \cite{popescu-letter}.
Richard Swan a donné une présentation accessible de la démonstration de Popescu,
voir \cite{swan} ; il s'est appuyé sur des notes d'Andr\'e et sur un article d'Ogoma, voir
\cite{Ogoma}.

\medskip\noindent
Notre présentation suit celle de Swan, mais nous démontrons d'abord un résultat
intermédiaire qui permet de nous placer dans une situation un peu plus simple.
En voici un aperçu. Nous résolvons d'abord le ``problème de relèvement'' suivant :
une déformation infinitésimale plate d'une colimite filtrante d'algèbres lisses
est une colimite filtrante d'algèbres lisses. Ce résultat dit essentiellement
qu'il suffit de démontrer le théorème principal pour les morphismes entre
anneaux noethériens réduits. Nous démontrons ensuite deux lemmes très ingénieux,
appelés le ``lemme de relèvement'' et le ``lemme de désingularisation''.
Nous montrons que, combinés, ces lemmes ramènent le théorème principal à prouver
qu'une algèbre noethérienne $\Lambda$, géométriquement régulière sur un corps $k$,
est une colimite filtrante de $k$-algèbres lisses. Nous examinons ensuite les
arguments locaux nécessaires à la démonstration de Popescu-Ogoma-Swan-Andr\'e.
Enfin, nous donnons la démonstration dans les trois dernières sections.

\medskip\noindent
Nous terminons cette introduction par quelques indications bibliographiques.
Soit $A$ un anneau local noethérien hensélien.
Nous disons que $A$ possède la {\it propriété d'approximation} si, quels que soient
$f_1, \ldots, f_m \in A[x_1, \ldots, x_n]$, le système d'équations
$f_1 = 0, \ldots, f_m = 0$ admet une solution dans le complété de $A$ si et
seulement s'il admet une solution dans $A$. Cette définition est due à Artin.
Artin a d'abord démontré la propriété d'approximation pour les systèmes
d'équations analytiques, voir \cite{Artin-Analytic-Approximation}.
Dans \cite{Artin-Algebraic-Approximation}, Artin a démontré la propriété
d'approximation pour les anneaux locaux essentiellement de type fini sur un
anneau de valuation discrète excellent.
Artin a conjecturé (page 26 de \cite{Artin-Algebraic-Approximation}) que tout
anneau local hensélien excellent devait posséder la propriété d'approximation.

\medskip\noindent
À un certain moment, on a conjecturé que tout morphisme d'anneaux régulier entre
anneaux noethériens est une colimite filtrante d'algèbres lisses (voir par exemple
\cite{Raynaud-Rennes}, \cite{popescu-global}, \cite{Artin-power-series},
\cite{Artin-Denef}). Nous ne savons pas à qui cette conjecture\footnote{La
question/conjecture formulée dans \cite{Artin-power-series},
\cite{Artin-Denef} et \cite{popescu-global} est plus forte ; il a été démontré
dans \cite{Cipu} qu'elle est équivalente à la version originale.}
est due. Le lien avec la propriété d'approximation est le suivant : si
$A \to A^\wedge$ est une colimite d'algèbres lisses, avec $A$ comme ci-dessus,
alors la propriété d'approximation est satisfaite (insérer ici une référence
ultérieure). En outre, le théorème principal s'applique au morphisme
$A \to A^\wedge$ si $A$ est un anneau local excellent, car l'une des conditions
définissant un anneau local excellent est que ses fibres formelles soient
géométriquement régulières. Notons que les anneaux locaux excellents ont été
définis par Grothendieck et que leur définition a été publiée en 1965.

\medskip\noindent
Dans \cite{Artin-power-series}, il a été démontré que
$R \to R^\wedge$ est une colimite filtrante d'algèbres lisses pour tout anneau
local $R$ essentiellement de type fini sur un corps.
Dans \cite{Rotthaus-Artin}, il a été démontré que $R \to R^\wedge$ est une
colimite filtrante d'algèbres lisses pour tout anneau local $R$ essentiellement
de type fini sur un anneau de valuation discrète excellent.
Enfin, le théorème principal a été démontré dans
\cite{popescu-GND}, \cite{popescu-GNDA}, \cite{popescu-letter}, 
\cite{Ogoma} et \cite{swan}, comme indiqué ci-dessus.

\medskip\noindent
Réciproquement, à l'aide de certains des résultats précédents, il a été démontré
dans \cite{Rotthaus-excellent} que tout anneau local noethérien possédant la
propriété d'approximation est excellent.

\medskip\noindent
L'article \cite{Spivakovsky} propose une autre approche du théorème principal,
mais il semble difficile à lire (par exemple, \cite[Lemma 5.2]{Spivakovsky}
semble être une reformulation incorrecte de \cite[Lemma 3]{Elkik}). Il existe
également un exposé Bourbaki sur ce sujet, voir \cite{Teissier}.












\section{Idéaux singuliers}
\label{section-singular-ideal}

\noindent
Soit $R \to A$ un morphisme d'anneaux. L'idéal singulier de $A$ sur $R$ est
l'idéal radical de $A$ qui définit le lieu singulier du morphisme
$\Spec(A) \to \Spec(R)$. En voici une définition précise.

\begin{definition}
\label{definition-singular-ideal}
Soit $R \to A$ un morphisme d'anneaux. L'{\it idéal singulier de $A$ sur $R$},
noté $H_{A/R}$, est l'unique idéal radical $H_{A/R} \subset A$ tel que
$$
V(H_{A/R}) = \{\mathfrak q \in \Spec(A) \mid R \to A
\text{ non lisse en }\mathfrak q\}
$$
\end{definition}

\noindent
Cette définition a un sens car l'ensemble des idéaux premiers où $R \to A$
est lisse est ouvert, voir
Algèbre, Définition \ref{algebra-definition-smooth-at-prime}.
Afin de trouver un ensemble explicite de générateurs de l'idéal singulier,
nous démontrons d'abord le lemme suivant.

\begin{lemma}
\label{lemma-find-strictly-standard}
Soit $R$ un anneau. Soit $A = R[x_1, \ldots, x_n]/(f_1, \ldots, f_m)$.
Soit $\mathfrak q \subset A$ un idéal premier. Supposons que $R \to A$ soit
lisse en l'idéal premier $\mathfrak q$. Il existe alors un $a \in A$, $a \not \in \mathfrak q$,
un entier $c$, $0 \leq c \leq \min(n, m)$, et des sous-ensembles
$U \subset \{1, \ldots, n\}$, $V \subset \{1, \ldots, m\}$
de cardinal $c$ tels que
$$
a = a' \det(\partial f_j/\partial x_i)_{j \in V, i \in U}
$$
pour un certain $a' \in A$, et
$$
a f_\ell \in (f_j, j \in V) + (f_1, \ldots, f_m)^2
$$
pour tout $\ell \in \{1, \ldots, m\}$.
\end{lemma}

\begin{proof}
Posons $I = (f_1, \ldots, f_m)$, de sorte que le complexe cotangent naïf de
$A$ sur $R$ soit homotopiquement équivalent à
$I/I^2 \to \bigoplus A\text{d}x_i$, voir
Algèbre, Lemme \ref{algebra-lemma-NL-homotopy}.
Nous utiliserons le fait que la formation du complexe cotangent naïf commute à
la localisation, voir Algèbre, Section \ref{algebra-section-netherlander},
et en particulier Algèbre, Lemme \ref{algebra-lemma-localize-NL}.
D'après Algèbre, Définitions \ref{algebra-definition-smooth} et
\ref{algebra-definition-smooth-at-prime}
nous voyons que $(I/I^2)_a \to \bigoplus A_a\text{d}x_i$ est une injection
scindée pour un certain $a \in A$, $a \not \in \mathfrak q$.
Après avoir renuméroté $x_1, \ldots, x_n$ et $f_1, \ldots, f_m$, nous pouvons
supposer que $f_1, \ldots, f_c$ forment une base de l'espace vectoriel
$I/I^2 \otimes_A \kappa(\mathfrak q)$ et que les images de
$\text{d}x_{c + 1}, \ldots, \text{d}x_n$ forment une base de
$\Omega_{A/R} \otimes_A \kappa(\mathfrak q)$. Ainsi, après avoir remplacé $a$
par $aa'$ pour un certain $a' \in A$, $a' \not \in \mathfrak q$, nous pouvons
supposer que $f_1, \ldots, f_c$ forment une base de $(I/I^2)_a$ et que les images
de $\text{d}x_{c + 1}, \ldots, \text{d}x_n$ forment une base de
$(\Omega_{A/R})_a$. Dans cette situation, $a^N$ satisfait les conditions du
lemme pour un entier $N$ suffisamment grand (avec $U = V = \{1, \ldots, c\}$).
\end{proof}

\noindent
Nous utiliserons la notion d'élément {\it strictement standard} de $A$ sur $R$.
Notre notion est légèrement plus faible que celle de l'article de Swan
\cite{swan}. Nous définissons aussi un élément {\it standard élémentaire}
comme étant du type obtenu dans le lemme précédent. Nous comparons les
différents types d'éléments dans le lemme \ref{lemma-compare-standard}.

\begin{definition}
\label{definition-strictly-standard}
Soit $R \to A$ un morphisme d'anneaux de présentation finie.
Nous disons qu'un élément $a \in A$ est {\it standard élémentaire dans $A$ sur $R$}
s'il existe une présentation
$A = R[x_1, \ldots, x_n]/(f_1, \ldots, f_m)$
et un entier $0 \leq c \leq \min(n, m)$ tels que
\begin{equation}
\label{equation-elementary-standard-one}
a = a' \det(\partial f_j/\partial x_i)_{i, j = 1, \ldots, c}
\end{equation}
pour un certain $a' \in A$, et
\begin{equation}
\label{equation-elementary-standard-two}
a f_{c + j} \in (f_1, \ldots, f_c) + (f_1, \ldots, f_m)^2
\end{equation}
pour $j = 1, \ldots, m - c$. Nous disons que $a \in A$ est
{\it strictement standard dans $A$ sur $R$} s'il existe une présentation
$A = R[x_1, \ldots, x_n]/(f_1, \ldots, f_m)$
et un entier $0 \leq c \leq \min(n, m)$ tels que
\begin{equation}
\label{equation-strictly-standard-one}
a = \sum\nolimits_{I \subset \{1, \ldots, n\},\ |I| = c}
a_I \det(\partial f_j/\partial x_i)_{j = 1, \ldots, c,\ i \in I}
\end{equation}
pour certains $a_I \in A$, et
\begin{equation}
\label{equation-strictly-standard-two}
a f_{c + j} \in (f_1, \ldots, f_c) + (f_1, \ldots, f_m)^2
\end{equation}
pour $j = 1, \ldots, m - c$.
\end{definition}

\noindent
Le lemme suivant est utile pour dégager des conséquences de
(\ref{equation-strictly-standard-one}).

\begin{lemma}
\label{lemma-parse-equation-strictly-standard-one}
Soit $R$ un anneau. Soit $A = R[x_1, \ldots, x_n]/(f_1, \ldots, f_m)$
et posons $I = (f_1, \ldots, f_m)$. Soit $a \in A$. Alors
(\ref{equation-strictly-standard-one}) implique qu'il existe un morphisme
$A$-linéaire
$\psi : \bigoplus\nolimits_{i = 1, \ldots, n} A \text{d}x_i \to A^{\oplus c}$
tel que la composée
$$
A^{\oplus c} \xrightarrow{(f_1, \ldots, f_c)}
I/I^2 \xrightarrow{f \mapsto \text{d}f}
\bigoplus\nolimits_{i = 1, \ldots, n} A \text{d}x_i
\xrightarrow{\psi}
A^{\oplus c}
$$
soit la multiplication par $a$. Réciproquement, si un tel $\psi$ existe,
alors $a^c$ satisfait (\ref{equation-strictly-standard-one}).
\end{lemma}

\begin{proof}
Il s'agit d'un cas particulier de
Algèbre, Lemme \ref{algebra-lemma-matrix-left-inverse}.
\end{proof}

\begin{lemma}[Elkik]
\label{lemma-elkik}
Soit $R \to A$ un morphisme d'anneaux de présentation finie.
L'idéal singulier $H_{A/R}$ est le radical de l'idéal engendré par les éléments
strictement standard de $A$ sur $R$, et aussi le radical de l'idéal engendré
par les éléments standard élémentaires de $A$ sur $R$.
\end{lemma}

\begin{proof}
Supposons que $a$ soit strictement standard dans $A$ sur $R$. Nous affirmons
que $A_a$ est lisse sur $R$, ce qui prouve que $a \in H_{A/R}$. En effet,
soient $A = R[x_1, \ldots, x_n]/(f_1, \ldots, f_m)$, $c$ et $a' \in A$
comme dans la définition \ref{definition-strictly-standard}.
Posons $I = (f_1, \ldots, f_m)$, de sorte que le complexe cotangent naïf de
$A$ sur $R$ soit donné par $I/I^2 \to \bigoplus A\text{d}x_i$.
L'hypothèse (\ref{equation-strictly-standard-two}) implique que $(I/I^2)_a$
est engendré par les classes de $f_1, \ldots, f_c$.
L'hypothèse (\ref{equation-strictly-standard-one}) implique que la différentielle
$(I/I^2)_a \to \bigoplus A_a\text{d}x_i$ admet un inverse à gauche, voir
le lemme \ref{lemma-parse-equation-strictly-standard-one}.
Ainsi, $R \to A_a$ est lisse par définition et d'après
Algèbre, Lemme \ref{algebra-lemma-localize-NL}.

\medskip\noindent
Soient $H_e, H_s \subset A$ les radicaux des idéaux engendrés par les éléments
standard élémentaires, resp.\ strictement standard, de $A$ sur $R$.
Par définition et d'après ce que nous venons de démontrer, nous avons
$H_e \subset H_s \subset H_{A/R}$. L'inclusion $H_{A/R} \subset H_e$ résulte
du lemme \ref{lemma-find-strictly-standard}.
\end{proof}

\begin{example}
\label{example-not-quasi-compact}
L'ensemble des points où un morphisme d'anneaux de présentation finie est lisse
n'est pas nécessairement un ouvert quasi-compact. Par exemple, soient
$R = k[x, y_1, y_2, y_3, \ldots]/(xy_i)$ et $A = R/(x)$.
Alors le lieu de lissité de $R \to A$ est $\bigcup D(y_i)$, qui n'est pas
quasi-compact.
\end{example}

\begin{lemma}
\label{lemma-strictly-standard-base-change}
Soit $R \to A$ un morphisme d'anneaux de présentation finie.
Soit $R \to R'$ un morphisme d'anneaux. Si $a \in A$ est standard élémentaire,
resp.\ strictement standard, dans $A$ sur $R$, alors $a \otimes 1$ est standard
élémentaire, resp.\ strictement standard, dans $A \otimes_R R'$ sur $R'$.
\end{lemma}

\begin{proof}
Si $A = R[x_1, \ldots, x_n]/(f_1, \ldots, f_m)$ est une présentation de $A$
sur $R$, alors
$A \otimes_R R' = R'[x_1, \ldots, x_n]/(f'_1, \ldots, f'_m)$
est une présentation de $A \otimes_R R'$ sur $R'$. Ici, $f'_j$ est l'image de
$f_j$ dans $R'[x_1, \ldots, x_n]$. Le résultat découle donc des définitions.
\end{proof}

\begin{lemma}
\label{lemma-final-solve}
Soient $R \to A \to \Lambda$ des morphismes d'anneaux, avec $A$ de présentation
finie sur $R$. Supposons que $H_{A/R} \Lambda = \Lambda$. Il existe alors une
factorisation $A \to B \to \Lambda$ telle que $B$ soit lisse sur $R$.
\end{lemma}

\begin{proof}
Choisissons $f_1, \ldots, f_r \in H_{A/R}$ et
$\lambda_1, \ldots, \lambda_r \in \Lambda$ tels que
$\sum f_i\lambda_i = 1$ dans $\Lambda$. Posons
$B = A[x_1, \ldots, x_r]/(f_1x_1 + \ldots + f_rx_r - 1)$
et définissons $B \to \Lambda$ en envoyant $x_i$ sur $\lambda_i$.
Pour vérifier que $B$ est lisse sur $R$, utilisons le fait que $A_{f_i}$ est
lisse sur $R$ par définition de $H_{A/R}$ et que $B_{f_i}$ est lisse sur
$A_{f_i}$. Détails omis.
\end{proof}





\section{Présentations d'algèbres}
\label{section-presentations}

\noindent
Certains des résultats de cette section sont dus à Elkik. Notons que l'algèbre
$C$ du lemme suivant est une algèbre symétrique sur $A$. De plus, si $R$ est
noethérien, alors $C$ est de présentation finie sur $R$.

\begin{lemma}
\label{lemma-improve-presentation}
Soit $R$ un anneau et soit $A$ une $R$-algèbre de présentation finie.
Il existe un morphisme de $R$-algèbres de type fini $A \to C$ qui admet une
rétraction et possède les deux propriétés suivantes :
\begin{enumerate}
\item pour tout $a \in A$ tel que $R \to A_a$ soit une intersection complète
locale (Compléments d'algèbre, Définition
\ref{more-algebra-definition-local-complete-intersection}),
l'anneau $C_a$ est lisse sur $A_a$ et admet une présentation
$C_a = R[y_1, \ldots, y_m]/J$ telle que $J/J^2$ soit libre sur $C_a$, et
\item pour tout $a \in A$ tel que $A_a$ soit lisse sur $R$, le module
$\Omega_{C_a/R}$ est libre sur $C_a$.
\end{enumerate}
\end{lemma}

\begin{proof}
Choisissons une présentation $A = R[x_1, \ldots, x_n]/I$ et écrivons
$I = (f_1, \ldots, f_m)$. Définissons le $A$-module $K$ par la suite exacte courte
$$
0 \to K \to A^{\oplus m} \to I/I^2 \to 0
$$
où le $j$-ième vecteur de base $e_j$ du terme central est envoyé sur la classe de
$f_j$ dans le terme de droite. Posons
$$
C = \text{Sym}^*_A(I/I^2).
$$
La rétraction est simplement la projection sur la composante de degré $0$ de
$C$. Nous avons une surjection
$R[x_1, \ldots, x_n, y_1, \ldots, y_m] \to C$ qui envoie $y_j$ sur la classe de
$f_j$ dans $I/I^2$. Le noyau $J$ de ce morphisme est engendré par les éléments
$f_1, \ldots, f_m$ et par les éléments $\sum h_j y_j$, où
$h_j \in R[x_1, \ldots, x_n]$, tels que $\sum h_j e_j$ définisse un élément de
$K$. D'après
Algèbre, Lemme \ref{algebra-lemma-exact-sequence-NL}
appliqué à $R \to A \to C$ et aux présentations ci-dessus, ainsi que
Compléments d'algèbre, Lemme
\ref{more-algebra-lemma-cotangent-complex-symmetric-algebra}
il existe une suite exacte
\begin{equation}
\label{equation-sequence}
I/I^2 \otimes_A C \to J/J^2 \to K \otimes_A C \to 0
\end{equation}
de $C$-modules. Soit $h \in R[x_1, \ldots, x_n]$ un élément d'image $a \in A$.
Nous utiliserons, pour les anneaux localisés, les présentations
$$
A_a = R[x_0, x_1, \ldots, x_n]/I'
\quad\text{et}\quad
C_a = R[x_0, x_1, \ldots, x_n, y_1, \ldots, y_m]/J'
$$
où $I' = (hx_0 - 1, I)$ et $J' = (hx_0 - 1, J)$. Ainsi,
$I'/(I')^2 = A_a \oplus (I/I^2)_a$ comme $A_a$-modules et
$J'/(J')^2 = C_a \oplus (J/J^2)_a$ comme $C_a$-modules.
Nous obtenons donc
\begin{equation}
\label{equation-sequence-localized}
C_a \oplus I/I^2 \otimes_A C_a \to
C_a \oplus (J/J^2)_a \to
K \otimes_A C_a \to 0
\end{equation}
qui est la suite fournie par
Algèbre, lemme \ref{algebra-lemma-exact-sequence-NL},
pour $R \to A_a \to C_a$ et les présentations ci-dessus.

\medskip\noindent
Supposons maintenant que $a \in A$ soit tel que $A_a$ soit une intersection
complète locale sur $R$. Alors $(I/I^2)_a$ est projectif de type fini sur $A_a$,
voir
Compléments d'algèbre, Lemme
\ref{more-algebra-lemma-quasi-regular-ideal-finite-projective}.
Nous voyons donc que $K_a \oplus (I/I^2)_a \cong A_a^{\oplus m}$ est libre.
En particulier, $K_a$ est lui aussi projectif de type fini.
D'après Compléments d'algèbre, Lemme \ref{more-algebra-lemma-transitive-lci-at-end},
la suite (\ref{equation-sequence-localized}) est exacte à gauche. Ainsi,
$$
J'/(J')^2 \cong
C_a \oplus I/I^2 \otimes_A C_a \oplus K \otimes_A C_a \cong
C_a^{\oplus m + 1}
$$
Cela démontre (1). Supposons enfin que, de plus, $A_a$ soit lisse sur $R$.
La même présentation montre alors que le module des différentielles de Kähler
$\Omega_{C_a/R}$ est le conoyau du morphisme
$$
J'/(J')^2 \longrightarrow
\bigoplus\nolimits_i C_a\text{d}x_i \oplus \bigoplus\nolimits_j C_a\text{d}y_j
$$
Dans la décomposition ci-dessus, le facteur direct $C_a$ de $J'/(J')^2$
correspond à $hx_0 - 1$ et s'envoie donc isomorphiquement sur le facteur direct
$C_a\text{d}x_0$. Le facteur direct $I/I^2 \otimes_A C_a$ de $J'/(J')^2$
s'injecte dans $\bigoplus_{i = 1, \ldots, n} C_a\text{d}x_i$, avec pour quotient
$\Omega_{A_a/R} \otimes_{A_a} C_a$. Le facteur direct $K \otimes_A C_a$
s'injecte dans $\bigoplus_{j \geq 1} C_a\text{d}y_j$, avec un quotient isomorphe
à $I/I^2 \otimes_A C_a$. Ainsi, le conoyau du dernier morphisme affiché est le
module
$I/I^2 \otimes_A C_a \oplus \Omega_{A_a/R} \otimes_{A_a} C_a$.
Puisque $(I/I^2)_a \oplus \Omega_{A_a/R}$ est libre (d'après la définition des
morphismes d'anneaux lisses), nous voyons que (2) est satisfaite.
\end{proof}

\noindent
La proposition suivante a été démontrée par Elkik dans \cite{Elkik} pour les
morphismes d'anneaux lisses sur des couples henséliens. Pour les morphismes
d'anneaux lisses, on la trouve dans \cite{Arabia}, où il est également démontré
que les morphismes d'anneaux entre algèbres lisses peuvent être relevés.

\begin{proposition}
\label{proposition-lift-smooth}
\begin{slogan}
Les algèbres lisses et syntomiques se relèvent à travers les surjections
\end{slogan}
Soit $R \to R_0$ un morphisme d'anneaux surjectif de noyau $I$.
\begin{enumerate}
\item Si $R_0 \to A_0$ est un morphisme d'anneaux syntomique, alors il existe
un morphisme d'anneaux syntomique $R \to A$ tel que $A/IA \cong A_0$.
\item Si $R_0 \to A_0$ est un morphisme d'anneaux lisse, alors il existe un
morphisme d'anneaux lisse $R \to A$ tel que $A/IA \cong A_0$.
\end{enumerate}
\end{proposition}

\begin{proof}
Supposons que $R_0 \to A_0$ soit syntomique, donc en particulier une
intersection complète locale (Compléments d'algèbre, Lemme
\ref{more-algebra-lemma-syntomic-lci}).
Choisissons une présentation $A_0 = R_0[x_1, \ldots, x_n]/J_0$. Posons
$C_0 = \text{Sym}^*_{A_0}(J_0/J_0^2)$. Notons que $J_0/J_0^2$ est un
$A_0$-module projectif de type fini (Algèbre, Lemme
\ref{algebra-lemma-syntomic-presentation-ideal-mod-squares}).
D'après le lemme \ref{lemma-improve-presentation}, le morphisme d'anneaux
$A_0 \to C_0$ est lisse et nous pouvons trouver une présentation
$C_0 = R_0[y_1, \ldots, y_m]/K_0$ dans laquelle $K_0/K_0^2$ est libre sur
$C_0$. D'après Algèbre, Lemme \ref{algebra-lemma-huber}, nous pouvons supposer
$C_0 = R_0[y_1, \ldots, y_m]/(\overline{f}_1, \ldots, \overline{f}_c)$
où les images de $\overline{f}_1, \ldots, \overline{f}_c$ forment une base de
$K_0/K_0^2$ sur $C_0$. Choisissons des relèvements
$f_1, \ldots, f_c \in R[y_1, \ldots, y_m]$ de
$\overline{f}_1, \ldots, \overline{f}_c$ et posons
$$
C = R[y_1, \ldots, y_m]/(f_1, \ldots, f_c)
$$
Par construction, $C_0 = C/IC$. D'après Algèbre, Lemme
\ref{algebra-lemma-localize-relative-complete-intersection},
quitte à remplacer $C$ par $C_g$, nous pouvons supposer que $C$ est une
intersection complète globale relative sur $R$.
Nous en concluons qu'il existe un $A_0$-module projectif de type fini $P_0$ tel
que $C_0 = \text{Sym}^*_{A_0}(P_0)$ soit isomorphe à $C/IC$ pour une certaine
$R$-algèbre syntomique $C$.

\medskip\noindent
Choisissons un entier $n$ et une décomposition en somme directe
$A_0^{\oplus n} = P_0 \oplus Q_0$.
D'après Compléments d'algèbre, Lemme \ref{more-algebra-lemma-lift-projective-module},
nous pouvons trouver un morphisme d'anneaux \'etale $C \to C'$ qui induit un
isomorphisme $C/IC \to C'/IC'$ et un $C'$-module projectif de type fini $Q$ tel
que $Q/IQ$ soit isomorphe à $Q_0 \otimes_{A_0} C/IC$.
Alors $D = \text{Sym}_{C'}^*(Q)$ est une $C'$-algèbre lisse (voir
Compléments d'algèbre, Lemme \ref{more-algebra-lemma-symmetric-algebra-smooth}).
Considérons le diagramme
$$
\xymatrix{
R \ar[d] \ar[rr] & &
C \ar[r] \ar[d] &
C' \ar[r] \ar[d] &
D \ar[d] \\
R/I \ar[r] &
A_0 \ar[r] &
C/IC \ar[r]^{\cong} &
C'/IC' \ar[r] &
D/ID
}
$$
Observons que notre choix de $Q$ donne
\begin{align*}
D/ID & =
\text{Sym}_{C/IC}^*(Q_0 \otimes_{A_0} C/IC) \\
& =
\text{Sym}_{A_0}^*(Q_0) \otimes_{A_0} C/IC \\
& =
\text{Sym}_{A_0}^*(Q_0) \otimes_{A_0}
\text{Sym}_{A_0}^*(P_0) \\
& =
\text{Sym}_{A_0}^*(Q_0 \oplus P_0) \\
& =
\text{Sym}_{A_0}^*(A_0^{\oplus n}) \\
& =
A_0[x_1, \ldots, x_n]
\end{align*}
Choisissons $f_1, \ldots, f_n \in D$ dont les images sont
$x_1, \ldots, x_n$ dans $D/ID = A_0[x_1, \ldots, x_n]$. Posons
$A = D/(f_1, \ldots, f_n)$. Notons que $A_0 = A/IA$. Nous affirmons que
$R \to A$ est syntomique dans un voisinage de $V(IA)$. Si cette affirmation est
vraie, nous pouvons trouver un $f \in A$ d'image $1 \in A_0$ tel que $A_f$ soit
syntomique sur $R$, ce qui achève la démonstration de (1).

\medskip\noindent
Démonstration de l'affirmation. Observons que $R \to D$ est syntomique en tant
que composée du morphisme d'anneaux syntomique $R \to C$, du morphisme
d'anneaux \'etale $C \to C'$ et du morphisme d'anneaux lisse $C' \to D$
(Algèbre, Lemmes
\ref{algebra-lemma-composition-syntomic} et
\ref{algebra-lemma-smooth-syntomic}).
La question est locale sur $\Spec(D)$ ; nous pouvons donc supposer que $D$ est
une intersection complète globale relative
(Algèbre, Lemme \ref{algebra-lemma-syntomic}).
Écrivons $D = R[y_1, \ldots, y_m]/(g_1, \ldots, g_s)$.
Soient $f'_1, \ldots, f'_n \in R[y_1, \ldots, y_m]$ des relèvements de
$f_1, \ldots, f_n$. Nous pouvons alors appliquer
Algèbre, Lemme \ref{algebra-lemma-localize-relative-complete-intersection}
pour obtenir l'affirmation.

\medskip\noindent
Démonstration de (2). Comme tout morphisme d'anneaux lisse est syntomique, nous
pouvons trouver un morphisme d'anneaux syntomique $R \to A$ tel que
$A_0 = A/IA$. Par hypothèse, les fibres de $R \to A$ sont lisses au-dessus des
idéaux premiers de $V(I)$ ; ainsi, $R \to A$ est lisse dans un voisinage ouvert
de $V(IA)$
(Algèbre, Lemme \ref{algebra-lemma-flat-fibre-smooth}).
Nous pouvons donc remplacer $A$ par une localisation pour obtenir le résultat
voulu.
\end{proof}

\noindent
Nous savons que tout morphisme d'anneaux syntomique $R \to A$ est localement
une intersection complète globale relative, voir
Algèbre, Lemme \ref{algebra-lemma-syntomic}.
Le lemme suivant affirme qu'un fibré vectoriel sur $\Spec(A)$ est une
intersection complète globale relative.

\begin{lemma}
\label{lemma-syntomic-complete-intersection}
Soit $R \to A$ un morphisme d'anneaux syntomique. Il existe alors un morphisme
lisse de $R$-algèbres $A \to C$ admettant une rétraction tel que $C$ soit une
intersection complète globale relative sur $R$, c'est-à-dire
$$
C \cong R[x_1, \ldots, x_n]/(f_1, \ldots, f_c)
$$
plate sur $R$ et dont toutes les fibres sont de dimension $n - c$.
\end{lemma}

\begin{proof}
Appliquons le lemme \ref{lemma-improve-presentation} pour obtenir $A \to C$.
D'après Algèbre, Lemme \ref{algebra-lemma-huber}, nous pouvons écrire
$C = R[x_1, \ldots, x_n]/(f_1, \ldots, f_c)$ de sorte que les images des $f_i$
forment une base de $J/J^2$.
Le morphisme d'anneaux $R \to C$ est syntomique (donc plat), car il est la
composée d'un morphisme d'anneaux syntomique et d'un morphisme d'anneaux lisse.
La dimension des fibres est $n - c$ d'après
Algèbre, Lemme \ref{algebra-lemma-lci}
(les fibres sont des intersections complètes locales, le lemme s'applique donc).
\end{proof}

\begin{lemma}
\label{lemma-smooth-standard-smooth}
Soit $R \to A$ un morphisme d'anneaux lisse. Il existe alors un morphisme lisse
de $R$-algèbres $A \to B$ admettant une rétraction tel que $B$ soit standard
lisse sur $R$, c'est-à-dire
$$
B \cong R[x_1, \ldots, x_n]/(f_1, \ldots, f_c)
$$
et $\det(\partial f_j/\partial x_i)_{i, j = 1, \ldots, c}$ est inversible dans
$B$.
\end{lemma}

\begin{proof}
Appliquons le lemme \ref{lemma-syntomic-complete-intersection} pour obtenir un
morphisme lisse de $R$-algèbres $A \to C$ admettant une rétraction, tel que
$C = R[x_1, \ldots, x_n]/(f_1, \ldots, f_c)$ soit une intersection complète
globale relative sur $R$. Comme $C$ est lisse sur $R$, nous avons une suite
exacte courte
$$
0 \to
\bigoplus\nolimits_{j = 1, \ldots, c} C f_j \to
\bigoplus\nolimits_{i = 1, \ldots, n} C\text{d}x_i \to
\Omega_{C/R} \to 0
$$
Puisque $\Omega_{C/R}$ est un $C$-module projectif, cette suite est scindée.
Choisissons un inverse à gauche $t$ du premier morphisme. Écrivons
$t(\text{d}x_i) = \sum c_{ij} f_j$
de sorte que $\sum_i \frac{\partial f_j}{\partial x_i} c_{i\ell} = \delta_{j\ell}$
(symbole de Kronecker). Soit
$$
B' = C[y_1, \ldots, y_c] =
R[x_1, \ldots, x_n, y_1, \ldots, y_c]/(f_1, \ldots, f_c)
$$
Le morphisme de $R$-algèbres $C \to B'$ admet la rétraction qui envoie $y_j$
sur zéro. Nous affirmons que le morphisme
$$
R[z_1, \ldots, z_n] \longrightarrow B',\quad
z_i \longmapsto x_i - \sum\nolimits_j c_{ij} y_j
$$
est \'etale en tout point de l'image de $\Spec(C) \to \Spec(B')$.
Dans $\Omega_{B'/R[z_1, \ldots, z_n]}$, nous avons
$$
0 =
\text{d}f_j - \sum\nolimits_i \frac{\partial f_j}{\partial x_i} \text{d}z_i
\equiv
\sum\nolimits_{i, \ell}
\frac{\partial f_j}{\partial x_i} c_{i\ell} \text{d}y_\ell
\equiv
\text{d}y_j \bmod (y_1, \ldots, y_c)\Omega_{B'/R[z_1, \ldots, z_n]}
$$
Puisque $0 = \text{d}z_i = \text{d}x_i$ modulo
$\sum B'\text{d}y_j + (y_1, \ldots, y_c)\Omega_{B'/R[z_1, \ldots, z_n]}$,
nous en concluons que
$$
\Omega_{B'/R[z_1, \ldots, z_n]}/
(y_1, \ldots, y_c)\Omega_{B'/R[z_1, \ldots, z_n]} = 0.
$$
Comme $\Omega_{B'/R[z_1, \ldots, z_n]}$ est un $B'$-module de type fini, le
lemme de Nakayama montre qu'il existe un $g \in 1 + (y_1, \ldots, y_c)$ tel que
$(\Omega_{B'/R[z_1, \ldots, z_n]})_g = 0$. Cela démontre que
$R[z_1, \ldots, z_n] \to B'_g$ est non ramifié, voir
Algèbre, Définition \ref{algebra-definition-unramified}.
Pour tout morphisme d'anneaux $R \to k$, où $k$ est un corps, nous obtenons un
morphisme d'anneaux non ramifié
$k[z_1, \ldots, z_n] \to (B'_g) \otimes_R k$ entre $k$-algèbres lisses de
dimension $n$. Il en résulte que
$k[z_1, \ldots, z_n] \to (B'_g) \otimes_R k$ est plat d'après
Algèbre, Lemmes \ref{algebra-lemma-CM-over-regular-flat} et
\ref{algebra-lemma-characterize-smooth-kbar}. Par le crit\`ere
de platitude par fibre
(Algèbre, Lemme \ref{algebra-lemma-criterion-flatness-fibre}),
nous concluons que $R[z_1, \ldots, z_n] \to B'_g$ est plat.
Enfin, Algèbre, Lemme \ref{algebra-lemma-characterize-etale} implique que
$R[z_1, \ldots, z_n] \to B'_g$ est \'etale.
Posons $B = B'_g$. Notons que $C \to B$ est lisse et admet une rétraction ; il
en est donc de même de $A \to B$. De plus,
$R[z_1, \ldots, z_n] \to B$ est \'etale.
D'après Algèbre, Lemme \ref{algebra-lemma-etale-standard-smooth}, nous pouvons
écrire
$$
B = R[z_1, \ldots, z_n, w_1, \ldots, w_m]/(g_1, \ldots, g_m)
$$
avec $\det(\partial g_j/\partial w_i)$ inversible dans $B$.
Cela démontre le lemme.
\end{proof}

\begin{lemma}
\label{lemma-colimit-standard-smooth}
Soit $R \to \Lambda$ un morphisme d'anneaux. Si $\Lambda$ est une colimite
filtrante de $R$-algèbres lisses, alors $\Lambda$ est une colimite filtrante de
$R$-algèbres standard lisses.
\end{lemma}

\begin{proof}
Soit $A \to \Lambda$ un morphisme de $R$-algèbres, avec $A$ de présentation
finie sur $R$. D'après
Algèbre, Lemme \ref{algebra-lemma-when-colimit}
nous devons factoriser ce morphisme par une algèbre standard lisse, et nous
savons que nous pouvons le factoriser sous la forme $A \to B \to \Lambda$, où
$B$ est lisse sur $R$. Choisissons un morphisme de $R$-algèbres $B \to C$
admettant une rétraction $C \to B$, tel que $C$ soit standard lisse sur $R$, voir
Lemme \ref{lemma-smooth-standard-smooth}.
La factorisation cherchée est alors $A \to B \to C \to B \to \Lambda$.
\end{proof}

\begin{lemma}
\label{lemma-standard-smooth-include-generators}
Soit $R \to A$ un morphisme d'anneaux standard lisse.
Soit $E \subset A$ un sous-ensemble fini de cardinal $|E| = n$.
Il existe alors une présentation
$A = R[x_1, \ldots, x_{n + m}]/(f_1, \ldots, f_c)$ avec $c \geq n$, dans
laquelle $\det(\partial f_j/\partial x_i)_{i, j = 1, \ldots, c}$ est inversible
dans $A$ et telle que $E$ soit l'ensemble des classes, dans ce quotient, de
$x_1, \ldots, x_n$.
\end{lemma}

\begin{proof}
Choisissons une présentation $A = R[y_1, \ldots, y_m]/(g_1, \ldots, g_d)$
telle que l'image de $\det(\partial g_j/\partial y_i)_{i, j = 1, \ldots, d}$
soit inversible dans $A$. Choisissons une énumération
$E = \{a_1, \ldots, a_n\}$ et des éléments $h_i \in R[y_1, \ldots, y_m]$ dont
l'image dans $A$ est $a_i$. Considérons la présentation
$$
A = R[x_1, \ldots, x_n, y_1, \ldots, y_m]/
(x_1 - h_1, \ldots, x_n - h_n, g_1, \ldots, g_d)
$$
et posons $c = n + d$.
\end{proof}

\begin{lemma}
\label{lemma-compare-standard}
Soit $R \to A$ un morphisme d'anneaux de présentation finie.
Soit $a \in A$. Considérons les conditions suivantes sur $a$ :
\begin{enumerate}
\item $A_a$ est lisse sur $R$,
\item $A_a$ est lisse sur $R$ et $\Omega_{A_a/R}$ est stablement libre,
\item $A_a$ est lisse sur $R$ et $\Omega_{A_a/R}$ est libre,
\item $A_a$ est standard lisse sur $R$,
\item $a$ est strictement standard dans $A$ sur $R$,
\item $a$ est standard élémentaire dans $A$ sur $R$.
\end{enumerate}
Alors :
\begin{enumerate}
\item[(a)] (4) $\Rightarrow$ (3) $\Rightarrow$ (2) $\Rightarrow$ (1),
\item[(b)] (6) $\Rightarrow$ (5),
\item[(c)] (6) $\Rightarrow$ (4),
\item[(d)] (5) $\Rightarrow$ (2),
\item[(e)] (2) $\Rightarrow$ les éléments $a^e$, $e \geq e_0$, sont
strictement standard dans $A$ sur $R$,
\item[(f)] (4) $\Rightarrow$ les éléments $a^e$, $e \geq e_0$, sont standard
élémentaires dans $A$ sur $R$.
\end{enumerate}
\end{lemma}

\begin{proof}
D'après les définitions et
Algèbre, Lemme \ref{algebra-lemma-standard-smooth},
la partie (a) est claire. La partie (b) résulte immédiatement de
la définition \ref{definition-strictly-standard}.

\medskip\noindent
Démonstration de (c). Choisissons une présentation
$A = R[x_1, \ldots, x_n]/(f_1, \ldots, f_m)$ telle que
(\ref{equation-elementary-standard-one}) et
(\ref{equation-elementary-standard-two}) soient satisfaites.
Choisissons $h \in R[x_1, \ldots, x_n]$ d'image $a$. Alors
$$
A_a = R[x_0, x_1, \ldots, x_n]/(x_0h - 1, f_1, \ldots, f_m).
$$
Écrivons $J = (x_0h - 1, f_1, \ldots, f_m)$.
D'après (\ref{equation-elementary-standard-two}), nous voyons que le
$A_a$-module $J/J^2$ est engendré par $x_0h - 1, f_1, \ldots, f_c$ sur
$A_a$. Ainsi, comme dans la démonstration de
Algèbre, Lemme \ref{algebra-lemma-huber}, nous pouvons choisir un $g \in 1 + J$
tel que
$$
A_a = R[x_0, \ldots, x_n, x_{n + 1}]/
(x_0h - 1, f_1, \ldots, f_c, gx_{n + 1} - 1).
$$
À ce stade, (\ref{equation-elementary-standard-one}) implique que $R \to A_a$
est standard lisse (utiliser les coordonnées
$x_0, x_1, \ldots, x_c, x_{n + 1}$ pour calculer les dérivées).

\medskip\noindent
Démonstration de (d). Choisissons une présentation
$A = R[x_1, \ldots, x_n]/(f_1, \ldots, f_m)$ telle que
(\ref{equation-strictly-standard-one}) et
(\ref{equation-strictly-standard-two}) soient satisfaites.
Écrivons $I = (f_1, \ldots, f_m)$.
Nous savons déjà que $A_a$ est lisse sur $R$, voir
Lemme \ref{lemma-elkik}.
D'après le lemme \ref{lemma-parse-equation-strictly-standard-one}, nous voyons que
$(I/I^2)_a$ est libre de base $f_1, \ldots, f_c$ et s'envoie isomorphiquement
sur un facteur direct de $\bigoplus A_a \text{d}x_i$. Puisque
$\Omega_{A_a/R} = (\Omega_{A/R})_a$
est le conoyau du morphisme
$(I/I^2)_a \to \bigoplus A_a \text{d}x_i$,
nous concluons qu'il est stablement libre.

\medskip\noindent
Démonstration de (e). Choisissons une présentation
$A = R[x_1, \ldots, x_n]/I$ avec $I$ de type fini.
Par hypothèse, nous avons une suite exacte courte
$$
0 \to (I/I^2)_a \to \bigoplus\nolimits_{i = 1, \ldots, n} A_a\text{d}x_i \to
\Omega_{A_a/R} \to 0
$$
qui est exacte scindée. Nous voyons donc que
$(I/I^2)_a \oplus \Omega_{A_a/R}$ est un $A_a$-module libre.
Puisque $\Omega_{A_a/R}$ est stablement libre, nous voyons que $(I/I^2)_a$
l'est également. Ainsi, en remplaçant la présentation choisie ci-dessus par
$A = R[x_1, \ldots, x_n, x_{n + 1}, \ldots, x_{n + r}]/J$, avec
$J = (I, x_{n + 1}, \ldots, x_{n + r})$ pour un certain $r$, nous obtenons que
$(J/J^2)_a$ est libre de type fini. Choisissons
$f_1, \ldots, f_c \in J$ dont les images forment une base de $(J/J^2)_a$.
Complétons cette famille en une famille de générateurs
$f_1, \ldots, f_m \in J$. Considérons la présentation
$A = R[x_1, \ldots, x_{n + r}]/(f_1, \ldots, f_m)$.
Alors, par construction, (\ref{equation-strictly-standard-two}) est satisfaite
pour $a^e$ dès que $e$ est suffisamment grand. De plus, puisque
$(J/J^2)_a \to \bigoplus\nolimits_{i = 1, \ldots, n + r} A_a\text{d}x_i$
est une injection scindée, nous pouvons trouver un inverse à gauche
$A_a$-linéaire. En exprimant cet inverse à gauche dans la base
$f_1, \ldots, f_c$ et en chassant les dénominateurs, nous trouvons un morphisme
linéaire
$\psi_0 : A^{\oplus n + r} \to A^{\oplus c}$ tel que
$$
A^{\oplus c} \xrightarrow{(f_1, \ldots, f_c)}
J/J^2 \xrightarrow{f \mapsto \text{d}f}
\bigoplus\nolimits_{i = 1, \ldots, n + r} A \text{d}x_i
\xrightarrow{\psi_0}
A^{\oplus c}
$$
soit la multiplication par $a^{e_0}$ pour un certain $e_0 \geq 1$. D'après
Lemme \ref{lemma-parse-equation-strictly-standard-one},
nous voyons que (\ref{equation-strictly-standard-one}) est satisfaite pour
$a^{ce_0}$, et donc pour $a^e$ pour tout $e$ tel que $e \geq ce_0$.

\medskip\noindent
Démonstration de (f). Choisissons une présentation
$A_a = R[x_1, \ldots, x_n]/(f_1, \ldots, f_c)$ telle que
$\det(\partial f_j/\partial x_i)_{i, j = 1, \ldots, c}$
soit inversible dans $A_a$. Nous pouvons supposer que, pour un certain $m < n$,
les classes des éléments $x_1, \ldots, x_m$ correspondent aux $a_i/1$, où
$a_1, \ldots, a_m \in A$ sont des générateurs de $A$ sur $R$, voir
Lemme \ref{lemma-standard-smooth-include-generators}.
Après avoir remplacé $x_i$ par $a^Nx_i$ pour $m < i \leq n$, nous pouvons
supposer que la classe de $x_i$ est $a_i/1 \in A_a$ pour un certain $a_i \in A$.
Considérons le morphisme d'anneaux
$$
\Psi : R[x_1, \ldots, x_n] \longrightarrow A,\quad
x_i \longmapsto a_i.
$$
Il s'agit d'un morphisme d'anneaux surjectif. En remplaçant $f_j$ par $a^Nf_j$,
nous pouvons supposer que $f_j \in R[x_1, \ldots, x_n]$ et que
$\Psi(f_j) = 0$ (car, après tout, $f_j(a_1/1, \ldots, a_n/1) = 0$ dans $A_a$).
Soit $J = \Ker(\Psi)$. Alors $A = R[x_1, \ldots, x_n]/J$ est une présentation,
et $f_1, \ldots, f_c \in J$ sont des éléments tels que $(J/J^2)_a$ soit libre
de base $f_1, \ldots, f_c$ et que
$\det(\partial f_j/\partial x_i)_{i, j = 1, \ldots, c}$ s'envoie sur un élément
inversible de $A_a$. Il en résulte que
(\ref{equation-elementary-standard-one}) et
(\ref{equation-elementary-standard-two})
sont satisfaites pour $a^e$ et tout $e$ suffisamment grand, comme souhaité.
\end{proof}
\section{Intermède : désingularisation de Néron}
\label{section-neron}

\noindent
Nous interrompons l'étude du cas général du théorème de Popescu pour
examiner un cas plus facile mais déjà très intéressant, à savoir celui où
$R \to \Lambda$ est un homomorphisme d'anneaux de valuation discrète.
Ce cas est étudié dans
\cite[Section 4]{Artin-Algebraic-Approximation}.

\begin{situation}
\label{situation-neron}
Ici, $R \subset \Lambda$ est une extension d'anneaux de valuation discrète
d'indice de ramification $1$ (Compléments d'algèbre, définition
\ref{more-algebra-definition-extension-discrete-valuation-rings}).
Nous nous donnons une factorisation
$$
R \to A \xrightarrow{\varphi} \Lambda
$$
où $R \to A$ est plat et de type fini. Posons $\mathfrak q = \Ker(\varphi)$
et $\mathfrak p = \varphi^{-1}(\mathfrak m_\Lambda)$.
\end{situation}

\noindent
Dans la situation \ref{situation-neron}, soit $\pi \in R$ une uniformisante.
Rappelons que la platitude de $A$ sur $R$ signifie que $\pi$
est un non-diviseur de zéro dans $A$
(Compléments d'algèbre, lemme
\ref{more-algebra-lemma-valuation-ring-torsion-free-flat}).
Notre hypothèse sur $R \subset \Lambda$ montre que l'image de $\pi$
est une uniformisante de $\Lambda$. Comme $\pi \in \mathfrak p$, nous pouvons considérer
l'algèbre d'éclatement affine de Néron (voir
Algèbre, section \ref{algebra-section-blow-up})
$$
\varphi' :
A' = A[\textstyle{\frac{\mathfrak p}{\pi}}]
\longrightarrow
\Lambda
$$
munie du morphisme induit vers $\Lambda$ qui envoie
$a/\pi^n$, $a \in \mathfrak p^n$, sur $\pi^{-n}\varphi(a)$ dans $\Lambda$.
Nous noterons $\mathfrak q' \subset \mathfrak p' \subset A'$
les idéaux premiers correspondants de $A'$. Remarquons que la classe
d'isomorphisme de $A'$ ne dépend pas du choix de l'uniformisante.
En répétant cette construction, nous obtenons une suite
$$
A \to A' \to A'' \to \ldots \to \Lambda
$$

\begin{lemma}
\label{lemma-neron-functorial}
Dans la situation \ref{situation-neron}, l'éclatement de Néron est fonctoriel
au sens suivant :
\begin{enumerate}
\item si $a \in A$, $a \not \in \mathfrak p$, alors l'éclatement de Néron
de $A_a$ est $A'_a$, et
\item si $B \to A$ est une surjection de $R$-algèbres plates de type fini
de noyau $I$, alors $A'$ est le quotient de $B'/IB'$ par sa
torsion annulée par une puissance de $\pi$.
\end{enumerate}
\end{lemma}

\begin{proof}
Les assertions (1) et (2) sont des cas particuliers de
Algèbre, lemme \ref{algebra-lemma-blowup-base-change}.
En effet, dès que nous avons $A_1 \to A_2 \to \Lambda$ avec
$\mathfrak p_1 A_2 = \mathfrak p_2$, l'algèbre $A_2'$ est
le quotient de $A_1' \otimes_{A_1} A_2$ par sa torsion annulée par une puissance de $\pi$.
\end{proof}

\begin{lemma}
\label{lemma-neron-blowup-smooth}
Dans la situation \ref{situation-neron}, supposons que $R \to A$ soit lisse
en $\mathfrak p$ et que $R/\pi R \subset \Lambda/\pi \Lambda$
soit une extension séparable de corps. Alors $R \to A'$ est lisse en
$\mathfrak p'$ et il existe une suite exacte courte
$$
0 \to
\Omega_{A/R} \otimes_A A'_{\mathfrak p'} \to
\Omega_{A'/R, \mathfrak p'} \to
(A'/\pi A')_{\mathfrak p'}^{\oplus c} \to 0
$$
où $c = \dim((A/\pi A)_\mathfrak p)$.
\end{lemma}

\begin{proof}
D'après le lemme \ref{lemma-neron-functorial}, nous pouvons remplacer $A$ par
un localisé en un élément qui n'appartient pas à $\mathfrak p$ ; nous le ferons
désormais sans autre mention. Posons $\kappa = R/\pi R$. Comme la lissité est
stable par changement de base
(Algèbre, lemme \ref{algebra-lemma-base-change-smooth}),
nous voyons que $A/\pi A$ est lisse sur $\kappa$ en $\mathfrak p$.
Par conséquent, $(A/\pi A)_\mathfrak p$ est un anneau local régulier
(Algèbre, lemme \ref{algebra-lemma-characterize-smooth-over-field}).
Choisissons $g_1, \ldots, g_c \in \mathfrak p$ dont les images forment
un système régulier de paramètres de $(A/\pi A)_\mathfrak p$.
Nous voyons alors que $\mathfrak p = (\pi, g_1, \ldots, g_c)$,
quitte à remplacer $A$ par un localisé.
Remarquons que $\pi, g_1, \ldots, g_c$ est une suite régulière
dans $A_\mathfrak p$ (d'abord $\pi$ est un non-diviseur de zéro, puis on applique
Algèbre, lemme \ref{algebra-lemma-regular-ring-CM}
au reste de la suite).
Après avoir remplacé $A$ par un localisé, nous pouvons supposer que
$\pi, g_1, \ldots, g_c$ est une suite régulière dans $A$
(Algèbre, lemme \ref{algebra-lemma-regular-sequence-in-neighbourhood}).
Il s'ensuit que
$$
A' = A[y_1, \ldots, y_c]/(\pi y_1 - g_1, \ldots, \pi y_c - g_c) =
A[y_1, \ldots, y_c]/I
$$
d'après Compléments d'algèbre, lemme \ref{more-algebra-lemma-blowup-regular-sequence}.
Dans la suite, nous utiliserons sans autre mention la définition de la lissité
au moyen du complexe cotangent naïf (Algèbre, définition \ref{algebra-definition-smooth})
et le critère d'Algèbre, lemme \ref{algebra-lemma-smooth-at-point}.
La suite exacte d'Algèbre, lemme \ref{algebra-lemma-exact-sequence-NL},
pour $R \to A[y_1, \ldots, y_c] \to A'$, s'écrit
$$
0 \to H_1(\NL_{A'/R}) \to I/I^2 \to
\Omega_{A/R} \otimes_A A' \oplus
\bigoplus\nolimits_{i = 1, \ldots, c} A' \text{d}y_i \to
\Omega_{A'/R} \to 0
$$
où la classe de $\pi y_i - g_i$ dans $I/I^2$ est envoyée
sur $- \text{d}g_i + \pi \text{d}y_i$ dans le terme suivant.
Ici, nous avons utilisé Algèbre, lemme \ref{algebra-lemma-NL-surjection},
pour calculer $\NL_{A'/A[y_1, \ldots, y_c]}$, ainsi que le fait que
$R \to A[y_1, \ldots, y_c]$ est lisse, de sorte que
$H_1(\NL_{A[y_1, \ldots, y_c]/R}) = 0$ et
$\Omega_{A[y_1, \ldots, y_c]/R}$ est un
$A[y_1, \ldots, y_c]$-module projectif fini (donc, a fortiori, plat),
qui est en fait la somme directe de
$\Omega_{A/R} \otimes_A A[y_1, \ldots, y_c]$ et d'un module libre
de base les $\text{d}y_i$. Pour achever la démonstration, il suffit de
montrer que $\text{d}g_1, \ldots, \text{d}g_c$ font partie d'une base du
module libre fini $\Omega_{A/R, \mathfrak p}$.
En effet, cela montrera que $(I/I^2)_{\mathfrak p'}$ est libre sur les
$\pi y_i - g_i$, que le localisé en $\mathfrak p'$ du morphisme central
de la suite est injectif, donc que $H_1(\NL_{A'/R})_{\mathfrak p'} = 0$, et que
le conoyau $\Omega_{A'/R, \mathfrak p'}$ est libre fini.
Pour cela, il suffit de montrer que les images des $\text{d}g_i$ sont
linéairement indépendantes sur $\kappa(\mathfrak p)$ dans
$\Omega_{A/R, \mathfrak p}/\pi = \Omega_{(A/\pi A)/\kappa, \mathfrak p}$
(l'égalité résulte d'Algèbre, lemme \ref{algebra-lemma-differentials-base-change}).
Comme $\kappa \subset \kappa(\mathfrak p) \subset \Lambda/\pi \Lambda$,
nous voyons que $\kappa(\mathfrak p)$ est séparable sur $\kappa$
(Algèbre, définition \ref{algebra-definition-separable-field-extension}).
L'indépendance linéaire recherchée résulte alors d'Algèbre,
lemme \ref{algebra-lemma-computation-differential}.
\end{proof}

\begin{lemma}
\label{lemma-neron-when-smooth}
Dans la situation \ref{situation-neron}, supposons que $R \to A$ soit lisse
en $\mathfrak q$ et que nous ayons une surjection de $R$-algèbres
$B \to A$ de noyau $I$. Supposons $R \to B$ lisse en
$\mathfrak p_B = (B \to A)^{-1}\mathfrak p$. Si le conoyau de
$$
I/I^2 \otimes_A \Lambda \to \Omega_{B/R} \otimes_B \Lambda
$$
est un $\Lambda$-module libre, alors $R \to A$ est lisse en $\mathfrak p$.
\end{lemma}

\begin{proof}
Le conoyau du morphisme $I/I^2 \to \Omega_{B/R} \otimes_B A$
est $\Omega_{A/R}$, voir Algèbre, lemme \ref{algebra-lemma-differential-seq}.
Soit $d = \dim_\mathfrak q(A/R)$ la dimension relative de $R \to A$
en $\mathfrak q$, c'est-à-dire la dimension de $\Spec(A[1/\pi])$ en $\mathfrak q$.
Voir Algèbre, définition \ref{algebra-definition-relative-dimension}.
Alors $\Omega_{A/R, \mathfrak q}$ est libre sur $A_\mathfrak q$ et de rang $d$
(Algèbre, lemme \ref{algebra-lemma-characterize-smooth-over-field}).
Ainsi, si l'hypothèse du lemme est satisfaite,
$\Omega_{A/R} \otimes_A \Lambda$ est libre de rang $d$.
Il s'ensuit que $\Omega_{A/R} \otimes_A \kappa(\mathfrak p)$
est de dimension $d$ (car c'est le cas après tensorisation par $\Lambda/\pi \Lambda$).
Comme $R \to A$ est plat et que $\mathfrak p$
est une spécialisation de $\mathfrak q$, nous voyons que
$\dim_\mathfrak p(A/R) \geq d$ d'après Algèbre, lemme
\ref{algebra-lemma-dimension-fibres-bounded-open-upstairs}.
Il en résulte que $R \to A$ est lisse en $\mathfrak p$ d'après
Algèbre, lemmes \ref{algebra-lemma-flat-fibre-smooth} et
\ref{algebra-lemma-characterize-smooth-over-field}.
\end{proof}

\begin{lemma}
\label{lemma-neron-desingularization}
Dans la situation \ref{situation-neron},
supposons que $R \to A$ soit lisse en $\mathfrak q$
et que $R/\pi R \subset \Lambda/\pi \Lambda$ soit une extension
séparable de corps. Après un nombre fini d'éclatements affines de Néron,
l'algèbre $A$ devient alors lisse sur $R$ en $\mathfrak p$.
\end{lemma}

\begin{proof}
Choisissons une $R$-algèbre $B$ et une surjection $B \to A$. Posons
$\mathfrak p_B = (B \to A)^{-1}(\mathfrak p)$ et notons $r$
la dimension relative de $R \to B$ en $\mathfrak p_B$. Nous choisissons $B$
de sorte que $R \to B$ soit lisse en $\mathfrak p_B$.
Par exemple, nous pouvons prendre pour $B$ une algèbre de polynômes en $r$
variables sur $R$. Considérons le complexe
$$
I/I^2 \otimes_A \Lambda \longrightarrow \Omega_{B/R} \otimes_B \Lambda
$$
du lemme \ref{lemma-neron-when-smooth}. La structure des modules finis sur
$\Lambda$ (Compléments d'algèbre, lemme \ref{more-algebra-lemma-modules-PID})
montre que le conoyau est de la forme
$$
\Lambda^{\oplus d} \oplus
\bigoplus\nolimits_{i = 1, \ldots, n} \Lambda/\pi^{e_i} \Lambda
$$
pour certains $d \geq 0$, $n \geq 0$ et $e_i \geq 1$. Remarquons que $d$
est la dimension relative de $A/R$ en $\mathfrak q$
(Algèbre, lemme \ref{algebra-lemma-characterize-smooth-over-field}).
Si le défaut $e = \sum_{i = 1, \ldots, n} e_i$ est nul, la conclusion résulte du
lemme \ref{lemma-neron-when-smooth}.

\medskip\noindent
Examinons maintenant ce qui se passe lorsque nous effectuons l'éclatement de Néron.
Rappelons que $A'$ est le quotient de $B'/IB'$ par sa torsion annulée par
une puissance de $\pi$ (lemme \ref{lemma-neron-functorial}) et que $R \to B'$ est lisse en
$\mathfrak p_{B'}$ (lemme \ref{lemma-neron-blowup-smooth}).
Après éclatement, nous nous retrouvons donc exactement dans la même situation. Représentons-la par
$$
\xymatrix{
0 \ar[r] & I' \ar[r] & B' \ar[r] & A' \ar[r] & 0 \\
0 \ar[r] & I \ar[u] \ar[r] & B \ar[u] \ar[r] & A \ar[r] \ar[u] & 0
}
$$
Comme $I \subset \mathfrak p_B$, nous voyons que $I \to I'$
se factorise par $\pi I'$. En considérant le morphisme induit de
complexes, nous obtenons
$$
\xymatrix{
I'/(I')^2 \otimes_{A'} \Lambda \ar[r] &
\Omega_{B'/R} \otimes_{B'} \Lambda \ar@{=}[r] & M' \\
I/I^2 \otimes_A \Lambda \ar[r] \ar[u] &
\Omega_{B/R} \otimes_B \Lambda \ar[u] \ar@{=}[r] & M
}
$$
Alors $M \subset M'$ sont des $\Lambda$-modules libres finis, et le quotient
$M'/M$ est annulé par $\pi$, voir le lemme \ref{lemma-neron-blowup-smooth}.
Soient $N \subset M$ et $N' \subset M'$ les images des morphismes horizontaux,
et posons $Q = M/N$ et $Q' = M'/N'$.
Nous obtenons un diagramme commutatif
$$
\xymatrix{
0 \ar[r] &
N' \ar[r] &
M' \ar[r] &
Q' \ar[r] &
0 \\
0 \ar[r] &
N \ar[r] \ar[u] &
M \ar[r] \ar[u] &
Q \ar[r] \ar[u] &
0
}
$$
Alors $N \subset N'$ sont des $\Lambda$-modules libres de rang $r - d$.
Comme $I$ s'envoie dans $\pi I'$, nous voyons que $N \subset \pi N'$.

\medskip\noindent
Soit $K = \Lambda_\pi$ le corps des fractions de $\Lambda$.
Nous avons un diagramme commutatif
$$
\xymatrix{
0 \ar[r] &
N' \ar[r] &
N'_K \cap M' \ar[r] &
Q'_{tor} \ar[r] &
0 \\
0 \ar[r] &
N \ar[r] \ar[u] &
N_K \cap M \ar[r] \ar[u] &
Q_{tor} \ar[r] \ar[u] &
0
}
$$
dont les lignes sont des suites exactes courtes. Cela montre que la variation du défaut
est donnée par
$$
e - e' =
\text{longueur}(Q_{tor}) - \text{longueur}(Q'_{tor})
=
\text{longueur}(N'/N) - \text{longueur}(N'_K \cap M' / N_K \cap M)
$$
Comme $M'/M$ est annulé par $\pi$, il en va de même de $N'_K \cap M' / N_K \cap M$,
et sa longueur est au plus $\dim_K(N_K)$.
Comme $N \subset \pi N'$, nous obtenons $\text{longueur}(N'/N) \ge \dim_K(N_K)$,
avec égalité si et seulement si $N = \pi N'$.

\medskip\noindent
Pour achever la démonstration, nous devons montrer que $N$ est strictement contenu
dans $\pi N'$ lorsque $A$ n'est pas lisse en $\mathfrak p$ ;
c'est le calcul clé de l'argument de Néron.
Pour cela, considérons la suite exacte
$$
I/I^2 \otimes_B \kappa(\mathfrak p_B)
\to \Omega_{B/R} \otimes_B \kappa(\mathfrak p_B)
\to \Omega_{A/R} \otimes_A \kappa(\mathfrak p) \to 0
$$
(qui résulte d'Algèbre, lemme \ref{algebra-lemma-differential-seq}).
Comme $R \to A$ n'est pas lisse en $\mathfrak p$, nous voyons que la dimension $s$ de
$\Omega_{A/R} \otimes_A \kappa(\mathfrak p)$
est strictement supérieure à $d$. D'autre part,
la première flèche se factorise par le morphisme injectif
$$
\mathfrak p_B B_{\mathfrak p_B}/\mathfrak p_B^2 B_{\mathfrak p_B}
\to \Omega_{B/R} \otimes_B \kappa(\mathfrak p_B)
$$
d'Algèbre, lemme \ref{algebra-lemma-computation-differential} ;
remarquons que $\kappa(\mathfrak p)$ est séparable sur $\kappa$
d'après notre hypothèse sur $R/\pi R \subset \Lambda/\pi \Lambda$.
Nous en concluons qu'il existe des générateurs
$g_1, \ldots, g_t \in I$ tels que $g_j \in \mathfrak p^2$
pour $j > r - s$. Les images des $g_j$ dans $A'$ appartiennent alors à $\pi^2 I'$
pour $j > r - s$. Comme $r - s < r - d$,
au moins un des générateurs minimaux
de $N$ devient divisible par $\pi^2$ dans $N'$.
Ainsi, $e$ diminue d'au moins $1$, ce qui conclut la démonstration.
\end{proof}

\noindent
Si $R \to \Lambda$ est une extension d'anneaux de valuation discrète,
alors $R \to \Lambda$ est régulier si et seulement si
(a) l'indice de ramification vaut $1$,
(b) l'extension des corps de fractions est séparable, et
(c) $R/\mathfrak m_R \subset \Lambda/\mathfrak m_\Lambda$
est séparable. Le résultat suivant est donc un cas particulier
de la désingularisation générale de Néron donnée dans le
théorème \ref{theorem-popescu}.

\begin{lemma}
\label{lemma-neron-colimit}
\begin{slogan}
Les extensions non ramifiées d'anneaux de valuation discrète sont ind-lisses : désingularisation de Néron
\end{slogan}
Soit $R \subset \Lambda$ une extension d'anneaux de valuation
discrète d'indice de ramification $1$ qui induit une extension séparable
des corps résiduels et des corps de fractions.
Alors $\Lambda$ est une colimite filtrante de $R$-algèbres lisses.
\end{lemma}

\begin{proof}
D'après Algèbre, lemme \ref{algebra-lemma-when-colimit}, il suffit de montrer
que tout $R \to A \to \Lambda$ comme dans la situation \ref{situation-neron}
se factorise sous la forme $A \to B \to \Lambda$, où $B$ est une
$R$-algèbre lisse. Après avoir remplacé $A$ par son image dans $\Lambda$,
nous pouvons supposer que $A$ est un domaine dont le corps des fractions $K$
est un sous-corps du corps des fractions de $\Lambda$.
En particulier, $K$ est séparable sur le corps des fractions de $R$
d'après nos hypothèses. Alors $R \to A$ est lisse en $\mathfrak q = (0)$ d'après
Algèbre, lemme \ref{algebra-lemma-smooth-at-generic-point}.
Après un nombre fini d'éclatements de Néron, nous pouvons supposer que $R \to A$
est lisse en $\mathfrak p$, voir le lemme \ref{lemma-neron-desingularization}.
Enfin, après avoir remplacé $A$ par son localisé
en un élément $a \in A$, $a \not \in \mathfrak p$, elle devient
lisse sur $R$, ce qui démontre le lemme.
\end{proof}












\section{Le problème de relèvement}
\label{section-lifting}

\noindent
Le but de cette section est de démontrer (proposition \ref{proposition-lift})
que la classe des algèbres qui sont des colimites filtrantes d'algèbres lisses
est stable par déformations infinitésimales plates. La démonstration est élémentaire
et n'utilise que les résultats sur les présentations d'algèbres lisses de la
section \ref{section-presentations}.

\begin{lemma}
\label{lemma-lift-once}
Soit $R \to \Lambda$ un homomorphisme d'anneaux. Soit $I \subset R$ un idéal.
Supposons que
\begin{enumerate}
\item $I^2 = 0$, et
\item $\Lambda/I\Lambda$ soit une colimite filtrante de $R/I$-algèbres lisses.
\end{enumerate}
Soit $\varphi : A \to \Lambda$ un morphisme de $R$-algèbres, où $A$ est de
présentation finie sur $R$. Il existe alors une factorisation
$$
A \to B/J \to \Lambda
$$
où $B$ est une $R$-algèbre lisse et $J \subset IB$ un idéal de type fini.
\end{lemma}

\begin{proof}
Choisissons une factorisation
$$
A/IA \to \bar B \to \Lambda/I\Lambda
$$
où $\bar B$ est standard lisse sur $R/I$ ; cela est possible d'après
l'hypothèse et le lemme \ref{lemma-colimit-standard-smooth}. Écrivons
$$
\bar B = A/IA[t_1, \ldots, t_r]/(\bar g_1, \ldots, \bar g_s)
$$
et supposons que $\bar B \to \Lambda/I\Lambda$ envoie $t_i$ sur la classe
de $\lambda_i$ modulo $I\Lambda$. Choisissons
$g_1, \ldots, g_s \in A[t_1, \ldots, t_r]$ relevant
$\bar g_1, \ldots, \bar g_s$. Écrivons
$\varphi(g_i)(\lambda_1, \ldots, \lambda_r) =
\sum \epsilon_{ij} \mu_{ij}$
pour certains $\epsilon_{ij} \in I$ et $\mu_{ij} \in \Lambda$. Définissons
$$
A' = A[t_1, \ldots, t_r, \delta_{i, j}]/
(g_i - \sum \epsilon_{ij} \delta_{ij})
$$
et considérons le morphisme
$$
A' \longrightarrow \Lambda,\quad
a \longmapsto \varphi(a),\quad
t_i \longmapsto \lambda_i,\quad
\delta_{ij} \longmapsto \mu_{ij}
$$
Nous avons
$$
A'/IA' = A/IA[t_1, \ldots, t_r]/(\bar g_1, \ldots, \bar g_s)[\delta_{ij}]
\cong \bar B[\delta_{ij}]
$$
C'est une algèbre standard lisse sur $R/I$, puisque $\bar B$ est standard
lisse. Choisissons une présentation
$A'/IA' = R/I[x_1, \ldots, x_n]/(\bar f_1, \ldots, \bar f_c)$ telle que
$\det(\partial \bar f_j/\partial x_i)_{i, j = 1, \ldots, c}$ soit inversible dans
$A'/IA'$. Choisissons des relèvements $f_1, \ldots, f_c \in R[x_1, \ldots, x_n]$ de
$\bar f_1, \ldots, \bar f_c$. Alors
$$
B = R[x_1, \ldots, x_n, x_{n + 1}]/
(f_1, \ldots, f_c,
x_{n + 1}\det(\partial f_j/\partial x_i)_{i, j = 1, \ldots, c} - 1)
$$
est lisse sur $R$. Comme les morphismes d'anneaux lisses sont formellement lisses
(Algèbre, proposition \ref{algebra-proposition-smooth-formally-smooth}),
il existe un morphisme de $R$-algèbres $B \to A'$ qui est un isomorphisme
modulo $I$. Le morphisme $B \to A'$ est alors surjectif par le lemme de Nakayama
(Algèbre, lemme \ref{algebra-lemma-NAK}).
Ainsi $A' = B/J$, où $J \subset IB$ est de type fini (voir
Algèbre, lemme \ref{algebra-lemma-finite-presentation-independent}).
\end{proof}

\begin{lemma}
\label{lemma-lift-twice}
Soit $R \to \Lambda$ un homomorphisme d'anneaux. Soit $I \subset R$ un idéal.
Supposons que
\begin{enumerate}
\item $I^2 = 0$,
\item $\Lambda/I\Lambda$ soit une colimite filtrante de $R/I$-algèbres lisses, et
\item $R \to \Lambda$ soit plat.
\end{enumerate}
Soit $\varphi : B \to \Lambda$ un morphisme de $R$-algèbres, où $B$
est lisse sur $R$. Soit $J \subset IB$ un idéal de type fini
tel que $\varphi(J) = 0$.
Il existe alors des morphismes de $R$-algèbres
$$
B \xrightarrow{\alpha} B' \xrightarrow{\beta} \Lambda
$$
tels que $B'$ soit lisse sur $R$, que $\alpha(J) = 0$ et
que $\beta \circ \alpha = \varphi$.
\end{lemma}

\begin{proof}
Si nous savons démontrer le lemme dans le cas $J = (h)$, nous pouvons le
démontrer par récurrence sur le nombre de générateurs de $J$. Supposons en effet
que $J$ puisse être engendré par $n$ éléments $h_1, \ldots, h_n$ et que le
lemme soit vrai dans tous les cas où $J$ est engendré par $n - 1$ éléments.
Nous appliquons d'abord le cas $n = 1$ pour construire $B \to B' \to \Lambda$
tel que le premier morphisme annule $h_n$. Soit ensuite $J'$ l'idéal de $B'$
engendré par les images de $h_1, \ldots, h_{n - 1}$ ; nous appliquons le cas
$n - 1$ pour construire $B' \to B'' \to \Lambda$.
Il est immédiat de vérifier que $B \to B'' \to \Lambda$ convient.

\medskip\noindent
Supposons $J = (h)$ et écrivons $h = \sum \epsilon_i b_i$
pour certains $\epsilon_i \in I$ et $b_i \in B$. Remarquons que
$0 = \varphi(h) = \sum \epsilon_i \varphi(b_i)$.
Comme $\Lambda$ est plat sur $R$, le critère équationnel de platitude
(Algèbre, lemme \ref{algebra-lemma-flat-eq}) implique qu'il existe
$\lambda_j \in \Lambda$, $j = 1, \ldots, m$, et $a_{ij} \in R$ tels que
$\varphi(b_i) = \sum_j a_{ij} \lambda_j$ et $\sum_i \epsilon_i a_{ij} = 0$.
Posons
$$
C = B[x_1, \ldots, x_m]/(b_i - \sum a_{ij} x_j)
$$
où $C \to \Lambda$ est défini par $\varphi$ et $x_j \mapsto \lambda_j$.
Choisissons une factorisation
$$
C \to B'/J' \to \Lambda
$$
comme dans le lemme \ref{lemma-lift-once}. Comme $B$ est lisse sur $R$, nous pouvons
relever le morphisme $B \to C \to B'/J'$ en un morphisme $\alpha : B \to B'$. Alors
$\varphi = \beta \circ \alpha$. Pour achever la démonstration, vérifions que
$\alpha(h) = 0$. En effet, le fait que $\alpha$ relève
$B \to C \to B'/J'$ implique que
$$
\alpha(b_i) = \sum a_{ij} \xi_j + \theta_i
$$
pour certains $\xi_j \in B'$ et $\theta_i \in J' \subset IB'$.
Par conséquent,
$$
\alpha(h) = \alpha(\sum \epsilon_i b_i) =
\sum \epsilon_i a_{ij} \xi_j + \sum \epsilon_i \theta_i = 0
$$
en vertu des relations ci-dessus et du fait que $I^2 = 0$.
\end{proof}

\begin{proposition}
\label{proposition-lift}
\begin{slogan}
L'ind-lissité d'une algèbre est stable par déformations infinitésimales
\end{slogan}
Soit $R \to \Lambda$ un homomorphisme d'anneaux. Soit $I \subset R$ un idéal.
Supposons que
\begin{enumerate}
\item $I$ soit nilpotent,
\item $\Lambda/I\Lambda$ soit une colimite filtrante de $R/I$-algèbres lisses, et
\item $R \to \Lambda$ soit plat.
\end{enumerate}
Alors $\Lambda$ est une colimite filtrante de $R$-algèbres lisses.
\end{proposition}

\begin{proof}
Comme $I^n = 0$ pour un certain $n$, une récurrence sur $n$ montre qu'il
suffit de considérer le cas où $I^2 = 0$. Soit
$\varphi : A \to \Lambda$ un morphisme de $R$-algèbres, où $A$ est de
présentation finie sur $R$. Nous devons trouver une factorisation $A \to B \to \Lambda$
avec $B$ lisse sur $R$, voir Algèbre, lemme \ref{algebra-lemma-when-colimit}.
D'après le lemme \ref{lemma-lift-once}, nous pouvons supposer que
$A = B/J$, où $B$ est lisse sur $R$ et $J \subset IB$
est un idéal de type fini. D'après le
lemme \ref{lemma-lift-twice},
il existe un diagramme commutatif
$$
\xymatrix{
B \ar[rr]_\alpha \ar[rd]_\varphi & & B' \ar[ld]^\beta \\
& \Lambda
}
$$
de $R$-algèbres, où $B'$ est lisse sur $R$ et $\alpha(J) = 0$.
Ainsi, $\alpha$ se factorise sous la forme $B \to A \to B'$, ce qui achève la démonstration.
\end{proof}






\section{Le lemme de relèvement}
\label{section-lifting-lemma}

\noindent
Voici un lemme diaboliquement ingénieux.

\begin{lemma}
\label{lemma-lifting}
Soit $R$ un anneau noethérien et soit $\Lambda$ une $R$-algèbre.
Soit $\pi \in R$ et supposons que $\text{Ann}_R(\pi) = \text{Ann}_R(\pi^2)$ et
$\text{Ann}_\Lambda(\pi) = \text{Ann}_\Lambda(\pi^2)$.
Supposons donnés des morphismes de $R$-algèbres
$R/\pi^2R \to \bar C \to \Lambda/\pi^2\Lambda$
avec $\bar C$ de présentation finie.
Il existe alors un homomorphisme de $R$-algèbres
$D \to \Lambda$ et un diagramme commutatif
$$
\xymatrix{
R/\pi^2R \ar[r] \ar[d] &
\bar C \ar[r] \ar[d] &
\Lambda/\pi^2\Lambda \ar[d] \\
R/\pi R \ar[r] &
D/\pi D \ar[r] &
\Lambda/\pi \Lambda
}
$$
ayant les propriétés suivantes :
\begin{enumerate}
\item[(a)] $D$ est de présentation finie,
\item[(b)] $R \to D$ est lisse en tout idéal premier $\mathfrak q$ tel que
$\pi \not \in \mathfrak q$,
\item[(c)] $R \to D$ est lisse en tout idéal premier $\mathfrak q$ tel que
$\pi \in \mathfrak q$ et situé au-dessus d'un idéal premier de $\bar C$ où
$R/\pi^2 R \to \bar C$ est lisse, et
\item[(d)] $\bar C/\pi \bar C \to D/\pi D$ est lisse en tout idéal premier
situé au-dessus d'un idéal premier de $\bar C$ où $R/\pi^2R \to \bar C$ est lisse.
\end{enumerate}
\end{lemma}

\begin{proof}
Choisissons une présentation
$$
\bar C = R[x_1, \ldots, x_n]/(f_1, \ldots, f_m)
$$
Notons également $I = (f_1, \ldots, f_m)$ et $\bar I$ l'image de
$I$ dans $R/\pi^2R[x_1, \ldots, x_n]$. Puisque $R$ est noethérien,
$\bar C$ l'est aussi. Le lieu lisse de $R/\pi^2 R \to \bar C$
est donc quasi-compact, voir
Topologie, lemme \ref{topology-lemma-Noetherian}.
En appliquant le
lemme \ref{lemma-find-strictly-standard},
nous pouvons choisir une famille finie d'éléments
$a_1, \ldots, a_r \in R[x_1, \ldots, x_n]$ telle que
\begin{enumerate}
\item la réunion des sous-espaces ouverts
$\Spec(\bar C_{a_k}) \subset \Spec(\bar C)$
recouvre le lieu lisse de $R/\pi^2 R \to \bar C$, et
\item pour chaque $k = 1, \ldots, r$, il existe une partie finie
$E_k \subset \{1, \ldots, m\}$ telle que
$(\bar I/\bar I^2)_{a_k}$ soit librement engendré par les classes des
$f_j$, $j \in E_k$.
\end{enumerate}
Posons $I_k = (f_j, j \in E_k) \subset I$ et notons $\bar I_k$
l'image de $I_k$ dans $R/\pi^2R[x_1, \ldots, x_n]$.
D'après (2) et le lemme de Nakayama, $(\bar I/\bar I_k)_{a_k}$
est annulé par $1 + b'_k$ pour un certain $b'_k \in \bar I_{a_k}$.
Supposons que $b'_k$ soit l'image de $b_k/(a_k)^N$ pour un certain $b_k \in I$
et un certain entier $N$. En remplaçant $a_k$ par $(a_k)^N + b_k$, nous obtenons
\begin{enumerate}
\item[(3)] $(\bar I_k)_{a_k} = (\bar I)_{a_k}$.
\end{enumerate}
Ainsi, quitte à remplacer $a_k$ par une puissance suffisamment grande, nous pouvons écrire
\begin{enumerate}
\item[(4)]
$a_k f_\ell = \sum\nolimits_{j \in E_k} h_{k, \ell}^jf_j + \pi^2 g_{k, \ell}$
\end{enumerate}
pour tout $\ell \in \{1, \ldots, m\}$ et certains
$h_{k, \ell}^j, g_{k, \ell} \in R[x_1, \ldots, x_n]$.
Si $\ell \in E_k$, nous choisissons $h_{k, \ell}^j = a_k\delta_{\ell, j}$
(symbole de Kronecker) et $g_{k, \ell} = 0$. Posons
$$
D = R[x_1, \ldots, x_n, z_1, \ldots, z_m]/
(f_j - \pi z_j, p_{k, \ell}).
$$
Ici, $j \in \{1, \ldots, m\}$, $k \in \{1, \ldots, r\}$,
$\ell \in \{1, \ldots, m\}$, et
$$
p_{k, \ell} = a_k z_\ell - \sum\nolimits_{j \in E_k} h_{k, \ell}^j z_j
- \pi g_{k, \ell}.
$$
Remarquons que, pour $\ell \in E_k$, nos choix ci-dessus donnent $p_{k, \ell} = 0$.

\medskip\noindent
Le morphisme $R \to D$ est le morphisme structural.
Supposons que $\bar C \to \Lambda/\pi^2\Lambda$ envoie $x_i$
sur la classe de $\lambda_i$ modulo $\pi^2$. Pour un élément
$f \in R[x_1, \ldots, x_n]$, nous notons $f(\lambda) \in \Lambda$
le résultat de la substitution de $\lambda_i$ à $x_i$. Nous savons alors que
$f_j(\lambda) = \pi^2 \mu_j$ pour un certain $\mu_j \in \Lambda$.
Définissons $D \to \Lambda$ par les règles $x_i \mapsto \lambda_i$ et
$z_j \mapsto \pi\mu_j$. Ce morphisme est bien défini car
\begin{align*}
p_{k, \ell} & \mapsto
a_k(\lambda) \pi \mu_\ell -
\sum\nolimits_{j \in E_k} h_{k, \ell}^j(\lambda) \pi \mu_j
- \pi g_{k, \ell}(\lambda) \\
& =
\pi\left(a_k(\lambda) \mu_\ell -
\sum\nolimits_{j \in E_k} h_{k, \ell}^j(\lambda) \mu_j
- g_{k, \ell}(\lambda)\right)
\end{align*}
En substituant $x_i = \lambda_i$ dans (4), nous voyons que l'expression
entre parenthèses est annulée par $\pi^2$, donc par $\pi$, puisque nous avons supposé
$\text{Ann}_\Lambda(\pi) = \text{Ann}_\Lambda(\pi^2)$.
Le morphisme $\bar C \to D/\pi D$ est déterminé par $x_i \mapsto x_i$
(il est manifestement bien défini). Il reste donc à démontrer (b), (c) et (d).

\medskip\noindent
À l'aide de (4), nous obtenons l'égalité clé suivante :
\begin{align*}
\pi p_{k, \ell} & =
\pi a_k z_\ell - \sum\nolimits_{j \in E_k} \pi h_{k, \ell}^jz_j
- \pi^2 g_{k, \ell} \\
& =
- a_k (f_\ell - \pi z_\ell) + a_k f_\ell +
\sum\nolimits_{j \in E_k} h_{k, \ell}^j (f_j - \pi z_j) -
\sum\nolimits_{j \in E_k} h_{k, \ell}^j f_j - \pi^2 g_{k, \ell} \\
& =
-a_k(f_\ell - \pi z_\ell) +
\sum\nolimits_{j \in E_k} h_{k, \ell}^j(f_j - \pi z_j)
\end{align*}
Le résultat est un élément de l'idéal engendré par les $f_j - \pi z_j$.
En particulier, nous voyons que $D[1/\pi]$ est isomorphe à
$R[1/\pi][x_1, \ldots, x_n, z_1, \ldots, z_m]/(f_j - \pi z_j)$
qui est isomorphe à $R[1/\pi][x_1, \ldots, x_n]$, donc lisse
sur $R$. Cela démontre (b).

\medskip\noindent
Pour $k \in \{1, \ldots, r\}$ fixé, considérons l'anneau
$$
D_k = R[x_1, \ldots, x_n, z_1, \ldots, z_m]/
(f_j - \pi z_j, j \in E_k, p_{k, \ell})
$$
Le nombre d'équations est $m = |E_k| + (m - |E_k|)$, puisque $p_{k, \ell}$
est nul si $\ell \in E_k$. Remarquons aussi que
\begin{align*}
(D_k/\pi D_k)_{a_k}
& =
R/\pi R[x_1, \ldots, x_n, 1/a_k, z_1, \ldots, z_m]/
(f_j, j \in E_k, p_{k, \ell}) \\
& =
(\bar C/\pi \bar C)_{a_k}[z_1, \ldots, z_m]/
(a_kz_\ell - \sum\nolimits_{j \in E_k} h_{k, \ell}^j z_j) \\
& \cong
(\bar C/\pi \bar C)_{a_k}[z_j, j \in E_k]
\end{align*}
En particulier, $(D_k/\pi D_k)_{a_k}$ est lisse sur $(\bar C/\pi \bar C)_{a_k}$.
Par notre choix de $a_k$, l'anneau $(\bar C/\pi \bar C)_{a_k}$ est lisse
sur $R/\pi R$ de dimension relative $n - |E_k|$, voir (2). Par conséquent,
pour un idéal premier $\mathfrak q_k \subset D_k$ contenant $\pi$ et situé au-dessus de
$\Spec(\bar C_{a_k})$, l'anneau de la fibre de $R \to D_k$
est lisse en $\mathfrak q_k$ de dimension $n$. Ainsi, $R \to D_k$ est syntomique
en $\mathfrak q_k$ d'après le décompte ci-dessus du nombre d'équations, voir
Algèbre, lemme \ref{algebra-lemma-localize-relative-complete-intersection}.
Il s'ensuit que $R \to D_k$ est lisse en $\mathfrak q_k$, voir
Algèbre, lemme \ref{algebra-lemma-flat-fibre-smooth}.

\medskip\noindent
Pour achever la démonstration, soit $\mathfrak q \subset D$ un idéal premier
contenant $\pi$ et situé au-dessus d'un idéal premier où $R/\pi^2 R \to \bar C$
est lisse. D'après (1), $a_k \not \in \mathfrak q$ pour un certain $k$.
Nous allons montrer que la surjection $D_k \to D$ induit
un isomorphisme sur les anneaux locaux en $\mathfrak q$. Nous savons, par le
paragraphe précédent, que les morphismes d'anneaux
$\bar C/\pi \bar C \to D_k/\pi D_k$ et $R \to D_k$ sont lisses en
l'idéal premier correspondant $\mathfrak q_k$ ; cela démontrera (c) et (d),
et achèvera donc la démonstration.

\medskip\noindent
Tout d'abord, pour tout $\ell$, l'équation démontrée ci-dessus
$\pi p_{k, \ell} = -a_k(f_\ell - \pi z_\ell) +
\sum_{j \in E_k} h_{k, \ell}^j (f_j - \pi z_j)$ montre que
$f_\ell - \pi z_\ell$ a une image nulle dans $(D_k)_{a_k}$, et donc dans
$(D_k)_{\mathfrak q_k}$.
Les relations (4) impliquent que $a_k f_\ell =
\sum_{j \in E_k} h_{k, \ell}^j f_j$ dans $I/I^2$.
Comme $(\bar I_k/\bar I_k^2)_{a_k}$ est libre sur les $f_j$, $j \in E_k$,
nous voyons que
$$
a_{k'} h_{k, \ell}^j -
\sum\nolimits_{j' \in E_{k'}} h_{k', \ell}^{j'} h_{k, j'}^j
$$
est nul dans $\bar C_{a_k}$ pour tous $k, k', \ell$ et $j \in E_k$.
Il existe donc un entier $N$ suffisamment grand tel que
$$
a_k^N\left(
a_{k'} h_{k, \ell}^j -
\sum\nolimits_{j' \in E_{k'}} h_{k', \ell}^{j'} h_{k, j'}^j
\right)
$$
appartienne à $I_k + \pi^2R[x_1, \ldots, x_n]$. En calculant modulo $\pi$, nous obtenons
\begin{align*}
&
a_kp_{k', \ell} - a_{k'}p_{k, \ell} + \sum h_{k', \ell}^{j'} p_{k, j'}
\\
&
=
- a_k \sum h_{k', \ell}^{j'} z_{j'}
+ a_{k'} \sum h_{k, \ell}^j z_j
+ \sum h_{k', \ell}^{j'} a_k z_{j'}
- \sum \sum h_{k', \ell}^{j'} h_{k, j'}^j z_j \\
&
=
\sum \left(
a_{k'} h_{k, \ell}^j
- \sum h_{k', \ell}^{j'} h_{k, j'}^j
\right) z_j
\end{align*}
avec la convention de sommation d'Einstein. En combinant ce calcul avec ce qui précède,
nous voyons que $a_k^{N + 1} p_{k', \ell}$ appartient à l'idéal engendré
par $I_k$ et $\pi$ dans $R[x_1, \ldots, x_n, z_1, \ldots, z_m]$.
Ainsi, $p_{k', \ell}$ s'envoie dans $\pi (D_k)_{a_k}$. D'autre part, l'équation
$$
\pi p_{k', \ell} =
-a_{k'} (f_\ell - \pi z_\ell) +
\sum\nolimits_{j' \in E_{k'}} h_{k', \ell}^{j'}(f_{j'} - \pi z_{j'})
$$
montre que $\pi p_{k', \ell}$ est nul dans $(D_k)_{a_k}$.
Comme nous avons supposé que $\text{Ann}_R(\pi) = \text{Ann}_R(\pi^2)$
et que $(D_k)_{\mathfrak q_k}$ est lisse, donc plat, sur $R$,
nous voyons que
$\text{Ann}_{(D_k)_{\mathfrak q_k}}(\pi) =
\text{Ann}_{(D_k)_{\mathfrak q_k}}(\pi^2)$.
Nous en concluons que $p_{k', \ell}$ s'annule également, et donc que
$D_{\mathfrak q} = (D_k)_{\mathfrak q_k}$, ce qui conclut la démonstration.
\end{proof}
\section{Le lemme de désingularisation}
\label{section-desingularization-lemma}

\noindent
Voici un autre lemme d'une ingéniosité diabolique.

\begin{lemma}
\label{lemma-desingularize}
Soit $R$ un anneau noethérien.
Soit $\Lambda$ une $R$-algèbre. Soit $\pi \in R$ et
supposons que $\text{Ann}_\Lambda(\pi) = \text{Ann}_\Lambda(\pi^2)$. Soit
$A \to \Lambda$ un morphisme de $R$-algèbres, avec $A$ de
présentation finie. Supposons que
\begin{enumerate}
\item l'image de $\pi$ est strictement standard dans $A$ sur $R$, et
\item il existe une section $\rho : A/\pi^4 A \to R/\pi^4 R$
compatible avec le morphisme vers $\Lambda/\pi^4 \Lambda$.
\end{enumerate}
Alors il existe des morphismes de $R$-algèbres $A \to B \to \Lambda$, avec
$B$ de présentation finie, tels que $\mathfrak a B \subset H_{B/R}$, où
$\mathfrak a = \text{Ann}_R(\text{Ann}_R(\pi^2)/\text{Ann}_R(\pi))$.
\end{lemma}

\begin{proof}
Choisissons une présentation
$$
A = R[x_1, \ldots, x_n]/(f_1, \ldots, f_m)
$$
et un entier $0 \leq c \leq \min(n, m)$ tels que
(\ref{equation-strictly-standard-one}) soit satisfaite pour $\pi$ et que
\begin{equation}
\label{equation-star}
\pi f_{c + j} \in (f_1, \ldots, f_c) + (f_1, \ldots, f_m)^2
\end{equation}
pour $j = 1, \ldots, m - c$. Disons que $\rho$ envoie $x_i$ sur la classe de
$r_i \in R$. Nous pouvons alors remplacer $x_i$ par $x_i - r_i$. Ainsi, nous pouvons
supposer que $\rho(x_i) = 0$ dans $R/\pi^4 R$. Il s'ensuit que
$f_j(0) \in \pi^4R$ et que $A \to \Lambda$ envoie $x_i$
sur $\pi^4\lambda_i$ pour un certain $\lambda_i \in \Lambda$. Écrivons
$$
f_j = f_j(0) + \sum\nolimits_{i = 1, \ldots, n} r_{ji} x_i + \text{termes d'ordre supérieur}
$$
Il en résulte que le terme constant de $\partial f_j/\partial x_i$ est
$r_{ji}$. Appliquons $\rho$ à (\ref{equation-strictly-standard-one})
pour $\pi$ ; nous obtenons
$$
\pi = \sum\nolimits_{I \subset \{1, \ldots, n\},\ |I| = c}
r_I \det(r_{ji})_{j = 1, \ldots, c,\ i \in I} \bmod \pi^4R
$$
pour certains $r_I \in R$. Nous avons donc
$$
u\pi = \sum\nolimits_{I \subset \{1, \ldots, n\},\ |I| = c}
r_I \det(r_{ji})_{j = 1, \ldots, c,\ i \in I}
$$
pour un certain $u \in 1 + \pi^3R$. D'après
Algèbre, lemme \ref{algebra-lemma-matrix-left-inverse},
il existe alors une matrice $n \times c$ $(s_{ik})$ telle que
$$
u\pi \delta_{jk} = \sum\nolimits_{i = 1, \ldots, n} r_{ji}s_{ik}\quad
\text{pour tous } j, k = 1, \ldots, c
$$
(symbole de Kronecker). Introduisons des variables auxiliaires
$v_1, \ldots, v_c, w_1, \ldots, w_n$ et posons
$$
h_i = x_i - \pi^2 \sum\nolimits_{j = 1, \ldots c} s_{ij} v_j - \pi^3 w_i
$$
Dans la suite, nous utiliserons que
$$
R[x_1, \ldots, x_n, v_1, \ldots, v_c, w_1, \ldots, w_n]/
(h_1, \ldots, h_n) = R[v_1, \ldots, v_c, w_1, \ldots, w_n]
$$
sans autre mention. Dans
$R[x_1, \ldots, x_n, v_1, \ldots, v_c, w_1, \ldots, w_n]/
(h_1, \ldots, h_n)$, nous avons
\begin{align*}
f_j & = f_j(x_1 - h_1, \ldots, x_n - h_n) \\
& =
\pi^2 \sum\nolimits_{k = 1}^c
\left(\sum\nolimits_{i = 1}^n r_{ji} s_{ik}\right) v_k
+
\pi^3 \sum\nolimits_{i = 1}^n r_{ji}w_i \bmod \pi^4 \\
& =
\pi^3 v_j + \pi^3 \sum\nolimits_{i = 1}^n r_{ji}w_i \bmod \pi^4
\end{align*}
pour $1 \leq j \leq c$. Nous pouvons donc choisir des éléments
$g_j \in R[v_1, \ldots, v_c, w_1, \ldots, w_n]$
tels que $g_j = v_j + \sum r_{ji}w_i \bmod \pi$
et que $f_j = \pi^3 g_j$ dans la $R$-algèbre :
$$
R[x_1, \ldots, x_n, v_1, \ldots, v_c, w_1, \ldots, w_n]/
(h_1, \ldots, h_n)
$$
Posons
$$
B = R[x_1, \ldots, x_n, v_1, \ldots, v_c, w_1, \ldots, w_n]/
(f_1, \ldots, f_m, h_1, \ldots, h_n, g_1, \ldots, g_c).
$$
Le morphisme $A \to B$ est évident. Définissons $B \to \Lambda$ en envoyant
$x_i \to \pi^4\lambda_i$, $v_i \mapsto 0$, et $w_i \mapsto \pi \lambda_i$.
Il est alors clair que les éléments $f_j$ et $h_i$ ont pour image zéro
dans $\Lambda$. De plus, $g_i$ a manifestement pour image un élément
$t$ de $\pi\Lambda$ tel que $\pi^3t = 0$ (car $f_i = \pi^3 g_i$ modulo
l'idéal engendré par les $h$). Notre hypothèse
$\text{Ann}_\Lambda(\pi) = \text{Ann}_\Lambda(\pi^2)$ entraîne donc que $t = 0$.
Il reste ainsi à démontrer l'assertion concernant la lissité.

\medskip\noindent
Remarquons que $B_\pi \cong A_\pi[v_1, \ldots, v_c]$, car les équations
$g_i = 0$ résultent de $f_i = 0$. Ainsi, $B_\pi$ est lisse sur $R$
puisque $A_\pi$ est lisse sur $R$ par l'hypothèse selon laquelle $\pi$ est
strictement standard dans $A$ sur $R$ ; voir le
Lemme \ref{lemma-elkik}.

\medskip\noindent
Posons $B' = R[v_1, \ldots, v_c, w_1, \ldots, w_n]/(g_1, \ldots, g_c)$.
Comme $g_i = v_i + \sum r_{ji}w_i \bmod \pi$, nous voyons que
$B'/\pi B' = R/\pi R[w_1, \ldots, w_n]$. Ainsi,
$R \to B'$ est lisse de dimension relative $n$ en tout
point de $V(\pi)$ d'après
Algèbre, lemmes
\ref{algebra-lemma-localize-relative-complete-intersection} et
\ref{algebra-lemma-flat-fibre-smooth}
(le premier lemme montre qu'il est syntomique en ces idéaux premiers, donc en particulier
plat, après quoi le second lemme montre qu'il est lisse).

\medskip\noindent
Soit $\mathfrak q \subset B$ un idéal premier tel que $\pi \in \mathfrak q$ et
tel que, pour un certain $r \in \mathfrak a$, $r \not \in \mathfrak q$.
Notons $\mathfrak q' = B' \cap \mathfrak q$.
Nous affirmons que la surjection $B' \to B$ induit un isomorphisme d'anneaux
locaux $(B')_{\mathfrak q'} \to B_\mathfrak q$. Cela
achèvera la démonstration du lemme. Remarquons que $B_\mathfrak q$ est le
quotient de $(B')_{\mathfrak q'}$ par l'idéal engendré par
$f_{c + j}$, $j = 1, \ldots, m - c$. Nous observons deux faits :
premièrement, l'image de $f_{c + j}$ dans $(B')_{\mathfrak q'}$ est
divisible par $\pi^2$ ;
deuxièmement, l'image de $\pi f_{c + j}$ dans $(B')_{\mathfrak q'}$
peut s'écrire $\sum b_{j_1 j_2} f_{c + j_1}f_{c + j_2}$ d'après
(\ref{equation-star}). Nous voyons donc que l'image de chaque $\pi f_{c + j}$
appartient à l'idéal engendré par les éléments $\pi^2 f_{c + j'}$.
Ainsi, $\pi f_{c + j} = 0$ dans $(B')_{\mathfrak q'}$, puisque celui-ci est un
anneau local noethérien ; voir
Algèbre, lemme \ref{algebra-lemma-intersect-powers-ideal-module-zero}.
Comme $R \to (B')_{\mathfrak q'}$ est plat, nous avons
$$
\left(\text{Ann}_R(\pi^2)/\text{Ann}_R(\pi)\right)
\otimes_R (B')_{\mathfrak q'}
=
\text{Ann}_{(B')_{\mathfrak q'}}(\pi^2)/\text{Ann}_{(B')_{\mathfrak q'}}(\pi)
$$
Puisque $r \in \mathfrak a$ est inversible dans
$(B')_{\mathfrak q'}$, nous voyons que ce module est nul.
Il s'ensuit que l'image de $f_{c + j}$ est nulle dans
$(B')_{\mathfrak q'}$, comme souhaité.
\end{proof}

\begin{lemma}
\label{lemma-desingularize-strictly-standard}
Soit $R$ un anneau noethérien. Soit $\Lambda$ une $R$-algèbre.
Soit $\pi \in R$ et supposons que $\text{Ann}_R(\pi) = \text{Ann}_R(\pi^2)$ et
$\text{Ann}_\Lambda(\pi) = \text{Ann}_\Lambda(\pi^2)$.
Soient $A \to \Lambda$ et $D \to \Lambda$ des morphismes de $R$-algèbres, avec
$A$ et $D$ de présentation finie. Supposons que
\begin{enumerate}
\item $\pi$ est strictement standard dans $A$ sur $R$, et
\item il existe un morphisme de $R$-algèbres $A/\pi^4 A \to D/\pi^4 D$ compatible
avec les morphismes vers $\Lambda/\pi^4 \Lambda$.
\end{enumerate}
Alors il existe un morphisme de $R$-algèbres $B \to \Lambda$, avec $B$ de
présentation finie, ainsi que des morphismes de $R$-algèbres $A \to B$ et $D \to B$
compatibles avec les morphismes vers $\Lambda$, tels que $H_{D/R}B \subset H_{B/D}$
et $H_{D/R}B \subset H_{B/R}$.
\end{lemma}

\begin{proof}
Appliquons le Lemme \ref{lemma-desingularize} à
$$
D \longrightarrow A \otimes_R D \longrightarrow \Lambda
$$
et à l'image de $\pi$ dans $D$. D'après le
Lemme \ref{lemma-strictly-standard-base-change},
$\pi$ est strictement standard dans $A \otimes_R D$ sur $D$.
Comme section $\rho : (A \otimes_R D)/\pi^4 (A \otimes_R D) \to D/\pi^4 D$,
prenons le morphisme induit par celui de (2). Ainsi, le
Lemme \ref{lemma-desingularize} s'applique et fournit une factorisation
$A \otimes_R D \to B \to \Lambda$, avec $B$ de présentation finie,
et $\mathfrak a B \subset H_{B/D}$, où
$$
\mathfrak a = \text{Ann}_D(\text{Ann}_D(\pi^2)/\text{Ann}_D(\pi)).
$$
Pour tout idéal premier $\mathfrak q$ de $D$ tel que $D_\mathfrak q$ soit plat sur $R$,
nous avons
$\text{Ann}_{D_\mathfrak q}(\pi^2)/\text{Ann}_{D_\mathfrak q}(\pi) = 0$
car la formation des annulateurs d'éléments commute au changement de base plat et
nous avons supposé $\text{Ann}_R(\pi) = \text{Ann}_R(\pi^2)$. Comme $D$ est
noethérien, $\text{Ann}_D(\pi^2)/\text{Ann}_D(\pi)$ est un
$D$-module de type fini ; la formation de son annulateur commute donc à la localisation.
Nous voyons ainsi que $\mathfrak a \not \subset \mathfrak q$. Par conséquent,
$D \to B$ est lisse en tout idéal premier de $B$ situé au-dessus de $\mathfrak q$.
Puisque, en tout idéal premier de $D$ où $R \to D$ est lisse,
$D_\mathfrak q$ est plat sur $R$, nous concluons que $H_{D/R}B \subset H_{B/D}$.
La dernière inclusion $H_{D/R}B \subset H_{B/R}$ résulte du fait que les composées
de morphismes d'anneaux lisses sont lisses
(Algèbre, lemme \ref{algebra-lemma-compose-smooth}).
\end{proof}

\begin{lemma}
\label{lemma-desingularize-lifting-apply}
Soit $R$ un anneau noethérien. Soit $\Lambda$ une $R$-algèbre.
Soit $\pi \in R$ et supposons que $\text{Ann}_R(\pi) = \text{Ann}_R(\pi^2)$ et
$\text{Ann}_\Lambda(\pi) = \text{Ann}_\Lambda(\pi^2)$.
Soit $A \to \Lambda$ un morphisme de $R$-algèbres, avec
$A$ de présentation finie, et supposons que $\pi$ est strictement standard
dans $A$ sur $R$. Soit
$$
A/\pi^8A \to \bar C \to \Lambda/\pi^8\Lambda
$$
une factorisation, avec $\bar C$ de présentation finie.
Alors il existe une factorisation $A \to B \to \Lambda$, avec $B$ de
présentation finie, telle que $R_\pi \to B_\pi$ soit lisse et que
$$
H_{\bar C/(R/\pi^8 R)} \cdot \Lambda/\pi^8\Lambda
\subset
\sqrt{H_{B/R} \Lambda} \bmod \pi^8\Lambda.
$$
\end{lemma}

\begin{proof}
Appliquons le Lemme \ref{lemma-lifting} pour obtenir $R \to D \to \Lambda$
avec une factorisation
$\bar C/\pi^4\bar C \to D/\pi^4 D \to \Lambda/\pi^4\Lambda$
telle que $R \to D$ soit lisse en tout idéal premier ne contenant pas $\pi$
et en tout idéal premier situé au-dessus d'un idéal premier de $\bar C/\pi^4\bar C$
où $R/\pi^8 R \to \bar C$ est lisse.
D'après le Lemme \ref{lemma-desingularize-strictly-standard},
il existe une $R$-algèbre $B$ de présentation finie et des
factorisations $A \to B \to \Lambda$ et $D \to B \to \Lambda$
telles que $H_{D/R}B \subset H_{B/R}$. Nous omettons de vérifier qu'il
s'agit bien d'une solution au problème posé par le lemme.
\end{proof}












\section{Préliminaire : réduction à un corps de base}
\label{section-reduction}

\noindent
Dans cette section, nous appliquons les lemmes des sections précédentes
pour montrer qu'il suffit de démontrer le résultat principal lorsque l'anneau
de base est un corps ; voir le Lemme \ref{lemma-reduce-to-field}.

\begin{situation}
\label{situation-global}
Ici, $R \to \Lambda$ est un morphisme d'anneaux régulier entre anneaux noethériens.
\end{situation}

\noindent
Soit $R \to \Lambda$ comme dans la Situation \ref{situation-global}.
Nous disons que {\it PT est vraie pour $R \to \Lambda$} si $\Lambda$ est une
colimite filtrante de $R$-algèbres lisses.

\begin{lemma}
\label{lemma-product}
Soient $R_i \to \Lambda_i$, $i = 1, 2$, comme dans la Situation \ref{situation-global}.
Si PT est vraie pour $R_i \to \Lambda_i$, $i = 1, 2$, alors PT est vraie pour
$R_1 \times R_2 \to \Lambda_1 \times \Lambda_2$.
\end{lemma}

\begin{proof}
Omis. Indication : un produit de colimites filtrantes est une colimite filtrante.
\end{proof}

\begin{lemma}
\label{lemma-delocalize-base}
Soient $R \to A \to \Lambda$ des morphismes d'anneaux, avec $A$ de présentation finie
sur $R$. Soit $S \subset R$ une partie
multiplicative. Soit $S^{-1}A \to B' \to S^{-1}\Lambda$ une factorisation telle que
$B'$ soit lisse sur $S^{-1}R$. Alors il existe une factorisation
$A \to B \to \Lambda$ telle que l'image d'un certain $s \in S$ soit un
élément standard élémentaire (Définition \ref{definition-strictly-standard})
dans $B$ sur $R$.
\end{lemma}

\begin{proof}
Appliquons d'abord le Lemme \ref{lemma-smooth-standard-smooth} à $S^{-1}R \to B'$.
Nous pouvons donc supposer que $B'$ est standard lisse sur $S^{-1}R$.
Écrivons $A = R[x_1, \ldots, x_n]/(g_1, \ldots, g_t)$ et supposons que
$x_i \mapsto \lambda_i$ dans $\Lambda$. Nous pouvons écrire
$B' = S^{-1}R[x_1, \ldots, x_{n + m}]/(f_1, \ldots, f_c)$
pour un certain $c \geq n$, où
$\det(\partial f_j/\partial x_i)_{i, j = 1, \ldots, c}$
est inversible dans $B'$ et tel que $A \to B'$ soit donné par $x_i \mapsto x_i$ ;
voir le Lemme \ref{lemma-standard-smooth-include-generators}.
Après avoir multiplié $x_i$, $i > n$, par un élément de $S$ et modifié en conséquence
les équations $f_j$, nous pouvons supposer que $B' \to S^{-1}\Lambda$ envoie
$x_i$ sur $\lambda_i/1$ pour un certain $\lambda_i \in \Lambda$, lorsque $i > n$.
Choisissons une relation
$$
1 =
a_0 \det(\partial f_j/\partial x_i)_{i, j = 1, \ldots, c}
+
\sum\nolimits_{j = 1, \ldots, c} a_jf_j
$$
pour certains $a_j \in S^{-1}R[x_1, \ldots, x_{n + m}]$. Puisque tout élément de $S$
est inversible dans $B'$, nous pouvons, en chassant les dénominateurs, supposer que
$f_j, a_j \in R[x_1, \ldots, x_{n + m}]$ et que
$$
s_0 = a_0 \det(\partial f_j/\partial x_i)_{i, j = 1, \ldots, c}
+
\sum\nolimits_{j = 1, \ldots, c} a_jf_j
$$
pour un certain $s_0 \in S$. Puisque $g_j$ a pour image zéro dans
$S^{-1}R[x_1, \ldots, x_{n + m}]/(f_1, \ldots, f_c)$
il existe des éléments $s_j \in S$ tels que $s_j g_j = 0$ dans
$R[x_1, \ldots, x_{n + m}]/(f_1, \ldots, f_c)$.
Puisque $f_j$ a pour image zéro dans $S^{-1}\Lambda$, il existe $s'_j \in S$
tel que $s'_j f_j(\lambda_1, \ldots, \lambda_{n + m}) = 0$ dans
$\Lambda$. Considérons l'anneau
$$
B = R[x_1, \ldots, x_{n + m}]/
(s'_1f_1, \ldots, s'_cf_c, g_1, \ldots, g_t)
$$
et la factorisation $A \to B \to \Lambda$, où $B \to \Lambda$ est donné par
$x_i \mapsto \lambda_i$. Nous affirmons que $s = s_0s_1 \ldots s_ts'_1 \ldots s'_c$
est standard élémentaire dans $B$ sur $R$, ce qui achève la démonstration.
En effet, $s_j g_j \in (f_1, \ldots, f_c)$, et donc
$sg_j \in (s'_1f_1, \ldots, s'_cf_c)$. Enfin, nous avons
$$
a_0\det(\partial s'_jf_j/\partial x_i)_{i, j = 1, \ldots, c}
+
\sum\nolimits_{j = 1, \ldots, c}
(s'_1 \ldots \hat{s'_j} \ldots s'_c) a_j s'_jf_j
=
s_0s'_1\ldots s'_c
$$
qui divise $s$, comme souhaité.
\end{proof}

\begin{lemma}
\label{lemma-reduce-to-field}
\begin{slogan}
Démontrer l'approximation de Popescu se ramène aux algèbres sur un corps
\end{slogan}
Si PT est vraie pour toute Situation \ref{situation-global} où $R$
est un corps, alors PT est vraie en général.
\end{lemma}

\begin{proof}
Supposons que PT soit vraie pour toute Situation \ref{situation-global} où $R$ est un corps.
Soit $R \to \Lambda$ une Situation \ref{situation-global} arbitraire.
Remarquons que $R/I \to \Lambda/I\Lambda$ est encore un morphisme d'anneaux régulier
entre anneaux noethériens ; voir
Compléments d'algèbre, lemme \ref{more-algebra-lemma-regular-base-change}.
Considérons l'ensemble d'idéaux
$$
\mathcal{I} = \{I \subset R \mid R/I \to \Lambda/I\Lambda
\text{ ne vérifie pas PT}\}
$$
Nous devons montrer que $\mathcal{I}$ est vide. Si cet ensemble est non vide,
il contient un élément maximal puisque $R$ est noethérien.
En remplaçant $R$ par $R/I$ et $\Lambda$ par $\Lambda/I$, nous obtenons une
situation où PT est vraie pour $R/I \to \Lambda/I\Lambda$ pour tout
idéal non nul de $R$. En particulier, en appliquant la
Proposition \ref{proposition-lift},
nous voyons que $R$ est un anneau réduit.

\medskip\noindent
Soit $A \to \Lambda$ un morphisme de $R$-algèbres, avec $A$ de
présentation finie. Nous devons trouver une factorisation $A \to B \to \Lambda$
où $B$ est lisse sur $R$ ; voir Algèbre, lemme \ref{algebra-lemma-when-colimit}.

\medskip\noindent
Soit $S \subset R$ l'ensemble des non-diviseurs de zéro et
considérons l'anneau total des fractions $Q = S^{-1}R$ de $R$. Nous savons que
$Q = K_1 \times \ldots \times K_n$ est un produit de corps ; voir
Algèbre, lemmes \ref{algebra-lemma-total-ring-fractions-no-embedded-points} et
\ref{algebra-lemma-Noetherian-irreducible-components}.
D'après le Lemme \ref{lemma-product} et notre hypothèse,
PT est vraie pour le morphisme d'anneaux $S^{-1}R \to S^{-1}\Lambda$.
Il existe donc une factorisation $S^{-1}A \to B' \to S^{-1}\Lambda$
où $B'$ est lisse sur $S^{-1}R$.

\medskip\noindent
Appliquons le Lemme \ref{lemma-delocalize-base} ;
nous obtenons une factorisation $A \to B \to \Lambda$ telle qu'un
certain $\pi \in S$ soit standard élémentaire dans $B$ sur $R$.
Après avoir remplacé $A$ par $B$, nous pouvons supposer que $\pi$ est
standard élémentaire, donc strictement standard dans $A$. Nous savons que
$R/\pi^8R \to \Lambda/\pi^8\Lambda$ vérifie PT.
Il existe donc une factorisation
$R/\pi^8 R \to A/\pi^8A \to \bar C \to \Lambda/\pi^8\Lambda$
où $R/\pi^8 R \to \bar C$ est lisse. D'après le
Lemme \ref{lemma-lifting},
il existe un morphisme de $R$-algèbres $D \to \Lambda$, avec $D$ lisse sur $R$,
et une factorisation
$R/\pi^4 R \to A/\pi^4A \to D/\pi^4D \to \Lambda/\pi^4\Lambda$.
D'après le Lemme \ref{lemma-desingularize-strictly-standard},
il existe $A \to B \to \Lambda$ avec $B$ lisse sur $R$,
ce qui achève la démonstration.
\end{proof}










\section{Arguments locaux}
\label{section-local}


\begin{situation}
\label{situation-local}
On se donne un anneau noethérien $R$, un morphisme de $R$-algèbres $A \to \Lambda$
et un idéal premier $\mathfrak q \subset \Lambda$. Nous supposons que $A$ est de
présentation finie sur $R$. Dans cette situation, nous notons
$\mathfrak h_A = \sqrt{H_{A/R} \Lambda}$.
\end{situation}

\noindent
Soit $R \to A \to \Lambda \supset \mathfrak q$ comme dans la
Situation \ref{situation-local}. Nous disons que
{\it $R \to A \to \Lambda \supset \mathfrak q$
peut être résolu} s'il existe une factorisation $A \to B \to \Lambda$
avec $B$ de présentation finie et
$\mathfrak h_A \subset \mathfrak h_B \not \subset \mathfrak q$.
Dans ce cas, nous appelons la factorisation $A \to B \to \Lambda$
une {\it résolution de $R \to A \to \Lambda \supset \mathfrak q$}.

\begin{lemma}
\label{lemma-lift-solution}
Soit $R \to A \to \Lambda \supset \mathfrak q$ comme dans la
Situation \ref{situation-local}. Soit $r \geq 1$, et supposons que
$\pi_1, \ldots, \pi_r \in R$ ont leurs images dans $\mathfrak q$. Supposons que
\begin{enumerate}
\item pour $i = 1, \ldots, r$, nous avons
$$
\text{Ann}_{R/(\pi_1^8, \ldots, \pi_{i - 1}^8)R}(\pi_i)
=
\text{Ann}_{R/(\pi_1^8, \ldots, \pi_{i - 1}^8)R}(\pi_i^2)
$$
et
$$
\text{Ann}_{\Lambda/(\pi_1^8, \ldots, \pi_{i - 1}^8)\Lambda}(\pi_i)
=
\text{Ann}_{\Lambda/(\pi_1^8, \ldots, \pi_{i - 1}^8)\Lambda}(\pi_i^2)
$$
\item pour $i = 1, \ldots, r$, l'élément $\pi_i$ a pour image un élément
strictement standard dans $A$ sur $R$.
\end{enumerate}
Alors, si
$$
R/(\pi_1^8, \ldots, \pi_r^8)R \to A/(\pi_1^8, \ldots, \pi_r^8)A
\to \Lambda/(\pi_1^8, \ldots, \pi_r^8)\Lambda \supset
\mathfrak q/(\pi_1^8, \ldots, \pi_r^8)\Lambda
$$
peut être résolu, il en va de même de $R \to A \to \Lambda \supset \mathfrak q$.
\end{lemma}

\begin{proof}
Nous allons le démontrer par récurrence sur $r$.

\medskip\noindent
Le cas $r = 1$. L'hypothèse affirme ici qu'il existe une
factorisation $A/\pi_1^8 \to \bar C \to \Lambda/\pi_1^8$
qui résout la situation modulo $\pi_1^8$. Les conditions (1) et
(2) sont les hypothèses nécessaires pour appliquer le
Lemme \ref{lemma-desingularize-lifting-apply}.
Nous pouvons donc ``relever'' la résolution $\bar C$
en une résolution de $R \to A \to \Lambda \supset \mathfrak q$.

\medskip\noindent
Le cas $r > 1$. Dans ce cas, appliquons l'hypothèse de récurrence pour $r - 1$
à la situation
$R/\pi_1^8 \to A/\pi_1^8 \to \Lambda/\pi_1^8
\supset \mathfrak q/\pi_1^8\Lambda$.
Remarquons que la propriété (2) est préservée par le
Lemme \ref{lemma-strictly-standard-base-change}.
\end{proof}

\begin{lemma}
\label{lemma-delocalize-weak}
\begin{reference}
\cite[Lemma 12.2]{swan} ou
\cite[Lemma 2]{popescu-GND}
\end{reference}
Soit $R \to A \to \Lambda \supset \mathfrak q$ comme dans la
Situation \ref{situation-local}. Soit $\mathfrak p = R \cap \mathfrak q$.
Supposons que $\mathfrak q$ soit minimal au-dessus de $\mathfrak h_A$ et que
$R_\mathfrak p \to A_\mathfrak p \to \Lambda_\mathfrak q
\supset \mathfrak q\Lambda_\mathfrak q$ puisse être résolu.
Alors il existe une factorisation $A \to C \to \Lambda$, avec $C$ de
présentation finie, telle que $H_{C/R} \Lambda \not \subset \mathfrak q$.
\end{lemma}

\begin{proof}
Soit $A_\mathfrak p \to C \to \Lambda_\mathfrak q$ une résolution de
$R_\mathfrak p \to A_\mathfrak p \to \Lambda_\mathfrak q
\supset \mathfrak q\Lambda_\mathfrak q$. Comme $\mathfrak q$ est, par hypothèse,
minimal au-dessus de $\mathfrak h_A$, cela
signifie que $H_{C/R_\mathfrak p} \Lambda_\mathfrak q = \Lambda_\mathfrak q$.
D'après le Lemme \ref{lemma-final-solve},
nous pouvons supposer que $C$ est lisse sur $R_\mathfrak p$.
D'après le Lemme \ref{lemma-smooth-standard-smooth}, nous pouvons supposer que
$C$ est standard lisse sur $R_\mathfrak p$.
Écrivons $A = R[x_1, \ldots, x_n]/(g_1, \ldots, g_t)$ et supposons que
$A \to \Lambda$ est donné par $x_i \mapsto \lambda_i$.
Écrivons $C = R_\mathfrak p[x_1, \ldots, x_{n + m}]/(f_1, \ldots, f_c)$
pour un certain $c \geq n$, de sorte que $A \to C$ envoie $x_i$ sur $x_i$ et que
$\det(\partial f_j/\partial x_i)_{i, j = 1, \ldots, c}$
soit inversible dans $C$ ; voir le
Lemme \ref{lemma-standard-smooth-include-generators}.
Après avoir chassé les dénominateurs, nous pouvons supposer que
$f_1, \ldots, f_c$ sont des éléments de $R[x_1, \ldots, x_{n + m}]$.
Bien entendu,
$\det(\partial f_j/\partial x_i)_{i, j = 1, \ldots, c}$
n'est pas inversible dans $R[x_1, \ldots, x_{n + m}]/(f_1, \ldots, f_c)$,
mais il le devient après inversion d'un certain élément $s_0 \in R$,
$s_0 \not \in \mathfrak p$.
Comme $g_j$ a pour image zéro par $R[x_1, \ldots, x_n] \to A \to C$,
il existe $s_j \in R$, $s_j \not \in \mathfrak p$, tel que
$s_j g_j$ soit nul dans $R[x_1, \ldots, x_{n + m}]/(f_1, \ldots, f_c)$.
Écrivons $f_j = F_j(x_1, \ldots, x_{n + m}, 1)$
pour un polynôme
$F_j \in R[x_1, \ldots, x_n, X_{n + 1}, \ldots, X_{n + m + 1}]$
homogène en $X_{n + 1}, \ldots, X_{n + m + 1}$.
Choisissons $\lambda_{n + i} \in \Lambda$, $i = 1, \ldots, m + 1$, avec
$\lambda_{n + m + 1} \not \in \mathfrak q$, de sorte que $x_{n + i}$ ait pour image
$\lambda_{n + i}/\lambda_{n + m + 1}$ dans $\Lambda_\mathfrak q$.
Alors
\begin{align*}
F_j(\lambda_1, \ldots, \lambda_{n + m + 1})
& =
(\lambda_{n + m + 1})^{\deg(F_j)} F_j(\lambda_1, \ldots, \lambda_n,
\frac{\lambda_{n + 1}}{\lambda_{n + m + 1}}, \ldots,
\frac{\lambda_{n + m}}{\lambda_{n + m + 1}}, 1) \\
& =
(\lambda_{n + m + 1})^{\deg(F_j)} f_j(\lambda_1, \ldots, \lambda_n,
\frac{\lambda_{n + 1}}{\lambda_{n + m + 1}}, \ldots,
\frac{\lambda_{n + m}}{\lambda_{n + m + 1}}) \\
& = 0
\end{align*}
dans $\Lambda_\mathfrak q$. Il existe donc
$\lambda_0 \in \Lambda$, $\lambda_0 \not \in \mathfrak q$, tel que
$\lambda_0 F_j(\lambda_1, \ldots, \lambda_{n + m + 1}) = 0$
dans $\Lambda$. Posons maintenant $B$ égal à
$$
R[x_0, \ldots, x_{n + m + 1}]/
(g_1, \ldots, g_t, x_0F_1(x_1, \ldots, x_{n + m + 1}), \ldots,
x_0F_c(x_1, \ldots, x_{n + m + 1}))
$$
que nous envoyons dans $\Lambda$ en envoyant $x_i$ sur $\lambda_i$.
Soit $b$ l'image de $x_0 x_{n + m + 1} s_0 s_1 \ldots s_t$ dans $B$.
Alors $B_b$ est isomorphe à
$$
R_{s_0s_1 \ldots s_t}[x_0, x_1, \ldots, x_{n + m + 1}, 1/x_0x_{n + m + 1}]/
(f_1, \ldots, f_c)
$$
qui est lisse sur $R$ par construction.
Puisque l'image de $b$ n'appartient pas à $\mathfrak q$, le résultat s'ensuit.
\end{proof}

\begin{lemma}
\label{lemma-delocalize-height-zero}
Soit $R \to A \to \Lambda \supset \mathfrak q$ comme dans la
Situation \ref{situation-local}. Soit $\mathfrak p = R \cap \mathfrak q$.
Supposons que
\begin{enumerate}
\item $\mathfrak q$ est minimal au-dessus de $\mathfrak h_A$,
\item $R_\mathfrak p \to A_\mathfrak p \to \Lambda_\mathfrak q
\supset \mathfrak q\Lambda_\mathfrak q$ peut être résolu, et
\item $\dim(\Lambda_\mathfrak q) = 0$.
\end{enumerate}
Alors $R \to A \to \Lambda \supset \mathfrak q$ peut être résolu.
\end{lemma}

\begin{proof}
Par (3), l'anneau $\Lambda_\mathfrak q$ est local artinien ; ainsi,
$\mathfrak q\Lambda_\mathfrak q$ est nilpotent. Par conséquent,
$(\mathfrak h_A)^N \Lambda_\mathfrak q = 0$ pour un certain $N > 0$.
Il existe donc $\lambda \in \Lambda$, $\lambda \not \in \mathfrak q$,
tel que $\lambda (\mathfrak h_A)^N = 0$ dans $\Lambda$.
Écrivons $H_{A/R} = (a_1, \ldots, a_r)$, de sorte que $\lambda a_i^N = 0$
dans $\Lambda$. D'après le Lemme \ref{lemma-delocalize-weak}, il existe une factorisation
$A \to C \to \Lambda$, avec $C$ de présentation finie, telle que
$\mathfrak h_C \not \subset \mathfrak q$.
Écrivons $C = A[x_1, \ldots, x_n]/(f_1, \ldots, f_m)$.
Posons
$$
B = A[x_1, \ldots, x_n, y_1, \ldots, y_r, z, t_{ij}]/
(f_j - \sum y_i t_{ij}, zy_i)
$$
où les $t_{ij}$ forment un ensemble de $rm$ variables.
Remarquons qu'il existe un morphisme $B \to C[y_i, z]/(y_iz)$ obtenu en posant les $t_{ij}$
égaux à zéro. Le morphisme $B \to \Lambda$ est la composée
$B \to C[y_i, z]/(y_iz) \to \Lambda$, où $C[y_i, z]/(y_iz) \to \Lambda$
est le morphisme donné $C \to \Lambda$, envoie $z$ sur $\lambda$ et envoie
$y_i$ sur l'image de $a_i^N$ dans $\Lambda$.

\medskip\noindent
Nous affirmons que $B$ est une solution pour $R \to A \to \Lambda \supset \mathfrak q$.
Remarquons d'abord que $B_z$ est isomorphe à $C[z, z^{-1}, t_{ij}]$.
Choisissons $c \in H_{C/R}$ dont l'image dans $\Lambda$ n'appartient pas à
$\mathfrak q$. Alors $B_{zc}$ est lisse sur $R$. D'autre part,
$B_{y_\ell} \cong A[x_i, y_i, y_\ell^{-1}, t_{ij}, i \not = \ell]$
qui est lisse sur $A$. Nous voyons ainsi que $zc$ et les $a_\ell y_\ell$
(les composées de morphismes lisses sont lisses) sont tous
des éléments de $H_{B/R}$. Cela démontre le lemme.
\end{proof}




\section{Corps résiduels séparables}
\label{section-separable}

\noindent
Dans cette section, nous expliquons comment résoudre un problème local dans le cas
d'une extension séparable de corps résiduels.

\begin{lemma}[Ogoma]
\label{lemma-ogoma}
Soit $A$ un anneau noethérien et soit $M$ un $A$-module de type fini.
Soit $S \subset A$ une partie multiplicative. Si $\pi \in A$ et
$\Ker(\pi : S^{-1}M \to S^{-1}M) =
\Ker(\pi^2 : S^{-1}M \to S^{-1}M)$
alors il existe $s \in S$ tel que, pour tout $n > 0$, nous ayons
$\Ker(s^n\pi : M \to M) = \Ker((s^n\pi)^2 : M \to M)$.
\end{lemma}

\begin{proof}
Posons $K = \Ker(\pi : M \to M)$ et
$K' = \{m \in M \mid \pi^2 m = 0\text{ dans }S^{-1}M\}$ et
$Q = K'/K$. Remarquons que $S^{-1}Q = 0$ par hypothèse. Comme $A$
est noethérien, $Q$ est un $A$-module de type fini.
Il existe donc $s \in S$ tel que $s$ annule $Q$.
Alors $s$ convient.
\end{proof}

\begin{lemma}
\label{lemma-find-sequence}
Soit $\Lambda$ un anneau noethérien. Soit $I \subset \Lambda$ un idéal.
Soit $I \subset \mathfrak q$, où ce dernier est un idéal premier. Soient $n, e$ des entiers positifs.
Supposons que $\mathfrak q^n\Lambda_\mathfrak q \subset I\Lambda_\mathfrak q$
et que $\Lambda_\mathfrak q$ soit un anneau local régulier de dimension $d$.
Alors il existe des éléments $\pi_1, \ldots, \pi_d \in \Lambda$ tels que
\begin{enumerate}
\item $(\pi_1, \ldots, \pi_d)\Lambda_\mathfrak q =
\mathfrak q\Lambda_\mathfrak q$,
\item $\pi_1^n, \ldots, \pi_d^n \in I$, et
\item pour $i = 1, \ldots, d$, nous avons
$$
\text{Ann}_{\Lambda/(\pi_1^e, \ldots, \pi_{i - 1}^e)\Lambda}(\pi_i) =
\text{Ann}_{\Lambda/(\pi_1^e, \ldots, \pi_{i - 1}^e)\Lambda}(\pi_i^2).
$$
\end{enumerate}
\end{lemma}

\begin{proof}
Posons $S = \Lambda \setminus \mathfrak q$, de sorte que
$\Lambda_\mathfrak q = S^{-1}\Lambda$.
Choisissons d'abord $\pi_1, \ldots, \pi_d$ satisfaisant (1), ce qui est possible
puisque $\Lambda_\mathfrak q$ est régulier. Par hypothèse,
$\pi_i^n \in I\Lambda_\mathfrak q$. Il existe donc
$s_1, \ldots, s_d \in S$ tels que $s_i\pi_i^n \in I$.
En remplaçant $\pi_i$ par $s_i\pi_i$, nous obtenons (2).
Remarquons que (1) et (2) sont préservées par toute multiplication ultérieure par des éléments de $S$.
Supposons que (3) soit satisfaite pour $i = 1, \ldots, t$, pour un certain
$t \in \{0, \ldots, d\}$. Remarquons que
$\pi_1, \ldots, \pi_d$ est une suite régulière dans $S^{-1}\Lambda$ ; voir
Algèbre, lemme \ref{algebra-lemma-regular-ring-CM}.
En particulier, $\pi_1^e, \ldots, \pi_t^e, \pi_{t + 1}$ est une
suite régulière dans $S^{-1}\Lambda = \Lambda_\mathfrak q$ d'après
Algèbre, lemme \ref{algebra-lemma-regular-sequence-powers}.
Nous obtenons donc
$$
\text{Ann}_{S^{-1}\Lambda/(\pi_1^e, \ldots, \pi_{i - 1}^e)}(\pi_i) =
\text{Ann}_{S^{-1}\Lambda/(\pi_1^e, \ldots, \pi_{i - 1}^e)}(\pi_i^2).
$$
Nous obtenons ainsi (3) pour $i = t + 1$ après avoir remplacé $\pi_{t + 1}$ par $s\pi_{t + 1}$
pour un certain $s \in S$, d'après le Lemme \ref{lemma-ogoma}. Par récurrence sur $t$, ceci
produit une suite satisfaisant (1), (2) et (3).
\end{proof}

\begin{lemma}
\label{lemma-resolve-special}
Soit $k \to A \to \Lambda \supset \mathfrak q$ comme dans la
Situation \ref{situation-local}, où
\begin{enumerate}
\item $k$ est un corps,
\item $\Lambda$ est noethérien,
\item $\mathfrak q$ est minimal au-dessus de $\mathfrak h_A$,
\item $\Lambda_\mathfrak q$ est un anneau local régulier, et
\item l'extension de corps $\kappa(\mathfrak q)/k$ est séparable.
\end{enumerate}
Alors $k \to A \to \Lambda \supset \mathfrak q$ peut être résolu.
\end{lemma}

\begin{proof}
Posons $d = \dim \Lambda_\mathfrak q$. Posons $R = k[x_1, \ldots, x_d]$.
Choisissons $n > 0$ tel que
$\mathfrak q^n\Lambda_\mathfrak q \subset \mathfrak h_A\Lambda_\mathfrak q$
ce qui est possible puisque $\mathfrak q$ est minimal au-dessus de $\mathfrak h_A$.
Choisissons des générateurs $a_1, \ldots, a_r$ de $H_{A/k}$. Posons
$$
B = A[x_1, \ldots, x_d, z_{ij}]/(x_i^n - \sum z_{ij}a_j)
$$
Chaque $B_{a_j}$ est lisse sur $R$ : c'est une algèbre de polynômes
sur $A_{a_j}[x_1, \ldots, x_d]$, et $A_{a_j}$ est lisse sur $k$.
Par conséquent, $B_{x_i}$ est lisse sur $R$. Soit $B \to C$ le morphisme de $R$-algèbres
construit dans le Lemme \ref{lemma-improve-presentation},
qui est muni d'une rétraction de $R$-algèbres $C \to B$. En particulier,
nous obtenons un morphisme $C \to \Lambda$ s'insérant dans le diagramme ci-dessous.
Par construction, $C_{x_i}$ est une $R$-algèbre lisse dont
$\Omega_{C_{x_i}/R}$ est libre. Il existe donc $c > 0$
tel que $x_i^c$ soit strictement standard dans $C/R$ ; voir le
Lemme \ref{lemma-compare-standard}.
Choisissons maintenant $\pi_1, \ldots, \pi_d \in \Lambda$ comme dans le
Lemme \ref{lemma-find-sequence},
où $n = n$, $e = 8c$, $\mathfrak q = \mathfrak q$ et $I = \mathfrak h_A$.
Écrivons $\pi_i^n = \sum \lambda_{ij} a_j$ pour certains $\lambda_{ij} \in \Lambda$.
Il existe un morphisme $B \to \Lambda$ donné par $x_i \mapsto \pi_i$
et $z_{ij} \mapsto \lambda_{ij}$. Posons $R = k[x_1, \ldots, x_d]$.
Considérons le diagramme
$$
\xymatrix{
R \ar[r] & B \ar[rd] \\
k \ar[u] \ar[r] & A \ar[u] \ar[r] & \Lambda
}
$$
Appliquons maintenant le
Lemme \ref{lemma-lift-solution}
à $R \to C \to \Lambda \supset \mathfrak q$
et à la suite d'éléments $x_1^c, \ldots, x_d^c$ de $R$.
L'hypothèse (2) est claire. L'hypothèse (1) est satisfaite pour $R$
par inspection et pour $\Lambda$ grâce à notre choix de
$\pi_1, \ldots, \pi_d$. (Remarquons que si
$\text{Ann}_\Lambda(\pi) = \text{Ann}_\Lambda(\pi^2)$, alors nous avons
$\text{Ann}_\Lambda(\pi) = \text{Ann}_\Lambda(\pi^c)$ pour tout $c > 0$.)
Il suffit donc de résoudre
$$
R/(x_1^e, \ldots, x_d^e) \to C/(x_1^e, \ldots, x_d^e) \to
\Lambda/(\pi_1^e, \ldots, \pi_d^e) \supset
\mathfrak q/(\pi_1^e, \ldots, \pi_d^e)
$$
pour $e = 8c$. D'après le
Lemme \ref{lemma-delocalize-height-zero},
il suffit de résoudre cette situation après localisation en $\mathfrak q$.
Mais puisque $x_1, \ldots, x_d$ ont pour images une suite régulière
dans $\Lambda_\mathfrak q$, nous voyons que $R_\mathfrak p \to \Lambda_\mathfrak q$
est plat ; voir Algèbre, lemme \ref{algebra-lemma-flat-over-regular}. Par conséquent,
$$
R_\mathfrak p/(x_1^e, \ldots, x_d^e) \to
\Lambda_\mathfrak q/(\pi_1^e, \ldots, \pi_d^e)
$$
est un morphisme d'anneaux plat entre anneaux locaux artiniens.
De plus, ce morphisme induit par hypothèse une extension séparable
de corps résiduels. Ce morphisme est donc une colimite filtrante
d'algèbres lisses d'après
Algèbre, lemme \ref{algebra-lemma-colimit-syntomic},
et la Proposition \ref{proposition-lift}.
L'existence de la solution souhaitée résulte de
Algèbre, lemme \ref{algebra-lemma-when-colimit}.
\end{proof}
\section{Extensions inséparables de corps résiduels}
\label{section-inseparable}

\noindent
Dans cette section, nous expliquons comment résoudre un problème local dans le cas
d'une extension inséparable de corps résiduels.

\begin{lemma}
\label{lemma-helper}
Soit $k$ un corps de caractéristique $p > 0$.
Soit $(\Lambda, \mathfrak m, K)$ une $k$-algèbre locale artinienne.
Supposons que $\dim H_1(L_{K/k}) < \infty$.
Alors $\Lambda$ est une colimite filtrante de $k$-algèbres
locales artiniennes $A$ telles que chaque morphisme $A \to \Lambda$ soit plat, que
$\mathfrak m_A \Lambda = \mathfrak m$ et que
$A$ soit essentiellement de type fini sur $k$.
\end{lemma}

\begin{proof}
Remarquons que la platitude de $A \to \Lambda$ implique que $A \to \Lambda$
est injectif ; le lemme affirme donc en réalité que $\Lambda$ est une
réunion filtrante de sous-anneaux $A \subset \Lambda$ de ce type.
Soit $n$ le plus petit entier tel que $\mathfrak m^n = 0$.
Nous démontrons ce lemme par récurrence sur $n$. Le cas $n = 1$ est clair,
car une extension de corps est la réunion de ses sous-extensions de type fini.

\medskip\noindent
Choisissons $\lambda_1, \ldots, \lambda_d \in \mathfrak m$ qui engendrent
$\mathfrak m$. Comme $K$ est formellement lisse sur $\mathbf{F}_p$ (voir
Algèbre, lemme \ref{algebra-lemma-formally-smooth-extensions-easy}), nous pouvons
trouver un morphisme d'anneaux $\sigma : K \to \Lambda$ qui est une section du
morphisme quotient $\Lambda \to K$. En général, $\sigma$ n'est {\bf pas}
un morphisme de $k$-algèbres. Étant donné $\sigma$, définissons
$$
\Psi_\sigma : K[x_1, \ldots, x_d] \longrightarrow \Lambda
$$
en utilisant $\sigma$ sur les éléments de $K$ et en envoyant $x_i$ sur $\lambda_i$.
Assertion : il existe un $\sigma : K \to \Lambda$
et un sous-corps $k \subset F \subset K$ de type fini sur $k$
tels que l'image de $k$ dans $\Lambda$ soit contenue dans
$\Psi_\sigma(F[x_1, \ldots, x_d])$.

\medskip\noindent
Nous démontrons l'assertion par récurrence sur le plus petit entier $n$ tel que
$\mathfrak m^n = 0$. Elle est claire pour $n = 1$. Si $n > 1$, posons
$I = \mathfrak m^{n - 1}$ et $\Lambda' = \Lambda/I$.
Par récurrence, nous pouvons supposer donnés
$\sigma' : K \to \Lambda'$ et $k \subset F' \subset K$, le corps intermédiaire
étant de type fini sur le corps de base, tels que l'image de $k \to \Lambda \to \Lambda'$
soit contenue dans $A' = \Psi_{\sigma'}(F'[x_1, \ldots, x_d])$.
Notons $\tau' : k \to A'$ le morphisme induit.
Choisissons un relèvement $\sigma : K \to \Lambda$ de $\sigma'$ (cela est possible
par la lissité formelle de $K/\mathbf{F}_p$ mentionnée ci-dessus).
Pour usage ultérieur, remarquons que nous pouvons remplacer $\sigma$ par
$\sigma + D$ pour une certaine dérivation $D : K \to I$.
Posons $A = F'[x_1, \ldots, x_d]/(x_1, \ldots, x_d)^n$.
Alors $\Psi_\sigma$ induit un morphisme d'anneaux
$\Psi_\sigma : A \to \Lambda$. La composition avec le
morphisme quotient $\Lambda \to \Lambda'$ induit un morphisme surjectif
$A \to A'$ de noyau nilpotent.
Choisissons un relèvement $\tau : k \to A$ de $\tau'$ (possible puisque $k/\mathbf{F}_p$
est formellement lisse). Nous obtenons ainsi deux morphismes $k \to \Lambda$, à savoir
$\Psi_\sigma \circ \tau : k \to \Lambda$ et le morphisme donné $i : k \to \Lambda$.
Ces morphismes coïncident modulo $I$ ; leur différence est donc une
dérivation $\theta = i - \Psi_\sigma \circ \tau : k \to I$.
Remarquons que si nous remplaçons $\sigma$ par $\sigma + D$, alors nous remplaçons
$\theta$ par $\theta - D|_k$.

\medskip\noindent
Choisissons une famille d'éléments $\{y_j\}_{j \in J}$ de $k$ dont les différentielles
$\text{d}y_j$ forment une base de $\Omega_{k/\mathbf{F}_p}$. La suite de Jacobi-Zariski
pour $\mathbf{F}_p \subset k \subset K$ est
$$
0 \to H_1(L_{K/k}) \to \Omega_{k/\mathbf{F}_p} \otimes K \to
\Omega_{K/\mathbf{F}_p} \to \Omega_{K/k} \to 0
$$
Comme $\dim H_1(L_{K/k}) < \infty$, nous pouvons trouver un sous-ensemble fini $J_0 \subset J$
tel que l'image du premier morphisme soit contenue dans
$\bigoplus_{j \in J_0} K\text{d}y_j$. Par conséquent, les images des éléments
$\text{d}y_j$, $j \in J \setminus J_0$, sont $K$-linéairement indépendantes
dans $\Omega_{K/\mathbf{F}_p}$. Nous pouvons donc choisir
$D : K \to I$ tel que $\theta - D|_k = \xi \circ \text{d}$,
où $\xi$ est la composition
$$
\Omega_{k/\mathbf{F}_p} = \bigoplus\nolimits_{j \in J} k \text{d}y_j
\longrightarrow \bigoplus\nolimits_{j \in J_0} k \text{d}y_j
\longrightarrow I
$$
Posons $f_j = \xi(\text{d}y_j) \in I$ pour $j \in J_0$.
Remplaçons $\sigma$ par $\sigma + D$ comme ci-dessus. Nous voyons alors que
$\theta(a) = \sum_{j \in J_0} a_j f_j$ pour $a \in k$, où
$\text{d}a = \sum a_j \text{d}y_j$ dans $\Omega_{k/\mathbf{F}_p}$.
Remarquons que $I$ est engendré par les monômes
$\lambda^E = \lambda_1^{e_1} \ldots \lambda_d^{e_d}$ de
degré total $|E| = \sum e_i = n - 1$ en $\lambda_1, \ldots, \lambda_d$.
Écrivons $f_j = \sum_E c_{j, E} \lambda^E$ avec $c_{j, E} \in K$.
Remplaçons $F'$ par $F = F'(c_{j, E})$. L'assertion est alors démontrée.

\medskip\noindent
Choisissons $\sigma$ et $F$ comme dans l'assertion. Le noyau de $\Psi_\sigma$ est
engendré par un nombre fini de polynômes
$g_1, \ldots, g_t \in K[x_1, \ldots, x_d]$ et, après avoir agrandi $F$ en lui adjoignant
un nombre fini d'éléments, nous pouvons supposer que leurs coefficients appartiennent à $F$.
Dans ce cas, il est clair que le morphisme
$A = F[x_1, \ldots, x_d]/(g_1, \ldots, g_t) \to
K[x_1, \ldots, x_d]/(g_1, \ldots, g_t) = \Lambda$ est plat.
D'après l'assertion, $A$ est une sous-$k$-algèbre de $\Lambda$.
Il est clair que $\Lambda$ est la colimite filtrante de ces
algèbres, puisque $K$ est la réunion filtrante des sous-corps $F$.
Enfin, ces algèbres sont essentiellement de type fini sur $k$ d'après
Algèbre, lemme
\ref{algebra-lemma-essentially-of-finite-type-into-artinian-local}.
\end{proof}

\begin{lemma}
\label{lemma-solution-modulo}
Soit $k$ un corps de caractéristique $p > 0$.
Soit $\Lambda$ une $k$-algèbre noethérienne géométriquement régulière.
Soit $\mathfrak q \subset \Lambda$ un idéal premier.
Soit $n \geq 1$ un entier et soit
$E \subset \Lambda_\mathfrak q/\mathfrak q^n\Lambda_\mathfrak q$
un sous-ensemble fini.
Alors nous pouvons trouver $m \geq 0$ et
$\varphi : k[y_1, \ldots, y_m] \to \Lambda$ ayant les propriétés suivantes :
\begin{enumerate}
\item en posant $\mathfrak p = \varphi^{-1}(\mathfrak q)$, nous avons
$\mathfrak q\Lambda_\mathfrak q = \mathfrak p \Lambda_\mathfrak q$
et $k[y_1, \ldots, y_m]_\mathfrak p \to \Lambda_\mathfrak q$ est plat,
\item il existe une factorisation par des morphismes d'anneaux locaux artiniens
$$
k[y_1, \ldots, y_m]_\mathfrak p/\mathfrak p^n k[y_1, \ldots, y_m]_\mathfrak p
\to D \to
\Lambda_\mathfrak q/\mathfrak q^n\Lambda_\mathfrak q
$$
où la première flèche est essentiellement lisse et la seconde est plate,
\item $E$ est contenu dans $D$ modulo $\mathfrak q^n\Lambda_\mathfrak q$.
\end{enumerate}
\end{lemma}

\begin{proof}
Posons $\bar \Lambda = \Lambda_\mathfrak q/\mathfrak q^n\Lambda_\mathfrak q$.
Remarquons que $\dim H_1(L_{\kappa(\mathfrak q)/k}) < \infty$ d'après
Compléments d'algèbre, proposition
\ref{more-algebra-proposition-characterization-geometrically-regular}.
Choisissons $A \subset \bar \Lambda$ contenant $E$, telle que $A$ soit locale
artinienne, essentiellement de type fini sur $k$, que le morphisme
$A \to \bar \Lambda$ soit plat et que $\mathfrak m_A$ engendre l'idéal maximal
de $\bar \Lambda$ ; voir le Lemme \ref{lemma-helper}.
Notons $F = A/\mathfrak m_A$ le corps résiduel, de sorte que $k \subset F \subset K$.
Choisissons $\lambda_1, \ldots, \lambda_t \in \Lambda$ dont les images
appartiennent à $A$ dans $\bar \Lambda$ et tels que, de plus, les images
de $\text{d}\lambda_1, \ldots, \text{d}\lambda_t$ forment une base
de $\Omega_{F/k}$. Considérons le morphisme
$\varphi' : k[y_1, \ldots, y_t] \to \Lambda$ qui envoie $y_j$ sur $\lambda_j$.
Posons $\mathfrak p' = (\varphi')^{-1}(\mathfrak q)$. D'après
Compléments d'algèbre, lemme
\ref{more-algebra-lemma-geometrically-regular-over-field}
le morphisme d'anneaux $k[y_1, \ldots, y_t]_{\mathfrak p'} \to \Lambda_\mathfrak q$
est plat et $\Lambda_\mathfrak q/\mathfrak p' \Lambda_\mathfrak q$ est
régulier. Nous pouvons donc choisir d'autres éléments
$\lambda_{t + 1}, \ldots, \lambda_m \in \Lambda$
dont les images appartiennent à $A \subset \bar \Lambda$ et dont les
images forment un système régulier de paramètres de
$\Lambda_\mathfrak q/\mathfrak p' \Lambda_\mathfrak q$.
Nous obtenons $\varphi : k[y_1, \ldots, y_m] \to \Lambda$ ayant
la propriété (1), tel que
$k[y_1, \ldots, y_m]_\mathfrak p/\mathfrak p^n k[y_1, \ldots, y_m]_\mathfrak p
\to \bar\Lambda$
se factorise par $A$. Ainsi,
$k[y_1, \ldots, y_m]_\mathfrak p/\mathfrak p^n k[y_1, \ldots, y_m]_\mathfrak p
\to A$ est plat d'après
Algèbre, lemme \ref{algebra-lemma-flatness-descends-more-general}.
Par construction, l'extension de corps résiduels $F/\kappa(\mathfrak p)$
est de type fini et $\Omega_{F/\kappa(\mathfrak p)} = 0$. Cette extension est
donc finie séparable d'après
Compléments d'algèbre, lemme \ref{more-algebra-lemma-cartier-equality}.
Ainsi,
$k[y_1, \ldots, y_m]_\mathfrak p/\mathfrak p^n k[y_1, \ldots, y_m]_\mathfrak p
\to A$
est fini d'après Algèbre, lemme
\ref{algebra-lemma-essentially-of-finite-type-into-artinian-local}.
Enfin, nous concluons qu'il est étale d'après
Algèbre, lemme \ref{algebra-lemma-characterize-etale}.
Comme un morphisme d'anneaux étale est certainement essentiellement lisse, la démonstration est achevée.
\end{proof}

\begin{lemma}
\label{lemma-enlarge-solution-modulo}
Soient $\varphi : k[y_1, \ldots, y_m] \to \Lambda$, $n$, $\mathfrak q$,
$\mathfrak p$ et
$$
k[y_1, \ldots, y_m]_\mathfrak p/\mathfrak p^n \to
D \to \Lambda_\mathfrak q/\mathfrak q^n \Lambda_\mathfrak q
$$
comme dans le Lemme \ref{lemma-solution-modulo}. Alors, pour tout
$\lambda \in \Lambda \setminus \mathfrak q$
il existe un entier $q > 0$ et une factorisation
$$
k[y_1, \ldots, y_m]_\mathfrak p/\mathfrak p^n \to
D \to D' \to \Lambda_\mathfrak q/\mathfrak q^n \Lambda_\mathfrak q
$$
tels que $D \to D'$ soit un morphisme essentiellement lisse d'anneaux locaux artiniens,
que la dernière flèche soit plate et que $\lambda^q$ appartienne à $D'$.
\end{lemma}

\begin{proof}
Posons $\bar \Lambda = \Lambda_\mathfrak q/\mathfrak q^n\Lambda_\mathfrak q$.
Soit $\bar \lambda$ l'image de $\lambda$ dans $\bar \Lambda$.
Soit $\alpha \in \kappa(\mathfrak q)$ l'image de $\lambda$ dans le
corps résiduel.
Soit $k \subset F \subset \kappa(\mathfrak q)$ le corps résiduel de $D$.
Si $\alpha$ appartient à $F$, nous pouvons trouver
$x \in D$ tel que $x \bar\lambda = 1 \bmod \mathfrak q$. Par conséquent,
$(x \bar \lambda)^q = 1 \bmod (\mathfrak q)^q$ si $q$ est divisible par $p$.
Ainsi, $\bar\lambda^q$ appartient à $D$. Si $\alpha$ est
transcendant sur $F$, nous pouvons prendre $D' = (D[\bar \lambda])_\mathfrak m$
égal au sous-anneau engendré par $D$ et $\bar \lambda$, localisé
en $\mathfrak m = D[\bar \lambda] \cap \mathfrak q \bar \Lambda$.
Cela convient car, dans ce cas, $D[\bar \lambda]$ est en fait une algèbre de polynômes
sur $D$. Enfin, si $\lambda \bmod \mathfrak q$ est
algébrique sur $F$, nous pouvons trouver une puissance de $p$, notée $q$, telle que
$\alpha^q$ soit séparable algébrique sur $F$ ; voir
Corps, section \ref{fields-section-algebraic}.
Remarquons que $D$ et $\bar\Lambda$ sont des anneaux locaux henséliens ; voir
Algèbre, lemme \ref{algebra-lemma-local-dimension-zero-henselian}.
Soit $D \to D'$ une extension étale finie
dont l'extension de corps résiduels est $F(\alpha^q)/F$ ; voir
Algèbre, lemme \ref{algebra-lemma-henselian-cat-finite-etale}.
Comme $\bar\Lambda$ est hensélien et que $F(\alpha^q)$ est contenu
dans son corps résiduel, nous pouvons trouver une factorisation
$D' \to \bar \Lambda$. La première partie de l'argument montre que
$\bar\lambda^{qq'} \in D'$ pour un certain $q' > 0$.
\end{proof}

\begin{lemma}
\label{lemma-resolve-general}
Soit $k \to A \to \Lambda \supset \mathfrak q$ comme dans la
Situation \ref{situation-local}, où
\begin{enumerate}
\item $k$ est un corps de caractéristique $p > 0$,
\item $\Lambda$ est noethérien et géométriquement régulier sur $k$,
\item $\mathfrak q$ est minimal au-dessus de $\mathfrak h_A$.
\end{enumerate}
Alors $k \to A \to \Lambda \supset \mathfrak q$ peut être résolu.
\end{lemma}

\begin{proof}
Le lemme se démontre par les étapes suivantes, dans l'ordre indiqué.
Nous justifierons chacune de ces étapes ci-dessous.
\begin{enumerate}
\item
\label{item-power}
Choisissons un entier $N > 0$ tel que
$\mathfrak q^N\Lambda_\mathfrak q \subset H_{A/k}\Lambda_\mathfrak q$.
\item
\label{item-generators}
Choisissons des générateurs $a_1, \ldots, a_t \in A$ de l'idéal $H_{A/k}$.
\item
\label{item-dimension}
Posons $d = \dim(\Lambda_\mathfrak q)$.
\item
\label{item-standardizer}
Posons $B = A[x_1, \ldots, x_d, z_{ij}]/(x_i^{2N} - \sum z_{ij}a_j)$.
\item
\label{item-elkik}
Considérons $B$ comme une $k[x_1, \ldots, x_d]$-algèbre et soit
$B \to C$ comme dans le
Lemme \ref{lemma-improve-presentation}.
Nous obtenons aussi une section $C \to B$.
\item
\label{item-strictly-standard}
Choisissons $c > 0$ tel que chaque $x_i^c$
soit strictement standard dans $C$ sur $k[x_1, \ldots, x_d]$.
\item
\label{item-set-n}
Posons $n = N + dc$ et $e = 8c$.
\item
\label{item-set-E}
Soit $E \subset \Lambda_\mathfrak q/\mathfrak q^n\Lambda_\mathfrak q$
l'ensemble des images de générateurs de $A$ comme $k$-algèbre.
\item
\label{item-NP}
Choisissons un entier $m$, un morphisme de $k$-algèbres
$\varphi : k[y_1, \ldots, y_m] \to \Lambda$
et une factorisation par des anneaux locaux artiniens
$$
k[y_1, \ldots, y_m]_\mathfrak p/\mathfrak p^n k[y_1, \ldots, y_m]_\mathfrak p
\to D \to
\Lambda_\mathfrak q/\mathfrak q^n\Lambda_\mathfrak q
$$
telle que la première flèche soit essentiellement lisse, la seconde plate,
que $E$ soit contenu dans $D$, et, avec $\mathfrak p = \varphi^{-1}(\mathfrak q)$,
que le morphisme $k[y_1, \ldots, y_m]_\mathfrak p \to \Lambda_\mathfrak q$ soit
plat et que $\mathfrak p \Lambda_\mathfrak q = \mathfrak q \Lambda_\mathfrak q$.
\item
\label{item-choose-pii}
Choisissons $\pi_1, \ldots, \pi_d \in \mathfrak p$ dont les images forment un
système régulier de paramètres de $k[y_1, \ldots, y_m]_\mathfrak p$.
\item
\label{item-set-R}
Soit $R = k[y_1, \ldots, y_m, t_1, \ldots, t_d]$ et $\gamma_i = \pi_i t_i$.
\item
\label{item-modify-pii}
Si nécessaire, modifions le choix des $\pi_i$ de sorte que,
pour $i = 1, \ldots, d$, nous ayons
$$
\text{Ann}_{R/(\gamma_1^e, \ldots, \gamma_{i - 1}^e)R}(\gamma_i)
=
\text{Ann}_{R/(\gamma_1^e, \ldots, \gamma_{i - 1}^e)R}(\gamma_i^2)
$$
\item
\label{item-choose-deltai}
Il existe $\delta_1, \ldots, \delta_d \in \Lambda$,
$\delta_i \not \in \mathfrak q$, et une factorisation
$D \to D' \to \Lambda_\mathfrak q/\mathfrak q^n\Lambda_\mathfrak q$
avec $D'$ local artinien, $D \to D'$ essentiellement lisse, la dernière flèche
$D' \to \Lambda_\mathfrak q/\mathfrak q^n\Lambda_\mathfrak q$
étant plate, et tels que, si l'on pose $\pi_i' = \delta_i \pi_i$, on ait pour
$i = 1, \ldots, d$ :
\begin{enumerate}
\item $(\pi_i')^{2N} = \sum a_j\lambda_{ij}$ dans $\Lambda$, où
$\lambda_{ij} \bmod \mathfrak q^n\Lambda_\mathfrak q$ est un élément de $D'$,
\item $\text{Ann}_{\Lambda/({\pi'}_1^e, \ldots, {\pi'}_{i - 1}^e)}({\pi'}_i) =
\text{Ann}_{\Lambda/({\pi'}_1^e, \ldots, {\pi'}_{i - 1}^e)}({\pi'}_i^2)$,
\item $\delta_i \bmod \mathfrak q^n\Lambda_\mathfrak q$ est un élément de $D'$.
\end{enumerate}
\item
\label{item-map-B-Lambda}
Définissons $B \to \Lambda$ en envoyant $x_i$ sur $\pi'_i$ et
$z_{ij}$ sur le $\lambda_{ij}$ trouvé ci-dessus. Définissons $C \to \Lambda$
en composant le morphisme $B \to \Lambda$ avec la rétraction $C \to B$.
\item
\label{item-set-map-R}
Définissons $R \to \Lambda$ par $\varphi$ sur $k[y_1, \ldots, y_m]$
et en envoyant $t_i$ sur $\delta_i$. Introduisons en outre un morphisme
$$
k[x_1, \ldots, x_d]
\longrightarrow
R = k[y_1, \ldots, y_m, t_1, \ldots, t_d]
$$
en envoyant $x_i$ sur $\gamma_i = \pi_i t_i$.
\item
\label{item-first-solve}
Il suffit de résoudre
$$
R
\to
C \otimes_{k[x_1, \ldots, x_d]} R
\to
\Lambda \supset \mathfrak q
$$
\item
\label{item-set-I}
Posons $I = (\gamma_1^e, \ldots, \gamma_d^e) \subset R$.
\item
\label{item-second-resolve}
Il suffit de résoudre
$$
R/I
\to
C \otimes_{k[x_1, \ldots, x_d]} R/I
\to
\Lambda/I\Lambda \supset \mathfrak q/I\Lambda
$$
\item
\label{item-set-r}
Notons $\mathfrak r \subset R = k[y_1, \ldots, y_m, t_1, \ldots, t_d]$
l'image inverse de $\mathfrak q$.
\item
\label{item-third-resolve}
Il suffit de résoudre
$$
(R/I)_\mathfrak r \to
C \otimes_{k[x_1, \ldots, x_d]} (R/I)_\mathfrak r \to
\Lambda_\mathfrak q/I\Lambda_\mathfrak q
\supset
\mathfrak q\Lambda_\mathfrak q/I\Lambda_\mathfrak q
$$
\item
\label{item-fourth-resolve}
Posons $J = (\pi_1^e, \ldots, \pi_d^e)$ dans $k[y_1, \ldots, y_m]$.
\item
\label{item-fifth-resolve}
Il suffit de résoudre
$$
(R/JR)_\mathfrak p \to
C \otimes_{k[x_1, \ldots, x_d]} (R/JR)_\mathfrak p \to
\Lambda_\mathfrak q/J\Lambda_\mathfrak q
\supset
\mathfrak q\Lambda_\mathfrak q/J\Lambda_\mathfrak q
$$
\item
\label{item-sixth-resolve}
Il suffit de résoudre
$$
(R/\mathfrak p^nR)_\mathfrak p \to
C \otimes_{k[x_1, \ldots, x_d]} (R/\mathfrak p^nR)_\mathfrak p \to
\Lambda_\mathfrak q/\mathfrak q^n\Lambda_\mathfrak q
\supset
\mathfrak q\Lambda_\mathfrak q/\mathfrak q^n\Lambda_\mathfrak q
$$
\item
\label{item-seventh-resolve}
Il suffit de résoudre
$$
(R/\mathfrak p^nR)_\mathfrak p \to
B \otimes_{k[x_1, \ldots, x_d]} (R/\mathfrak p^nR)_\mathfrak p \to
\Lambda_\mathfrak q/\mathfrak q^n\Lambda_\mathfrak q
\supset
\mathfrak q\Lambda_\mathfrak q/\mathfrak q^n\Lambda_\mathfrak q
$$
\item
\label{item-eighth-resolve}
On munit l'anneau $D'[t_1, \ldots, t_d]$ d'une structure de
$R_\mathfrak p/\mathfrak p^nR_\mathfrak p$-algèbre au moyen du morphisme donné
$k[y_1, \ldots, y_m]_\mathfrak p/\mathfrak p^n k[y_1, \ldots, y_m]_\mathfrak p
\to D'$ et en envoyant $t_i$ sur $t_i$. Il suffit de trouver une factorisation
$$
B \otimes_{k[x_1, \ldots, x_d]} (R/\mathfrak p^nR)_\mathfrak p
\to D'[t_1, \ldots, t_d] \to
\Lambda_\mathfrak q/\mathfrak q^n\Lambda_\mathfrak q
$$
où la seconde flèche envoie $t_i$ sur $\delta_i$ et induit le
morphisme donné $D' \to \Lambda_\mathfrak q/\mathfrak q^n\Lambda_\mathfrak q$.
\item
\label{item-done}
Une telle factorisation existe grâce au choix de $D'$ fait ci-dessus.
\end{enumerate}
Nous donnons maintenant la justification de chacune des étapes, en omettant
celles qui ne font qu'introduire des notations.

\medskip\noindent
Pour (\ref{item-power}). Cela est possible puisque $\mathfrak q$ est minimal
au-dessus de $\mathfrak h_A = \sqrt{H_{A/k}\Lambda}$.

\medskip\noindent
Pour (\ref{item-strictly-standard}). Remarquons que $A_{a_j}$ est lisse
sur $k$. Ainsi $B_{a_j}$, qui est isomorphe à une algèbre de polynômes
sur $A_{a_j}[x_1, \ldots, x_d]$, est lisse sur
$k[x_1, \ldots, x_d]$. Par conséquent, $B_{x_i}$ est lisse sur $k[x_1, \ldots, x_d]$.
D'après le Lemme \ref{lemma-improve-presentation},
nous voyons que $C_{x_i}$ est lisse sur $k[x_1, \ldots, x_d]$
et que son module des différentielles de Kähler est libre de type fini. Une certaine puissance de
$x_i$ est donc strictement standard dans $C$ sur $k[x_1, \ldots, x_d]$
d'après le Lemme \ref{lemma-compare-standard}.

\medskip\noindent
Pour (\ref{item-NP}). Cela résulte de l'application du Lemme \ref{lemma-solution-modulo}.

\medskip\noindent
Pour (\ref{item-choose-pii}). Puisque
$k[y_1, \ldots, y_m]_\mathfrak p \to \Lambda_\mathfrak q$
est plat et que $\mathfrak p \Lambda_\mathfrak q = \mathfrak q \Lambda_\mathfrak q$
par construction,
nous avons $\dim(k[y_1, \ldots, y_m]_\mathfrak p) = d$ d'après
Algèbre, lemme \ref{algebra-lemma-dimension-base-fibre-equals-total}.
Nous pouvons donc trouver $\pi_1, \ldots, \pi_d \in \mathfrak p$ formant
un système régulier de paramètres de $k[y_1, \ldots, y_m]_\mathfrak p$.

\medskip\noindent
Pour (\ref{item-modify-pii}). D'après
Algèbre, lemme \ref{algebra-lemma-regular-ring-CM},
toute permutation de la suite $\pi_1, \ldots, \pi_d$ est une
suite régulière dans $k[y_1, \ldots, y_m]_\mathfrak p$. Par conséquent,
$\gamma_1 = \pi_1 t_1, \ldots, \gamma_d = \pi_d t_d$ est une suite
régulière dans
$R_\mathfrak p = k[y_1, \ldots, y_m]_\mathfrak p[t_1, \ldots, t_d]$ ; voir
Algèbre, lemme \ref{algebra-lemma-regular-sequence-in-polynomial-ring}.
Soit $S = k[y_1, \ldots, y_m] \setminus \mathfrak p$, de sorte que
$R_\mathfrak p = S^{-1}R$. Remarquons que $\pi_1, \ldots, \pi_d$
et $\gamma_1, \ldots, \gamma_d$
restent des suites régulières si nous multiplions nos $\pi_i$ par des éléments de $S$.
Supposons que
$$
\text{Ann}_{R/(\gamma_1^e, \ldots, \gamma_{i - 1}^e)R}(\gamma_i)
=
\text{Ann}_{R/(\gamma_1^e, \ldots, \gamma_{i - 1}^e)R}(\gamma_i^2)
$$
soit satisfaite pour $i = 1, \ldots, t$, pour un certain $t \in \{0, \ldots, d\}$. Remarquons que
$\gamma_1^e, \ldots, \gamma_t^e, \gamma_{t + 1}$ est une
suite régulière dans $S^{-1}R$ d'après
Algèbre, lemme \ref{algebra-lemma-regular-sequence-powers}.
Nous voyons donc que
$$
\text{Ann}_{S^{-1}R/(\gamma_1^e, \ldots, \gamma_{i - 1}^e)}(\gamma_i) =
\text{Ann}_{S^{-1}R/(\gamma_1^e, \ldots, \gamma_{i - 1}^e)}(\gamma_i^2).
$$
Nous obtenons ainsi
$$
\text{Ann}_{R/(\gamma_1^e, \ldots, \gamma_t^e)R}(\gamma_{t + 1})
=
\text{Ann}_{R/(\gamma_1^e, \ldots, \gamma_t^e)R}(\gamma_{t + 1}^2)
$$
après avoir remplacé $\pi_{t + 1}$ par $s\pi_{t + 1}$ pour un certain $s \in S$, d'après le
Lemme \ref{lemma-ogoma}. Par récurrence sur $t$, ceci produit la
suite souhaitée.

\medskip\noindent
Pour (\ref{item-choose-deltai}). Posons $S = \Lambda \setminus \mathfrak q$,
de sorte que $\Lambda_\mathfrak q = S^{-1}\Lambda$. Posons
$\bar \Lambda = \Lambda_\mathfrak q/\mathfrak q^n \Lambda_\mathfrak q$.
Supposons que nous ayons un $t \in \{0, \ldots, d\}$, des éléments
$\delta_1, \ldots, \delta_t \in S$ et une factorisation
$D \to D' \to \bar \Lambda$ comme dans (\ref{item-choose-deltai}),
tels que (a), (b), (c) soient satisfaites pour $i = 1, \ldots, t$. Nous avons
$\pi_{t + 1}^N \in H_{A/k}\Lambda_\mathfrak q$
car $\mathfrak q^N \Lambda_\mathfrak q \subset H_{A/k}\Lambda_\mathfrak q$
d'après (\ref{item-power}). Par conséquent,
$\pi_{t + 1}^N \in H_{A/k} \bar\Lambda$, puis
$\pi_{t + 1}^N \in H_{A/k}D'$, puisque $D' \to \bar \Lambda$
est fidèlement plat ; voir
Algèbre, lemme \ref{algebra-lemma-faithfully-flat-universally-injective}.
Rappelons que $H_{A/k} = (a_1, \ldots, a_t)$.
Écrivons $\pi_{t + 1}^N = \sum a_j d_j$ dans $D'$ et choisissons
$c_j \in \Lambda_\mathfrak q$ relevant $d_j \in D'$. Alors
$\pi_{t + 1}^N = \sum c_j a_j + \epsilon$, avec
$\epsilon \in \mathfrak q^n\Lambda_\mathfrak q \subset
\mathfrak q^{n - N}H_{A/k}\Lambda_\mathfrak q$.
Écrivons $\epsilon = \sum a_j c'_j$ pour certains
$c'_j \in \mathfrak q^{n - N}\Lambda_\mathfrak q$.
Ainsi, $\pi_{t + 1}^{2N} = \sum (\pi_{t + 1}^N c_j + \pi_{t + 1}^N c'_j) a_j$.
Remarquons que $\pi_{t + 1}^Nc'_j$ a pour image zéro dans $\bar \Lambda$ ; cette observation
élémentaire mais essentielle garantira plus loin que (a) est satisfaite.
Choisissons maintenant $s \in S$ tel qu'il existe des
$\mu_{t + 1j} \in \Lambda$ vérifiant, d'une part,
$\pi_{t + 1}^N c_j + \pi_{t + 1}^N c'_j = \mu_{t + 1j}/s^{2N}$
dans $S^{-1}\Lambda$ et, d'autre part,
$(s \pi_{t + 1})^{2N} = \sum \mu_{t + 1j}a_j$
dans $\Lambda$ (détail mineur omis). Nous pouvons en outre remplacer $s$ par
une puissance et agrandir $D'$ de sorte que $s$ ait pour image un élément de $D'$.
Avec ces choix, $\mu_{t + 1j}$ a pour image $s^{2N}d_j$, qui est
un élément de $D'$. Remarquons que $\pi_1, \ldots, \pi_d$ forment un système
régulier de paramètres dans $S^{-1}\Lambda$ par notre
choix de $\varphi$. Ainsi, $\pi_1, \ldots, \pi_d$ forment une suite régulière
dans $\Lambda_\mathfrak q$ d'après
Algèbre, lemme \ref{algebra-lemma-regular-ring-CM}.
Il s'ensuit que ${\pi'}_1^e, \ldots, {\pi'}_t^e, s\pi_{t + 1}$ est une
suite régulière dans $S^{-1}\Lambda$ d'après
Algèbre, lemme \ref{algebra-lemma-regular-sequence-powers}.
Nous obtenons donc
$$
\text{Ann}_{S^{-1}\Lambda/({\pi'}_1^e, \ldots, {\pi'}_t^e)}(s\pi_{t + 1}) =
\text{Ann}_{S^{-1}\Lambda/({\pi'}_1^e, \ldots, {\pi'}_t^e)}((s\pi_{t + 1})^2).
$$
Nous pouvons donc appliquer le Lemme \ref{lemma-ogoma} afin de trouver $s' \in S$
tel que
$$
\text{Ann}_{\Lambda/({\pi'}_1^e, \ldots, {\pi'}_t^e)}((s')^qs\pi_{t + 1})
=
\text{Ann}_{\Lambda/({\pi'}_1^e, \ldots, {\pi'}_t^e)}(((s')^qs\pi_{t + 1})^2).
$$
pour tout $q > 0$. D'après le Lemme \ref{lemma-enlarge-solution-modulo},
nous pouvons choisir $q$ et agrandir $D'$ de sorte que $(s')^q$ ait pour image un élément
de $D'$. En posant $\delta_{t + 1} = (s')^qs$, nous concluons que
(a), (b), (c) sont satisfaites pour $i = 1, \ldots, t + 1$. Pour (a), remarquons que
$\lambda_{t + 1j} = (s')^{2Nq}\mu_{t + 1j}$ convient.
Une récurrence sur $t$ achève la démonstration.

\medskip\noindent
Pour (\ref{item-first-solve}). Par construction, le radical de
$H_{(C \otimes_{k[x_1, \ldots, x_d]} R)/R} \Lambda$ contient
$\mathfrak h_A$. En effet, les éléments $a_j \in H_{A/k}$
ont pour images des éléments de $H_{B/k[x_1, \ldots, x_d]}$, puis des éléments
de $H_{C/k[x_1, \ldots, x_d]}$ ; par conséquent, les $a_j \otimes 1$ ont pour images des éléments de
$H_{C \otimes_{k[x_1, \ldots, x_d]} R/R}$. De plus, si nous avons une solution
$C \otimes_{k[x_1, \ldots, x_d]} R \to T \to \Lambda$ de
$$
R
\to
C \otimes_{k[x_1, \ldots, x_d]} R
\to
\Lambda \supset \mathfrak q
$$
alors $H_{T/R} \subset H_{T/k}$, puisque $R$ est lisse sur $k$.
Ainsi, $T$ sera aussi une solution de
la situation initiale $k \to A \to \Lambda \supset \mathfrak q$.

\medskip\noindent
Pour (\ref{item-second-resolve}). Cela résulte de l'application du
Lemme \ref{lemma-lift-solution} à
$R \to C \otimes_{k[x_1, \ldots, x_d]} R
\to \Lambda \supset \mathfrak q$ et à la suite d'éléments
$\gamma_1^c, \ldots, \gamma_d^c$. Remarquons que, puisque les $x_i^c$
sont strictement standard dans $C$ sur $k[x_1, \ldots, x_d]$, les éléments
$\gamma_i^c$ sont strictement standard dans $C \otimes_{k[x_1, \ldots, x_d]} R$
sur $R$ d'après le Lemme \ref{lemma-strictly-standard-base-change}.
L'autre hypothèse du Lemme \ref{lemma-lift-solution} est satisfaite grâce aux étapes
(\ref{item-modify-pii}) et (\ref{item-choose-deltai}).

\medskip\noindent
Pour (\ref{item-third-resolve}). Appliquons le Lemme \ref{lemma-delocalize-height-zero}
à la situation de (\ref{item-second-resolve}). Dans la suite des
arguments, l'anneau d'arrivée est local artinien ; nous cherchons donc
une factorisation par une algèbre $T$ lisse sur l'anneau de départ.

\medskip\noindent
Pour (\ref{item-fifth-resolve}).
Supposons que $C \otimes_{k[x_1, \ldots, x_d]} (R/JR)_\mathfrak p \to
T \to \Lambda_\mathfrak q/J\Lambda_\mathfrak q$ soit une solution de
$$
(R/JR)_\mathfrak p \to
C \otimes_{k[x_1, \ldots, x_d]} (R/JR)_\mathfrak p \to
\Lambda_\mathfrak q/J\Lambda_\mathfrak q
\supset
\mathfrak q\Lambda_\mathfrak q/J\Lambda_\mathfrak q
$$
Alors $C \otimes_{k[x_1, \ldots, x_d]} (R/I)_\mathfrak r \to T_\mathfrak r \to
\Lambda_\mathfrak q/I\Lambda_\mathfrak q$
est une solution de la situation de (\ref{item-third-resolve}).

\medskip\noindent
Pour (\ref{item-sixth-resolve}). Notre entier $n = N + dc$ est assez grand pour que
$\mathfrak p^nk[y_1, \ldots, y_m]_\mathfrak p \subset J_\mathfrak p$
et $\mathfrak q^n \Lambda_\mathfrak q \subset J\Lambda_\mathfrak q$.
Par conséquent, si nous avons une solution
$C \otimes_{k[x_1, \ldots, x_d]} (R/\mathfrak p^nR)_\mathfrak p \to
T \to \Lambda_\mathfrak q/\mathfrak q^n\Lambda_\mathfrak q$
de (\ref{item-fifth-resolve},
nous pouvons prendre $T/JT$ comme solution de
(\ref{item-sixth-resolve}).

\medskip\noindent
Pour (\ref{item-seventh-resolve}). Cela est vrai parce que nous avons une
section $C \to B$ dans la catégorie des $R$-algèbres.

\medskip\noindent
Pour (\ref{item-eighth-resolve}). Cela est vrai parce que $D'$ est
essentiellement lisse sur l'anneau local artinien
$k[y_1, \ldots, y_m]_\mathfrak p/\mathfrak p^n k[y_1, \ldots, y_m]_\mathfrak p$
et que
$$
R_\mathfrak p/\mathfrak p^nR_\mathfrak p =
k[y_1, \ldots, y_m]_\mathfrak p/
\mathfrak p^n k[y_1, \ldots, y_m]_\mathfrak p[t_1, \ldots, t_d].
$$
Ainsi, $D'[t_1, \ldots, t_d]$ est une colimite filtrante de
$R_\mathfrak p/\mathfrak p^nR_\mathfrak p$-algèbres lisses et
$B \otimes_{k[x_1, \ldots, x_d]} (R_\mathfrak p/\mathfrak p^nR_\mathfrak p)$
se factorise par l'une d'elles.

\medskip\noindent
Pour (\ref{item-done}). La dernière subtilité de la démonstration est que nous ne pouvons
pas simplement utiliser le morphisme $B \to D'$ qui envoie $x_i$ sur l'image de $\pi_i'$
dans $D'$ et $z_{ij}$ sur l'image de $\lambda_{ij}$ dans $D'$, 
car nous avons besoin que le diagramme
$$
\xymatrix{
B \ar[r] & D'[t_1, \ldots, t_d] \\
k[x_1, \ldots, x_d] \ar[r] \ar[u] &
R_\mathfrak p/\mathfrak p^nR_\mathfrak p \ar[u]
}
$$
soit commutatif et que la composée
$B \to D'[t_1, \ldots, t_d] \to
\Lambda_\mathfrak q/\mathfrak q^n\Lambda_\mathfrak q$
soit le morphisme de (\ref{item-map-B-Lambda}).
Cela nous impose d'envoyer $x_i$ sur l'image de
$\pi_i t_i$ dans $D'[t_1, \ldots, t_d]$.
Nous envoyons donc $z_{ij}$ sur l'image de
$\lambda_{ij} t_i^{2N} / \delta_i^{2N}$ dans $D'[t_1, \ldots, t_d]$
et tout est clair.
\end{proof}







\section{Le théorème principal}
\label{section-main}

\noindent
Dans cette section, nous concluons la discussion.

\begin{theorem}[Popescu]
\label{theorem-popescu}
Tout morphisme régulier entre anneaux noethériens est une colimite filtrante
de morphismes d'anneaux lisses.
\end{theorem}

\begin{proof}
D'après le lemme \ref{lemma-reduce-to-field},
il suffit de démontrer l'énoncé pour $k \to \Lambda$,
où $\Lambda$ est noethérien et géométriquement régulier sur $k$.
Soit $k \to A \to \Lambda$ une factorisation, avec $A$ une
$k$-algèbre de type fini. Il suffit de construire une factorisation
$A \to B \to \Lambda$, avec $B$ de type fini, telle que
$\mathfrak h_B = \Lambda$ ; voir le lemme \ref{lemma-final-solve}.
Nous pouvons donc procéder par récurrence noethérienne sur l'idéal
$\mathfrak h_A$. Choisissons un idéal premier
$\mathfrak q \supset \mathfrak h_A$ tel que
$\mathfrak q$ soit minimal au-dessus de $\mathfrak h_A$.
Il suffit alors de résoudre $k \to A \to \Lambda \supset \mathfrak q$
(au sens défini dans le texte qui suit la situation \ref{situation-local}).
Si la caractéristique de $k$ est nulle, cela résulte du
lemme \ref{lemma-resolve-special}.
Si la caractéristique de $k$ est $p > 0$, cela résulte du
lemme \ref{lemma-resolve-general}.
\end{proof}




\section{La propriété d'approximation pour les G-anneaux}
\label{section-approximation-G-rings}

\noindent
Soit $R$ un anneau local noethérien. Alors $R$ est un G-anneau si et
seulement si le morphisme d'anneaux $R \to R^\wedge$ est régulier ; voir
Compléments d'algèbre, lemme
\ref{more-algebra-lemma-check-G-ring-maximal-ideals}.
Dans ce cas, l'hensélisation $R^h$ et l'hensélisation stricte
$R^{sh}$ de $R$ sont des G-anneaux ; voir
Compléments d'algèbre, lemme
\ref{more-algebra-lemma-henselization-G-ring}.
En outre, toute algèbre essentiellement de type fini sur un corps, sur un
anneau local complet, sur $\mathbf{Z}$ ou sur un anneau de Dedekind de
caractéristique nulle est un G-anneau ; voir
Compléments d'algèbre, proposition
\ref{more-algebra-proposition-ubiquity-G-ring}.
Nous disposons ainsi d'un vaste choix d'anneaux auxquels s'applique le
résultat ci-dessous.

\medskip\noindent
Soit $R$ un anneau. Soient $f_1, \ldots, f_m \in R[x_1, \ldots, x_n]$.
Soit $S$ une $R$-algèbre. Dans cette situation, nous disons qu'un vecteur
$(a_1, \ldots, a_n) \in S^n$ est une {\it solution dans $S$}
si et seulement si
$$
f_j(a_1, \ldots, a_n) = 0 \text{ dans } S, \text{ pour }
j = 1, \ldots, m
$$
L'une des questions importantes de la géométrie algébrique consiste
naturellement à déterminer quand des systèmes d'équations polynomiales
ont des solutions. Le théorème suivant affirme que l'existence de solutions
dans le complété d'un anneau local noethérien suffit souvent à montrer
qu'il en existe dans l'hensélisation de cet anneau.

\begin{theorem}
\label{theorem-approximation-property}
Soit $R$ un anneau local noethérien. Soient
$f_1, \ldots, f_m \in R[x_1, \ldots, x_n]$.
Supposons que $(a_1, \ldots, a_n) \in (R^\wedge)^n$ soit une solution
dans $R^\wedge$. Si $R$ est un G-anneau hensélien, alors, pour tout entier
$N$, il existe une solution $(b_1, \ldots, b_n) \in R^n$ dans $R$ telle que
$a_i - b_i \in \mathfrak m^NR^\wedge$.
\end{theorem}

\begin{proof}
Choisissons $c_i \in R$ tel que $a_i - c_i \in \mathfrak m^N$.
Choisissons des générateurs $\mathfrak m^N = (d_1, \ldots, d_M)$.
Écrivons $a_i = c_i + \sum a_{i, l} d_l$.
Considérons l'anneau de polynômes $R[x_{i, l}]$ et les éléments
$$
g_j = f_j(c_1 + \sum x_{1, l} d_l , \ldots, c_n + \sum x_{n, l} d_l)
\in R[x_{i, l}]
$$
Le système d'équations $g_j = 0$ admet la solution $(a_{i, l})$.
Supposons que nous puissions montrer que $g_j$ admet une solution
$(b_{i, l})$ dans $R$.
Il s'ensuit que $b_i = c_i + \sum b_{i, l}d_l$ est une solution
de $f_j = 0$ congrue à $a_i$ modulo $\mathfrak m^N$.
Il suffit donc de montrer que l'existence d'une solution sur $R^\wedge$
entraîne l'existence d'une solution sur $R$.

\medskip\noindent
Soit $A \subset R^\wedge$ la sous-$R$-algèbre engendrée par
$a_1, \ldots, a_n$. Puisque nous avons supposé que $R$ est un G-anneau,
c'est-à-dire que $R \to R^\wedge$ est régulier, il existe une factorisation
$$
A \to B \to R^\wedge
$$
où $B$ est lisse sur $R$ ; voir le théorème \ref{theorem-popescu}.
Notons $\kappa = R/\mathfrak m$ le corps résiduel. C'est aussi
le corps résiduel de $R^\wedge$, si bien que nous obtenons un diagramme
commutatif
$$
\xymatrix{
B \ar[rd] \ar@{..>}[r] & R' \ar@{..>}[d] \\
R \ar[r] \ar[u] & \kappa
}
$$
Comme la flèche verticale est lisse,
Compléments d'algèbre, lemme
\ref{more-algebra-lemma-lift-section-smooth-morphism},
il existe un morphisme d'anneaux \'etale $R \to R'$ qui induit un isomorphisme
$R/\mathfrak m \to R'/\mathfrak mR'$ et un morphisme de $R$-algèbres
$B \to R'$ rendant commutatif le diagramme ci-dessus.
Puisque $R$ est hensélien, le morphisme $R \to R'$ possède une section ; voir
Algèbre, lemme \ref{algebra-lemma-characterize-henselian}.
Soit $b_i \in R$ l'image de $a_i$ par les morphismes d'anneaux
$A \to B \to R' \to R$. Puisque tous ces morphismes sont des morphismes de
$R$-algèbres, $(b_1, \ldots, b_n)$ est une solution dans $R$.
\end{proof}

\noindent
Étant donnés un anneau local noethérien $(R, \mathfrak m)$, un morphisme
d'anneaux \'etale $R \to R'$ et un idéal maximal
$\mathfrak m' \subset R'$ au-dessus de $\mathfrak m$, tels que
$\kappa(\mathfrak m) = \kappa(\mathfrak m')$. Nous avons alors les inclusions
$$
R \subset R_{\mathfrak m'} \subset R^h \subset R^\wedge,
$$
d'après Algèbre, lemme \ref{algebra-lemma-henselian-functorial-prepare}, et
Compléments d'algèbre, lemme
\ref{more-algebra-lemma-henselization-noetherian}.

\begin{theorem}
\label{theorem-approximation-property-variant}
Soit $R$ un anneau local noethérien. Soient
$f_1, \ldots, f_m \in R[x_1, \ldots, x_n]$.
Supposons que $(a_1, \ldots, a_n) \in (R^\wedge)^n$ soit une solution.
Si $R$ est un G-anneau, alors, pour tout entier $N$, il existe
\begin{enumerate}
\item un morphisme d'anneaux \'etale $R \to R'$,
\item un idéal maximal $\mathfrak m' \subset R'$ au-dessus de $\mathfrak m$,
\item une solution $(b_1, \ldots, b_n) \in (R')^n$ dans $R'$,
\end{enumerate}
tels que $\kappa(\mathfrak m) = \kappa(\mathfrak m')$ et
$a_i - b_i \in (\mathfrak m')^NR^\wedge$.
\end{theorem}

\begin{proof}
Nous pourrions déduire ce théorème du théorème
\ref{theorem-approximation-property}, en utilisant le fait que
l'hensélisation $R^h$ est un G-anneau d'après
Compléments d'algèbre, lemme
\ref{more-algebra-lemma-henselization-G-ring}, et en écrivant $R^h$ comme une
colimite filtrante d'extensions \'etales $R'$. Nous préférons en donner une
démonstration en reprenant, dans ce cas, celle du théorème précédent.

\medskip\noindent
Choisissons $c_i \in R$ tel que $a_i - c_i \in \mathfrak m^N$.
Choisissons des générateurs $\mathfrak m^N = (d_1, \ldots, d_M)$.
Écrivons $a_i = c_i + \sum a_{i, l} d_l$.
Considérons l'anneau de polynômes $R[x_{i, l}]$ et les éléments
$$
g_j = f_j(c_1 + \sum x_{1, l} d_l , \ldots, c_n + \sum x_{n, l} d_l)
\in R[x_{i, l}]
$$
Le système d'équations $g_j = 0$ admet la solution $(a_{i, l})$.
Supposons que nous puissions montrer que $g_j$ admet une solution
$(b_{i, l})$ dans $R'$, pour un morphisme d'anneaux \'etale $R \to R'$ muni
d'un idéal maximal $\mathfrak m'$ tel que
$\kappa(\mathfrak m) = \kappa(\mathfrak m')$.
Il s'ensuit que $b_i = c_i + \sum b_{i, l}d_l$ est une solution
de $f_j = 0$ congrue à $a_i$ modulo $(\mathfrak m')^N$.
Il suffit donc de montrer que l'existence d'une solution sur $R^\wedge$
entraîne l'existence d'une solution sur une extension d'anneaux \'etale
qui induit une extension triviale des corps résiduels en un idéal premier
au-dessus de $\mathfrak m$.

\medskip\noindent
Soit $A \subset R^\wedge$ la sous-$R$-algèbre engendrée par
$a_1, \ldots, a_n$. Puisque nous avons supposé que $R$ est un G-anneau,
c'est-à-dire que $R \to R^\wedge$ est régulier, il existe une factorisation
$$
A \to B \to R^\wedge
$$
où $B$ est lisse sur $R$ ; voir le théorème \ref{theorem-popescu}.
Notons $\kappa = R/\mathfrak m$ le corps résiduel. C'est aussi
le corps résiduel de $R^\wedge$, si bien que nous obtenons un diagramme
commutatif
$$
\xymatrix{
B \ar[rd] \ar@{..>}[r] & R' \ar@{..>}[d] \\
R \ar[r] \ar[u] & \kappa
}
$$
Comme la flèche verticale est lisse,
Compléments d'algèbre, lemme
\ref{more-algebra-lemma-lift-section-smooth-morphism},
il existe un morphisme d'anneaux \'etale $R \to R'$ qui induit un isomorphisme
$R/\mathfrak m \to R'/\mathfrak mR'$ et un morphisme de $R$-algèbres
$B \to R'$ rendant commutatif le diagramme ci-dessus.
Soit $b_i \in R'$ l'image de $a_i$ par les morphismes d'anneaux
$A \to B \to R'$. Puisque tous ces morphismes sont des morphismes de
$R$-algèbres, $(b_1, \ldots, b_n)$ est une solution dans $R'$.
\end{proof}

\begin{example}
\label{example-describe-henselian}
Soit $(R, \mathfrak m)$ un anneau local noethérien d'hensélisation $R^h$.
Le morphisme entre les complétés $R^\wedge \to (R^h)^\wedge$ est un
isomorphisme ; voir Compléments d'algèbre, lemme
\ref{more-algebra-lemma-henselization-noetherian}.
Comme $R^h$ est lui aussi noethérien (ibid.), nous pouvons considérer $R^h$
comme un sous-anneau de son complété (car la complétion est fidèlement
plate). Nous pouvons ainsi identifier $R^h$ à un sous-anneau de $R^\wedge$.

\medskip\noindent
Cherchons à comprendre quels sont les éléments de $R^\wedge$ qui
appartiennent à $R^h$. Pour simplifier, supposons que $R$ soit un domaine
ayant $K$ pour corps des fractions. Il est clair que tout élément $f$ de $R^h$
est algébrique sur $R$, au sens où il existe une équation de la forme
$a_n f^n + \ldots + a_1 f + a_0 = 0$, pour certains $a_i \in R$ tels que
$n > 0$ et $a_n \not = 0$.

\medskip\noindent
Réciproquement, soient $f \in R^\wedge$, $n \in \mathbf{N}$ et
$a_0, \ldots, a_n \in R$, avec $a_n \not = 0$, tels que
$a_n f^n + \ldots + a_1 f + a_0 = 0$. Si $R$ est un G-anneau, alors, pour
tout $N > 0$, il existe un élément $g \in R^h$ tel que
$a_n g^n + \ldots + a_1 g + a_0 = 0$ et $f - g \in \mathfrak m^N R^\wedge$,
voir le théorème \ref{theorem-approximation-property-variant}.
Nous voudrions en déduire que $f = g$ lorsque $N \gg 0$.
Si ce n'est pas le cas, nous obtenons une infinité de racines $g$ de $P(T)$
dans $R^h$. C'est impossible, car (1) $R^h \subset R^h \otimes_R K$
et (2) $R^h \otimes_R K$ est un produit fini d'extensions de corps de $K$.
En effet, $R \to K$ est injectif et $R \to R^h$ est plat ; par conséquent,
$R^h \to R^h \otimes_R K$ est injectif, et (2) résulte de
Compléments d'algèbre, lemme
\ref{more-algebra-lemma-fibres-henselization}.

\medskip\noindent
Conclusion : si $R$ est un domaine local noethérien ayant $K$ pour corps des
fractions et un G-anneau, alors $R^h \subset R^\wedge$ est l'ensemble des
éléments algébriques sur $K$.
\end{example}

\noindent
Voici une autre variante du théorème principal de cette section.

\begin{lemma}
\label{lemma-approximation-property-variant}
Soit $R$ un anneau noethérien. Soit $\mathfrak p \subset R$ un idéal
premier. Soient
$f_1, \ldots, f_m \in R[x_1, \ldots, x_n]$.
Supposons que $(a_1, \ldots, a_n) \in ((R_\mathfrak p)^\wedge)^n$ soit une
solution. Si $R_\mathfrak p$ est un G-anneau, alors, pour tout entier $N$,
il existe
\begin{enumerate}
\item un morphisme d'anneaux \'etale $R \to R'$,
\item un idéal premier $\mathfrak p' \subset R'$ au-dessus de $\mathfrak p$,
\item une solution $(b_1, \ldots, b_n) \in (R')^n$ dans $R'$,
\end{enumerate}
tels que $\kappa(\mathfrak p) = \kappa(\mathfrak p')$ et
$a_i - b_i \in (\mathfrak p')^N(R'_{\mathfrak p'})^\wedge$.
\end{lemma}

\begin{proof}
D'après le théorème \ref{theorem-approximation-property-variant},
nous pouvons trouver une solution $(b'_1, \ldots, b'_n)$ dans un anneau
$R''$ \'etale sur $R_\mathfrak p$, muni d'un idéal premier
$\mathfrak p''$ au-dessus de $\mathfrak p$, tel que
$\kappa(\mathfrak p) = \kappa(\mathfrak p'')$ et
$a_i - b'_i \in (\mathfrak p'')^N(R''_{\mathfrak p''})^\wedge$.
Nous pouvons écrire
$R'' = R' \otimes_R R_\mathfrak p$
pour une $R$-algèbre \'etale $R'$ (voir Algèbre, lemme
\ref{algebra-lemma-etale}). Après avoir remplacé $R'$ par une localisation
principale si nécessaire, nous pouvons supposer que
$(b'_1, \ldots, b'_n)$ provient d'une solution $(b_1, \ldots, b_n)$
dans $R'$. Posons $\mathfrak p' = R' \cap \mathfrak p''$ ; alors
$R''_{\mathfrak p''} = R'_{\mathfrak p'}$, ce qui achève la démonstration.
\end{proof}



\section{Approximation pour les couples henséliens}
\label{section-approximation-pairs}

\noindent
Nous pouvons généraliser la discussion de la
section \ref{section-approximation-G-rings} au cas des couples henséliens.
Les couples henséliens ont été définis dans
Compléments d'algèbre, section
\ref{more-algebra-section-henselian-pairs}.

\begin{lemma}
\label{lemma-henselian-pair}
\begin{slogan}
Approximation pour les couples henséliens.
\end{slogan}
Soit $(A, I)$ un couple hensélien, avec $A$ noethérien.
Soit $A^\wedge$ le complété $I$-adique de $A$. Supposons que l'une au
moins des conditions suivantes soit satisfaite :
\begin{enumerate}
\item $A \to A^\wedge$ est un morphisme d'anneaux régulier ;
\item $A$ est un G-anneau noethérien ;
\item $(A, I)$ est l'hensélisation
(Compléments d'algèbre, lemme \ref{more-algebra-lemma-henselization})
d'un couple $(B, J)$, où $B$ est un G-anneau noethérien.
\end{enumerate}
Soient $f_1, \ldots, f_m \in A[x_1, \ldots, x_n]$
et $\hat{a}_1, \ldots, \hat{a}_n \in A^\wedge$ tels que
$f_j(\hat{a}_1, \ldots, \hat{a}_n) = 0$
pour $j = 1, \ldots, m$. Alors, pour tout $N \geq 1$, il existe
$a_1, \ldots, a_n \in A$ tels que
$\hat{a}_i - a_i \in I^N$ et $f_j(a_1, \ldots, a_n) = 0$
pour $j = 1, \ldots, m$.
\end{lemma}

\begin{proof}
D'après Compléments d'algèbre, lemme
\ref{more-algebra-lemma-henselization-pair-G-ring}
nous voyons que (3) implique (2). D'après Compléments d'algèbre, lemme
\ref{more-algebra-lemma-map-G-ring-to-completion-regular}
nous voyons que (2) implique (1). Il suffit donc de démontrer le lemme
dans le cas où $A \to A^\wedge$ est un morphisme d'anneaux régulier.

\medskip\noindent
Soient $\hat{a}_1, \ldots, \hat{a}_n$ comme dans l'énoncé du lemme.
D'après le théorème \ref{theorem-popescu}, nous pouvons trouver une
factorisation $A \to B \to A^\wedge$ telle que $A \to B$ soit lisse, ainsi que
des éléments $b_1, \ldots, b_n \in B$ tels que
$f_j(b_1, \ldots, b_n) = 0$ dans $B$.
Notons $\sigma : B \to A^\wedge \to A/I^N$ la composée.
D'après Compléments d'algèbre, lemme
\ref{more-algebra-lemma-lift-section-smooth-morphism}, il existe un morphisme
d'anneaux \'etale $A \to A'$ qui induit un isomorphisme
$A/I^N \to A'/I^NA'$ et un morphisme de $A$-algèbres
$\tilde \sigma : B \to A'$ relevant $\sigma$.
Comme $(A, I)$ est hensélien, il existe un morphisme de $A$-algèbres
$\chi : A' \to A$ ; voir Compléments d'algèbre, lemme
\ref{more-algebra-lemma-characterize-henselian-pair}.
En posant alors $a_i = \chi(\tilde \sigma(b_i))$, nous obtenons une solution.
\end{proof}



\begin{multicols}{2}[\section{Autres chapitres}]
\noindent
Préliminaires
\begin{enumerate}
\item \hyperref[introduction-section-phantom]{Introduction}
\item \hyperref[conventions-section-phantom]{Conventions}
\item \hyperref[sets-section-phantom]{Théorie des ensembles}
\item \hyperref[categories-section-phantom]{Catégories}
\item \hyperref[topology-section-phantom]{Topologie}
\item \hyperref[sheaves-section-phantom]{Faisceaux sur les espaces}
\item \hyperref[sites-section-phantom]{Sites et faisceaux}
\item \hyperref[stacks-section-phantom]{Champs}
\item \hyperref[fields-section-phantom]{Corps}
\item \hyperref[algebra-section-phantom]{Algèbre commutative}
\item \hyperref[brauer-section-phantom]{Groupes de Brauer}
\item \hyperref[homology-section-phantom]{Algèbre homologique}
\item \hyperref[derived-section-phantom]{Catégories dérivées}
\item \hyperref[simplicial-section-phantom]{Méthodes simpliciales}
\item \hyperref[more-algebra-section-phantom]{Compléments d'algèbre}
\item \hyperref[smoothing-section-phantom]{Lissage des morphismes d'anneaux}
\item \hyperref[modules-section-phantom]{Faisceaux de modules}
\item \hyperref[sites-modules-section-phantom]{Modules sur les sites}
\item \hyperref[injectives-section-phantom]{Objets injectifs}
\item \hyperref[cohomology-section-phantom]{Cohomologie des faisceaux}
\item \hyperref[sites-cohomology-section-phantom]{Cohomologie sur les sites}
\item \hyperref[dga-section-phantom]{Algèbre différentielle graduée}
\item \hyperref[dpa-section-phantom]{Algèbre à puissances divisées}
\item \hyperref[sdga-section-phantom]{Faisceaux différentiels gradués}
\item \hyperref[hypercovering-section-phantom]{Hyperrecouvrements}
\end{enumerate}
Schémas
\begin{enumerate}
\setcounter{enumi}{25}
\item \hyperref[schemes-section-phantom]{Schémas}
\item \hyperref[constructions-section-phantom]{Constructions de schémas}
\item \hyperref[properties-section-phantom]{Propriétés des schémas}
\item \hyperref[morphisms-section-phantom]{Morphismes de schémas}
\item \hyperref[coherent-section-phantom]{Cohomologie des schémas}
\item \hyperref[divisors-section-phantom]{Diviseurs}
\item \hyperref[limits-section-phantom]{Limites de schémas}
\item \hyperref[varieties-section-phantom]{Variétés}
\item \hyperref[topologies-section-phantom]{Topologies sur les schémas}
\item \hyperref[descent-section-phantom]{Descente}
\item \hyperref[perfect-section-phantom]{Catégories dérivées de schémas}
\item \hyperref[more-morphisms-section-phantom]{Compléments sur les morphismes}
\item \hyperref[flat-section-phantom]{Compléments sur la platitude}
\item \hyperref[groupoids-section-phantom]{Schémas en groupoïdes}
\item \hyperref[more-groupoids-section-phantom]{Compléments sur les schémas en groupoïdes}
\item \hyperref[etale-section-phantom]{Morphismes étales de schémas}
\end{enumerate}
Thèmes de théorie des schémas
\begin{enumerate}
\setcounter{enumi}{41}
\item \hyperref[chow-section-phantom]{Homologie de Chow}
\item \hyperref[intersection-section-phantom]{Théorie de l'intersection}
\item \hyperref[pic-section-phantom]{Schémas de Picard des courbes}
\item \hyperref[weil-section-phantom]{Théories cohomologiques de Weil}
\item \hyperref[adequate-section-phantom]{Modules adéquats}
\item \hyperref[dualizing-section-phantom]{Complexes dualisants}
\item \hyperref[duality-section-phantom]{Dualité pour les schémas}
\item \hyperref[discriminant-section-phantom]{Discriminants et différentes}
\item \hyperref[derham-section-phantom]{Cohomologie de de Rham}
\item \hyperref[local-cohomology-section-phantom]{Cohomologie locale}
\item \hyperref[algebraization-section-phantom]{Géométrie algébrique et formelle}
\item \hyperref[curves-section-phantom]{Courbes algébriques}
\item \hyperref[resolve-section-phantom]{Résolution des surfaces}
\item \hyperref[models-section-phantom]{Réduction semi-stable}
\item \hyperref[functors-section-phantom]{Foncteurs et morphismes}
\item \hyperref[equiv-section-phantom]{Catégories dérivées de variétés}
\item \hyperref[pione-section-phantom]{Groupes fondamentaux de schémas}
\item \hyperref[etale-cohomology-section-phantom]{Cohomologie étale}
\item \hyperref[crystalline-section-phantom]{Cohomologie cristalline}
\item \hyperref[proetale-section-phantom]{Cohomologie pro-étale}
\item \hyperref[relative-cycles-section-phantom]{Cycles relatifs}
\item \hyperref[more-etale-section-phantom]{Compléments de cohomologie étale}
\item \hyperref[trace-section-phantom]{La formule des traces}
\end{enumerate}
Espaces algébriques
\begin{enumerate}
\setcounter{enumi}{64}
\item \hyperref[spaces-section-phantom]{Espaces algébriques}
\item \hyperref[spaces-properties-section-phantom]{Propriétés des espaces algébriques}
\item \hyperref[spaces-morphisms-section-phantom]{Morphismes d'espaces algébriques}
\item \hyperref[decent-spaces-section-phantom]{Espaces algébriques décents}
\item \hyperref[spaces-cohomology-section-phantom]{Cohomologie des espaces algébriques}
\item \hyperref[spaces-limits-section-phantom]{Limites d'espaces algébriques}
\item \hyperref[spaces-divisors-section-phantom]{Diviseurs sur les espaces algébriques}
\item \hyperref[spaces-over-fields-section-phantom]{Espaces algébriques sur des corps}
\item \hyperref[spaces-topologies-section-phantom]{Topologies sur les espaces algébriques}
\item \hyperref[spaces-descent-section-phantom]{Descente et espaces algébriques}
\item \hyperref[spaces-perfect-section-phantom]{Catégories dérivées d'espaces}
\item \hyperref[spaces-more-morphisms-section-phantom]{Compléments sur les morphismes d'espaces}
\item \hyperref[spaces-flat-section-phantom]{Platitude sur les espaces algébriques}
\item \hyperref[spaces-groupoids-section-phantom]{Groupoïdes dans les espaces algébriques}
\item \hyperref[spaces-more-groupoids-section-phantom]{Compléments sur les groupoïdes dans les espaces}
\item \hyperref[bootstrap-section-phantom]{Amorçage}
\item \hyperref[spaces-pushouts-section-phantom]{Sommes amalgamées d'espaces algébriques}
\end{enumerate}
Thèmes de géométrie
\begin{enumerate}
\setcounter{enumi}{81}
\item \hyperref[spaces-chow-section-phantom]{Groupes de Chow des espaces}
\item \hyperref[groupoids-quotients-section-phantom]{Quotients de groupoïdes}
\item \hyperref[spaces-more-cohomology-section-phantom]{Compléments sur la cohomologie des espaces}
\item \hyperref[spaces-simplicial-section-phantom]{Espaces simpliciaux}
\item \hyperref[spaces-duality-section-phantom]{Dualité pour les espaces}
\item \hyperref[formal-spaces-section-phantom]{Espaces algébriques formels}
\item \hyperref[restricted-section-phantom]{Algébrisation des espaces formels}
\item \hyperref[spaces-resolve-section-phantom]{Résolution des surfaces revisitée}
\end{enumerate}
Théorie des déformations
\begin{enumerate}
\setcounter{enumi}{89}
\item \hyperref[formal-defos-section-phantom]{Théorie formelle des déformations}
\item \hyperref[defos-section-phantom]{Théorie des déformations}
\item \hyperref[cotangent-section-phantom]{Le complexe cotangent}
\item \hyperref[examples-defos-section-phantom]{Problèmes de déformation}
\end{enumerate}
Champs algébriques
\begin{enumerate}
\setcounter{enumi}{93}
\item \hyperref[algebraic-section-phantom]{Champs algébriques}
\item \hyperref[examples-stacks-section-phantom]{Exemples de champs}
\item \hyperref[stacks-sheaves-section-phantom]{Faisceaux sur les champs algébriques}
\item \hyperref[criteria-section-phantom]{Critères de représentabilité}
\item \hyperref[artin-section-phantom]{Axiomes d'Artin}
\item \hyperref[quot-section-phantom]{Espaces de Quot et de Hilbert}
\item \hyperref[stacks-properties-section-phantom]{Propriétés des champs algébriques}
\item \hyperref[stacks-morphisms-section-phantom]{Morphismes de champs algébriques}
\item \hyperref[stacks-limits-section-phantom]{Limites de champs algébriques}
\item \hyperref[stacks-cohomology-section-phantom]{Cohomologie des champs algébriques}
\item \hyperref[stacks-perfect-section-phantom]{Catégories dérivées de champs}
\item \hyperref[stacks-introduction-section-phantom]{Introduction aux champs algébriques}
\item \hyperref[stacks-more-morphisms-section-phantom]{Compléments sur les morphismes de champs}
\item \hyperref[stacks-geometry-section-phantom]{La géométrie des champs}
\end{enumerate}
Thèmes de théorie des espaces de modules
\begin{enumerate}
\setcounter{enumi}{107}
\item \hyperref[moduli-section-phantom]{Champs de modules}
\item \hyperref[moduli-curves-section-phantom]{Espaces de modules de courbes}
\end{enumerate}
Divers
\begin{enumerate}
\setcounter{enumi}{109}
\item \hyperref[examples-section-phantom]{Exemples}
\item \hyperref[exercises-section-phantom]{Exercices}
\item \hyperref[guide-section-phantom]{Guide bibliographique}
\item \hyperref[desirables-section-phantom]{Desiderata}
\item \hyperref[coding-section-phantom]{Style de programmation}
\item \hyperref[obsolete-section-phantom]{Obsolète}
\item \hyperref[fdl-section-phantom]{Licence de documentation libre GNU}
\item \hyperref[index-section-phantom]{Index généré automatiquement}
\end{enumerate}
\end{multicols}


\bibliography{my}
\bibliographystyle{amsalpha}

\end{document}
